% Lesson 47 — Vocabulary: Words and expresssions related to campaigning and voting for leaders — Anglais Tle SH — Cours signé OnBuch+
% Programme officiel MINESEC. Compiler avec : tectonic EN47-vocabulary-words-and-expresssions-related-to-campaignin.tex
\newcommand{\DOCMATIERE}{Anglais}
\newcommand{\DOCNIVEAU}{Tle SH}
\newcommand{\DOCMODULE}{Chapter 4 : Citizenship and Human Rights}
\newcommand{\DOCLECON}{EN47}
\newcommand{\DOCTITRE}{Lesson 47 — Vocabulary: Words and expresssions related to campaigning and voting for leaders}
% OnBuch+ — préambule « cours signés » ENRICHI (compilé avec tectonic / XeLaTeX)
% Système visuel « Pop » d'OnBuch+ : fond crème (jamais blanc), bordures encre
% épaisses, ombres « dures » décalées, titres Archivo Black en capitales.
% Version MATHÉMATIQUES : graphes (pgfplots/tikz), boîtes pédagogiques
% (définition, propriété, méthode, attention, le savais-tu, exercice résolu,
% exercice, TP, bilan), page de garde et signature OnBuch+ ancrée en bas.
%
% Macros attendues dans main.tex AVANT \input{preamble} :
%   \DOCMATIERE  \DOCNIVEAU  \DOCMODULE  \DOCLECON (numéro)  \DOCTITRE
\documentclass[11pt]{article}

\usepackage{fontspec}
\usepackage[english,french]{babel} % français = langue principale ; l'anglais via \eng{..} / textebox
\usepackage[a4paper,margin=2cm,top=3.4cm,bottom=2.8cm,headheight=1.4cm,footskip=1.3cm]{geometry}
\usepackage{amsmath,amssymb,amsfonts}
\usepackage{mathtools} % \xmapsto, \coloneqq... souvent utilisés par les modèles
\usepackage{cancel} % simplifications barrées (bilans de forces)
% \abs{z} (module, valeur absolue) et \norm{u} (norme), utilisés par les modèles
\providecommand{\abs}{}\let\abs\relax\DeclarePairedDelimiter{\abs}{\lvert}{\rvert}
\providecommand{\norm}{}\let\norm\relax\DeclarePairedDelimiter{\norm}{\lVert}{\rVert}
\usepackage[table]{xcolor} % [table] : fournit aussi \rowcolors (alternance de lignes)
\usepackage{pagecolor}
\usepackage{tikz}
\usetikzlibrary{arrows.meta,positioning,calc,shapes.geometric,shapes.misc,shapes.arrows,decorations.pathmorphing,decorations.pathreplacing,decorations.markings,patterns,shadows,mindmap,trees,fit,backgrounds,matrix,intersections,angles,quotes}
\usepackage{pgfplots}
\usepackage[version=4]{mhchem} % équations nucléaires \ce{^{235}_{92}U}
\usepackage[european,siunitx]{circuitikz} % schémas électriques
\pgfplotsset{compat=1.18}
\usepgfplotslibrary{fillbetween,groupplots} % aires entre courbes (calcul intégral)
\usepackage{enumitem}
\usepackage{fancyhdr}
% [most] charge toutes les bibliothèques tcolorbox (listings, minted, raster,
% documentation...) et gonfle inutilement la mémoire TeX sur un document long
% (~300 boîtes) : on ne charge que ce qui est réellement utilisé.
\usepackage[skins,breakable]{tcolorbox}
\usepackage{graphicx}
\usepackage{colortbl}
\usepackage{booktabs}
\usepackage{tabularx}
\usepackage{diagbox} % case d'en-tête diagonale des tableaux à double entrée
% Colonnes extensibles centrée / gauche / droite (C, L, R), souvent utilisées par les modèles
\newcolumntype{C}{>{\centering\arraybackslash}X}
\newcolumntype{L}{>{\raggedright\arraybackslash}X}
\newcolumntype{R}{>{\raggedleft\arraybackslash}X}
\usepackage{array}
\usepackage{multicol}
\usepackage{multirow}
\usepackage{siunitx}
% parse-numbers=false : accepte aussi \SI{2,5 \times 10^{-3}}{...}
\sisetup{locale=FR,output-decimal-marker={,},group-minimum-digits=4,parse-numbers=false}
\DeclareSIUnit\ohm{\ensuremath{\symup{\Omega}}} % U+2126 absent de la police
\DeclareSIUnit\division{div} % réglages d'oscilloscope (ms/div, V/div)
\DeclareSIUnit\tr{tr} % tours (vitesse de rotation, ex. \SI{1500}{\tr\per\minute})
\DeclareSIUnit\molar{M} % concentration molaire (ex. \SI{3}{\molar}, \SI{1}{\micro\molar})
\DeclareSIUnit\an{an} % année (le modèle confond parfois avec \year, un compteur TeX) — ex. \SI{5}{\kilo\meter\per\an}
\DeclareSIUnit\FCFA{FCFA} % franc CFA (ex. \SI{500}{\FCFA\per\kilo\gram})
\DeclareSIUnit\degreeBrix{\textdegree Brix} % teneur en sucre d'un sirop/jus (réfractométrie)
\DeclareSIUnit\month{mois} % le modèle utilise parfois \month, un compteur TeX, comme unité de durée
\usepackage{qrcode}
% \degree : commande fréquemment utilisée par les modèles pour un angle ou une
% température en dehors de \SI{}{} (non fournie par les paquets chargés).
\providecommand{\degree}{^{\circ}}
% \qed (fin de démonstration), utilisé par les modèles sans amsthm
\providecommand{\qed}{\hfill$\square$}
% \barre : double barre des valeurs interdites dans un tableau de variations (tkz-tab)
\providecommand{\barre}{\ensuremath{\|}}
% environnement proof (amsthm non chargé) utilisé par les modèles
\ifdefined\proof\else\newenvironment{proof}[1][Démonstration]{\par\noindent\textit{#1.}\ }{\qed\par}\fi
\usepackage{lastpage}
\usepackage{hyperref}
\hypersetup{hidelinks}

% ============================================================
% Palette « Pop » OnBuch+ — mêmes hex que la classe P (plus_ui.dart)
% ============================================================
\definecolor{poporange}{HTML}{F59321}
\definecolor{popdark}{HTML}{E07A0C}
\definecolor{popcream}{HTML}{FFF6E8}
\definecolor{popink}{HTML}{1C1714}
\definecolor{poppurple}{HTML}{7A5AE0}
\definecolor{popgreen}{HTML}{2E9E6A}
\definecolor{poppink}{HTML}{E0457B}
\definecolor{popblue}{HTML}{2F7DE1}
\definecolor{popgold}{HTML}{C98A12}
\definecolor{popmuted}{HTML}{8A7F76}
\definecolor{popline}{HTML}{EADFD2}
\definecolor{popgray}{HTML}{8A7F76}
\colorlet{popgrey}{popgray}
% Alias tolérants (le modèle nomme parfois les couleurs différemment)
\colorlet{popwhite}{white}
\colorlet{popblack}{popink}
\colorlet{popred}{poppink}
\colorlet{popyellow}{popgold}
% Teintes claires (fonds de tableaux/encadrés) — le modèle nomme parfois
% les couleurs avec un suffixe L (ex. popblueL) sans les avoir définies.
\colorlet{poporangeL}{poporange!20}
\colorlet{popdarkL}{popdark!20}
\colorlet{poppurpleL}{poppurple!20}
\colorlet{popgreenL}{popgreen!20}
\colorlet{poppinkL}{poppink!20}
\colorlet{popblueL}{popblue!20}
\colorlet{popgoldL}{popgold!20}
\colorlet{popredL}{popred!20}
\colorlet{popgrayL}{popgray!20}
% Teintes claires pour fonds de tableaux / zones
\colorlet{popblueL}{popblue!10!white}
\colorlet{poporangeL}{poporange!14!white}
\colorlet{popgreenL}{popgreen!12!white}
\colorlet{poppurpleL}{poppurple!12!white}
\colorlet{poppinkL}{poppink!12!white}
\colorlet{popgoldL}{popgold!14!white}

\pagecolor{popcream}

% ============================================================
% Typographie
% ============================================================
\defaultfontfeatures{Ligatures=TeX}
\setmainfont{PlusJakartaSans-Medium.ttf}[Path=../../../fonts/,BoldFont=PlusJakartaSans-Bold.ttf]
% Symboles, API et emojis absents de la police principale : repli sur DejaVu Sans (caractères actifs)
\newfontface\symfont{DejaVuSans.ttf}[Path=../../../fonts/]
\makeatletter
\newcommand\symactive[1]{\catcode#1=13 \begingroup\lccode`\~=#1 \lowercase{\endgroup\def~}{{\symfont\char#1}}}
\newcommand\symblank[1]{\catcode#1=13 \begingroup\lccode`\~=#1 \lowercase{\endgroup\def~}{}}
\newcommand\symhyph[1]{\catcode#1=13 \begingroup\lccode`\~=#1 \lowercase{\endgroup\def~}{\mbox{-}}}
\newcount\symc
\newcommand\symrange[2]{\symc=#1 \@whilenum\symc<#2 \do{\expandafter\symactive\expandafter{\the\symc}\advance\symc by 1 }}
\newcommand\symblankrange[2]{\symc=#1 \@whilenum\symc<#2 \do{\expandafter\symblank\expandafter{\the\symc}\advance\symc by 1 }}
\makeatother
\symrange{"0250}{"02FF}\symactive{"2713}\symactive{"2714}\symactive{"2717}\symactive{"2718}\symactive{"2610}\symactive{"2611}\symactive{"2612}
\symhyph{"2011}\symhyph{"2E1A}\symblankrange{"1F300}{"1FAFF}\symblank{"FE0F}
% Maths en sans-serif (Fira Math) avec chiffres et lettres droites de Plus
% Jakarta, pour que les formules se fondent dans le texte.
\usepackage[math-style=ISO,bold-style=ISO]{unicode-math}
\setmathfont{FiraMath-Regular.otf}[Path=../../../fonts/]
\setmathfont{PlusJakartaSans-Medium.ttf}[Path=../../../fonts/,range={up/num,up/{Latin,latin},\mathup}]
\usepackage{icomma} % virgule décimale sans espace en mode math
% Caractères mathématiques Unicode absents des polices de texte (Plus Jakarta,
% Archivo Black) : rendus par leur équivalent mathématique, en texte comme en
% formule — sinon ils s'affichent en carré vide (ex. « Arithmétique dans ℤ »).
\usepackage{newunicodechar}
\newunicodechar{ℝ}{\ensuremath{\mathbb{R}}}
\newunicodechar{ℤ}{\ensuremath{\mathbb{Z}}}
\newunicodechar{ℂ}{\ensuremath{\mathbb{C}}}
\newunicodechar{ℕ}{\ensuremath{\mathbb{N}}}
\newunicodechar{ℚ}{\ensuremath{\mathbb{Q}}}
\newunicodechar{𝔻}{\ensuremath{\mathbb{D}}}
\newunicodechar{α}{\ensuremath{\alpha}}
\newunicodechar{β}{\ensuremath{\beta}}
\newunicodechar{γ}{\ensuremath{\gamma}}
\newunicodechar{δ}{\ensuremath{\delta}}
\newunicodechar{ε}{\ensuremath{\varepsilon}}
\newunicodechar{θ}{\ensuremath{\theta}}
\newunicodechar{λ}{\ensuremath{\lambda}}
\newunicodechar{μ}{\ensuremath{\mu}}
\newunicodechar{σ}{\ensuremath{\sigma}}
\newunicodechar{τ}{\ensuremath{\tau}}
\newunicodechar{φ}{\ensuremath{\varphi}}
\newunicodechar{ψ}{\ensuremath{\psi}}
\newunicodechar{ω}{\ensuremath{\omega}}
\newunicodechar{Σ}{\ensuremath{\Sigma}}
% majuscules produites par \MakeUppercase dans les titres (α-aminés -> Α)
\newunicodechar{Α}{\ensuremath{\alpha}}
\newunicodechar{Β}{\ensuremath{\beta}}
\newunicodechar{Γ}{\ensuremath{\Gamma}}
\newunicodechar{Δ}{\ensuremath{\Delta}}
\newunicodechar{Ε}{\ensuremath{\varepsilon}}
\newunicodechar{Θ}{\ensuremath{\Theta}}
\newunicodechar{Λ}{\ensuremath{\Lambda}}
\newunicodechar{Μ}{\ensuremath{\mu}}
\newunicodechar{Τ}{\ensuremath{\tau}}
\newunicodechar{Φ}{\ensuremath{\Phi}}
\newunicodechar{Ψ}{\ensuremath{\Psi}}
\newunicodechar{Ω}{\ensuremath{\Omega}}
\newunicodechar{↦}{\ensuremath{\mapsto}}
\newunicodechar{∈}{\ensuremath{\in}}
\newunicodechar{∉}{\ensuremath{\notin}}
\newunicodechar{∧}{\ensuremath{\wedge}}
\newunicodechar{⇔}{\ensuremath{\Leftrightarrow}}
\newunicodechar{⟺}{\ensuremath{\Longleftrightarrow}}
\newunicodechar{⇒}{\ensuremath{\Rightarrow}}
\newunicodechar{⟹}{\ensuremath{\Longrightarrow}}
\newunicodechar{∘}{\ensuremath{\circ}}
\newunicodechar{∪}{\ensuremath{\cup}}
\newunicodechar{∩}{\ensuremath{\cap}}
\newunicodechar{≡}{\ensuremath{\equiv}}
\newunicodechar{⊂}{\ensuremath{\subset}}
\newunicodechar{⊥}{\ensuremath{\perp}}
\newunicodechar{∃}{\ensuremath{\exists}}
\newunicodechar{∀}{\ensuremath{\forall}}
\newunicodechar{∛}{\ensuremath{\sqrt[3]{\,}}}
\newunicodechar{∜}{\ensuremath{\sqrt[4]{\,}}}
\newunicodechar{⃗}{\ensuremath{{}^{\to}}}
\newunicodechar{ⁿ}{\ifmmode^{n}\else\textsuperscript{\ensuremath{n}}\fi}
\newunicodechar{ˣ}{\ifmmode^{x}\else\textsuperscript{\ensuremath{x}}\fi}
\newunicodechar{ᵇ}{\ifmmode^{b}\else\textsuperscript{\ensuremath{b}}\fi}
\newunicodechar{ʳ}{\ifmmode^{r}\else\textsuperscript{\ensuremath{r}}\fi}
\newunicodechar{ᵘ}{\ifmmode^{u}\else\textsuperscript{\ensuremath{u}}\fi}
\newunicodechar{ʸ}{\ifmmode^{y}\else\textsuperscript{\ensuremath{y}}\fi}
\newunicodechar{ᵃ}{\ifmmode^{a}\else\textsuperscript{\ensuremath{a}}\fi}
\newunicodechar{ᵉ}{\ifmmode^{e}\else\textsuperscript{\ensuremath{e}}\fi}
\newunicodechar{ᵏ}{\ifmmode^{k}\else\textsuperscript{\ensuremath{k}}\fi}
\newunicodechar{ᵖ}{\ifmmode^{p}\else\textsuperscript{\ensuremath{p}}\fi}
\newunicodechar{ᴺ}{\ifmmode^{N}\else\textsuperscript{\ensuremath{N}}\fi}
\newunicodechar{ᶜ}{\ifmmode^{c}\else\textsuperscript{\ensuremath{c}}\fi}
\newunicodechar{ʰ}{\ifmmode^{h}\else\textsuperscript{\ensuremath{h}}\fi}
\newunicodechar{ᵠ}{\ifmmode^{q}\else\textsuperscript{\ensuremath{q}}\fi}
\newunicodechar{ₙ}{\ifmmode_{n}\else\textsubscript{\ensuremath{n}}\fi}
\newunicodechar{ₐ}{\ifmmode_{a}\else\textsubscript{\ensuremath{a}}\fi}
\newunicodechar{ᵢ}{\ifmmode_{i}\else\textsubscript{\ensuremath{i}}\fi}
\newunicodechar{ᵣ}{\ifmmode_{r}\else\textsubscript{\ensuremath{r}}\fi}
\newunicodechar{ₖ}{\ifmmode_{k}\else\textsubscript{\ensuremath{k}}\fi}
\newunicodechar{ᵦ}{\ifmmode_{\beta}\else\textsubscript{\ensuremath{\beta}}\fi}
\newunicodechar{ₓ}{\ifmmode_{x}\else\textsubscript{\ensuremath{x}}\fi}
\newfontfamily\popdisplay{ArchivoBlack-Regular.ttf}[Path=../../../fonts/]
\newfontfamily\poplogo{SpaceGrotesk-Bold.ttf}[Path=../../../fonts/]
\newfontfamily\popsemi{PlusJakartaSans-SemiBold.ttf}[Path=../../../fonts/]
\newfontfamily\popxbold{PlusJakartaSans-ExtraBold.ttf}[Path=../../../fonts/]
\newfontfamily\popxboldls{PlusJakartaSans-ExtraBold.ttf}[Path=../../../fonts/,LetterSpace=3.0]

\setlist[enumerate]{leftmargin=1.7em,itemsep=5pt,topsep=4pt}
\setlist[itemize]{leftmargin=1.7em,itemsep=4pt,topsep=4pt,label={\color{poporange}\textbullet}}
\setlength{\parindent}{0pt}
\setlength{\parskip}{5pt}
\renewcommand{\arraystretch}{1.25}

% pgfplots : style Pop par défaut
\pgfplotsset{
  every axis/.append style={axis line style={popink,line width=1pt},tick style={popink},
    grid style={popline,line width=0.5pt},label style={font=\small},tick label style={font=\footnotesize}},
  popcurve/.style={poporange,line width=1.8pt,no markers,smooth},
}
% Même style disponible dans un \draw TikZ simple (hors pgfplots)
\tikzset{popcurve/.style={poporange,line width=1.8pt}}
\tikzset{popgrid/.style={popline,very thin}}
\colorlet{popcurve}{poporange} % le nom du style est parfois utilisé comme couleur

% ============================================================
% Pill / squiggle
% ============================================================
\newcommand{\poppill}[3]{%
  \tikz[baseline=-0.55ex]\node[draw=popink,line width=1.2pt,rounded corners=2.6pt,%
    fill=#2,inner xsep=6.5pt,inner ysep=3pt]{%
    \popxboldls\fontsize{7}{7}\selectfont\color{#3}\MakeUppercase{#1}};%
}
\newcommand{\popsquiggle}[1]{%
  \tikz[baseline=-0.4ex]{
    \draw[#1,line width=2.6pt,line cap=round]
      plot [smooth] coordinates {(0,0.09) (0.34,0.01) (0.68,0.21) (1.02,0.01) (1.36,0.10)};
  }%
}
% Mot-clé surligné (marqueur orange sous le texte)
\newcommand{\cle}[1]{{\popxbold\color{popdark}#1}}

% ============================================================
% En-tête / pied
% ============================================================
\pagestyle{fancy}
\fancyhf{}
\renewcommand{\headrulewidth}{0pt}
\renewcommand{\footrulewidth}{0pt}
\lhead{\raisebox{-0.15ex}{\poplogo\fontsize{13}{13}\selectfont\color{popink}OnBuch\color{poporange}+}}
\rhead{\poppill{\DOCMATIERE\ \textbullet\ \DOCNIVEAU\ \textbullet\ Leçon \DOCLECON}{white}{popink}}
\lfoot{\popxbold\fontsize{7.6}{7.6}\selectfont\color{popmuted}\MakeUppercase{Cours signé OnBuch+}\ \textbullet\ {\popsemi\DOCTITRE}}
\rfoot{\tikz[baseline=-0.6ex]\node[rounded corners=8.5pt,fill=popink,minimum height=17pt,minimum width=17pt,inner xsep=5pt,inner sep=0]{\popxbold\fontsize{8}{8}\selectfont\color{popcream}\thepage\,/\,\pageref*{LastPage}};}

% ============================================================
% Titres
% ============================================================
\newcommand{\chaptitre}[2]{%
  \begin{tcolorbox}[enhanced,colback=poporange,colframe=popink,boxrule=2.6pt,arc=11pt,
      drop shadow=popink,left=15pt,right=15pt,top=12pt,bottom=11pt,before skip=3pt,after skip=18pt]
    \poppill{#2}{popink}{popcream}\par\vspace{9pt}
    {\raggedright\hyphenpenalty=10000\exhyphenpenalty=10000\popdisplay\fontsize{22}{24}\selectfont\color{popcream}\MakeUppercase{#1}\par}
  \end{tcolorbox}
}
% Section numérotée : pastille orange avec le numéro + titre Archivo
\newcounter{popsec}
\newcommand{\coursec}[1]{%
  \stepcounter{popsec}\setcounter{popsub}{0}%
  \par\vspace{16pt}\noindent
  \begin{minipage}{\linewidth}
  \tikz[baseline=-0.7ex]\node[draw=popink,line width=1.6pt,fill=poporange,rounded corners=4pt,minimum size=22pt,inner sep=0]{\popdisplay\fontsize{12}{12}\selectfont\color{popcream}\thepopsec};%
  \hspace{8pt}{\popdisplay\fontsize{15}{17}\selectfont\color{popink}\MakeUppercase{#1}}\par
  \vspace{4pt}\noindent{\color{poporange}\rule{36pt}{3.6pt}}
  \end{minipage}\par\nopagebreak\vspace{8pt}
}
\newcounter{popsub}
\newcommand{\courssub}[1]{%
  \stepcounter{popsub}%
  \par\vspace{9pt}\noindent{\popxbold\fontsize{11.5}{13}\selectfont\color{popink}{\color{poporange}\thepopsec.\thepopsub}\hspace{6pt}#1}\par\nopagebreak\vspace{4pt}
}

% ============================================================
% Boîtes pédagogiques — même recette : fond blanc, bordure épaisse colorée,
% ombre dure encre, badge plein. Titre optionnel [#1].
% ============================================================
\tcbset{popbase/.style={breakable,enhanced,colback=white,boxrule=2.3pt,arc=9pt,
  shadow={3pt}{-3pt}{0pt}{popink},left=14pt,right=14pt,top=11pt,bottom=11pt,before skip=11pt,after skip=15pt,
  fontupper=\popsemi\fontsize{10.6}{14.5}\selectfont\color{popink},
  fontlower=\popsemi\fontsize{10.6}{14.5}\selectfont\color{popink}}}
% Coche dessinée (le glyphe ✓ n'existe pas dans les polices du document)
\newcommand{\popcheck}[1]{\tikz[baseline=-0.5ex]\draw[#1,line width=1.6pt,line cap=round,line join=round](-0.13,0.0)--(-0.04,-0.09)--(0.13,0.1);}
\providecommand{\coche}{\popcheck{popgreen}}
\providecommand{\resultat}[1]{\par\smallskip\noindent\fcolorbox{popgreen}{popgreenL}{\parbox{\dimexpr\linewidth-2\fboxsep-2\fboxrule\relax}{\textbf{Résultat.} #1}}\par\smallskip}
\newcommand{\popboxhead}[3]{\poppill{#1}{#2}{white}\ifx\relax#3\relax\else\hspace{6pt}{\popxbold\fontsize{10.6}{12}\selectfont\color{popink}#3}\fi\par\vspace{7pt}}

\newtcolorbox{aretenir}[1][]{popbase,colframe=popgreen,before upper={\popboxhead{À retenir}{popgreen}{#1}}}
% Texte en anglais (document de lecture, dialogue, extrait) : cadre
% d'étude. Un titre optionnel (source, auteur) s'affiche dans l'en-tête.
\newtcolorbox{textebox}[1][]{popbase,colframe=popblue,before upper={\popboxhead{Text}{popblue}{#1}\begin{otherlanguage}{english}},after upper={\end{otherlanguage}}}
% Marques ✓ / ✗ (forme correcte / fautive) souvent utilisées par les modèles
\providecommand{\cross}{\textcolor{poppink}{\ensuremath{\times}}}
\providecommand{\xmark}{\cross}\providecommand{\crossmark}{\cross}\providecommand{\cmark}{\ensuremath{\checkmark}}
% Ligne de réponse / justification à compléter (inventée par les modèles)
\providecommand{\justification}[1]{\par\noindent\textit{Justification~:} \dotfill\par}
% \textipa (package tipa non chargé) : la police ne couvre pas l'API -> simple passage
\providecommand{\textipa}[1]{#1}
% Soulignement (ulem non chargé) : \uline, \ul
\providecommand{\uline}[1]{\underline{#1}}\providecommand{\ul}[1]{\underline{#1}}
% \glossaire{\item ...} inventée par les modèles : liste de glossaire
\providecommand{\glossaire}[1]{\begin{itemize}#1\end{itemize}}
% \alert (beamer) inventée par les modèles : texte mis en évidence
\providecommand{\alert}[1]{\textbf{\textcolor{poppink}{#1}}}
% variantes ASCII d'ordinaux écrites en macro
\providecommand{\iere}{\textsuperscript{re}}\providecommand{\ieme}{\textsuperscript{e}}\providecommand{\ere}{\textsuperscript{re}}\providecommand{\eme}{\textsuperscript{e}}\providecommand{\er}{\textsuperscript{er}}\providecommand{\ie}{\textsuperscript{e}}
% Ordinaux français écrits en macro par les modèles : 1\ière, 3\ième
\newcommand{\ière}{\textsuperscript{re}}\newcommand{\ième}{\textsuperscript{e}}
% \exemple (inventée par les modèles) : étiquette « Exemple : »
\providecommand{\exemple}{\textbf{Exemple~:}\ }
% \square utilisable aussi hors mode math (cases à cocher)
\AtBeginDocument{\let\squareMath\square\renewcommand{\square}{\ensuremath{\squareMath}}}
% Mot ou phrase en anglais à l'intérieur d'un paragraphe français
\newcommand{\eng}[1]{\ifmmode\text{\textit{#1}}\else\begingroup\selectlanguage{english}\itshape #1\endgroup\fi}
\let\esp\eng % alias toléré
% flèches d'intonation inventées par les modèles
\providecommand{\textsearrow}{}\renewcommand{\textsearrow}{\ensuremath{\searrow}}\providecommand{\textnearrow}{}\renewcommand{\textnearrow}{\ensuremath{\nearrow}}
% \fr{..} : mot français (inventé par les modèles)
\providecommand{\fr}[1]{\textit{#1}}
\newtcolorbox{exemplebox}[1][]{popbase,colframe=poporange,before upper={\popboxhead{Exemple}{poporange}{#1}}}
\newtcolorbox{definition}[1][]{popbase,colframe=popblue,before upper={\popboxhead{Définition}{popblue}{#1}}}
\newtcolorbox{propriete}[1][]{popbase,colframe=poppurple,before upper={\popboxhead{Propriété}{poppurple}{#1}}}
\newtcolorbox{methode}[1][]{popbase,colframe=popgold,before upper={\popboxhead{Méthode}{popgold}{#1}}}
\newtcolorbox{attention}[1][]{popbase,colframe=poppink,before upper={\popboxhead{Attention}{poppink}{#1}}}
\newtcolorbox{experience}[1][]{popbase,colframe=popblue,colback=popblueL,before upper={\popboxhead{Expérience}{popblue}{#1}}}
% « Le savais-tu ? » : carte violette pleine (contexte, culture, Cameroun)
\newtcolorbox{savaistu}[1][]{popbase,colback=poppurple,colframe=popink,
  fontupper=\popsemi\fontsize{10.4}{14.2}\selectfont\color{white},
  before upper={\poppill{Le savais-tu\,?}{white}{poppurple}\ifx\relax#1\relax\else\hspace{6pt}{\popxbold\color{white}#1}\fi\par\vspace{7pt}}}
% Exercice résolu : énoncé en haut, \tcblower puis la solution sur fond crème
\newcounter{popexo}\newcounter{popexr}
\newtcolorbox{exoresolu}[1][]{popbase,colframe=poporange,colbacklower=popcream,
  segmentation style={popink,line width=1.2pt,dashed},
  before upper={\stepcounter{popexr}\popboxhead{Exercice résolu \thepopexr}{poporange}{#1}},
  before lower={\poppill{Solution}{popink}{popcream}\par\vspace{6pt}}}
% Exercice (non résolu en place) — la correction est dans la partie Corrigés
\newtcolorbox{exercice}[1][]{popbase,colframe=popink,
  before upper={\stepcounter{popexo}\popboxhead{Exercice \thepopexo}{popink}{#1}}}
% Corrigé
\newtcolorbox{corrige}[1][]{popbase,colframe=popgreen,colback=popgreenL,
  before upper={\popboxhead{Corrigé}{popgreen}{#1}}}
% Objectifs / prérequis / bilan
\newtcolorbox{objectifs}{popbase,colframe=popink,colback=poporangeL,
  before upper={\popboxhead{Objectifs de la leçon}{popink}{}}}
\newtcolorbox{prerequis}{popbase,colframe=popink,colback=popblueL,
  before upper={\popboxhead{Prérequis}{popblue}{}}}
\newtcolorbox{bilan}{popbase,colframe=popink,colback=white,boxrule=2.8pt,
  before upper={\popboxhead{Fiche bilan}{poporange}{L'essentiel à connaître pour l'examen}}}
% Figure encadrée avec légende
\newtcolorbox{popfigure}[1][]{enhanced,breakable=false,colback=white,colframe=popline,boxrule=1.4pt,arc=8pt,
  left=10pt,right=10pt,top=10pt,bottom=8pt,before skip=10pt,after skip=12pt,halign=center,
  lower separated=false,halign lower=center,
  fontlower=\popsemi\fontsize{9}{11.5}\selectfont\color{popmuted},
  #1}
\newcommand{\legende}[1]{\par\vspace{6pt}{\popsemi\fontsize{9}{11.5}\selectfont\color{popmuted}#1}}

% ============================================================
% Signature OnBuch+
% ============================================================
\newcommand{\poplogoinline}[1]{{\poplogo\fontsize{#1}{#1}\selectfont\color{popink}OnBuch\color{poporange}+}}
\newcommand{\signatureonbuch}{%
  \begin{tcolorbox}[enhanced,colback=popink,colframe=popink,boxrule=2.2pt,arc=10pt,
      drop shadow=poporange,left=14pt,right=14pt,top=10pt,bottom=10pt,before skip=0pt,after skip=0pt]
    \tikz[baseline=-0.7ex]\node[circle,fill=popgreen,draw=popcream,line width=1.3pt,minimum size=22pt,inner sep=0]{\popcheck{white}};%
    \hspace{9pt}\begin{minipage}[c]{0.62\linewidth}
      {\popxbold\fontsize{12}{13}\selectfont\color{popcream}Signé\ \textbullet\ {\poplogo OnBuch\color{poporange}+}}\par\vspace{2pt}
      {\raggedright\popsemi\fontsize{8}{10.5}\selectfont\color{popline}\MakeUppercase{Équipe pédagogique OnBuch+}\\\MakeUppercase{Conforme au programme officiel MINESEC}\par}
    \end{minipage}\hfill
    {\poplogo\fontsize{22}{22}\selectfont\color{popcream}OnBuch\color{poporange}+}
  \end{tcolorbox}%
}

% ============================================================
% Page de garde
% #1 titre · #2 matière · #3 niveau · #4 module · #5 n° leçon · #6 durée ·
% #7 description / objectifs (texte ou itemize)
% ============================================================
\newcommand{\pagegarde}[7]{%
  \thispagestyle{empty}
  \newgeometry{a4paper,margin=1.7cm,top=1.6cm,bottom=1.4cm}%
  \begin{tikzpicture}[remember picture,overlay]
    \fill[poporange!18] ([xshift=-3.2cm,yshift=-3.6cm]current page.north east) circle (4.6cm);
    \fill[poppurple!14] ([xshift=2.4cm,yshift=7.6cm]current page.south west) circle (3.8cm);
  \end{tikzpicture}%
  {\poplogo\fontsize{30}{30}\selectfont\color{popink}OnBuch\color{poporange}+}\hfill
  \raisebox{0.6ex}{\poppill{Cours signé \textbullet\ Premium}{popink}{popcream}}\par
  \vspace{1pt}\popsquiggle{poppurple}\par
  \vspace{8pt}
  \begin{tcolorbox}[enhanced,colback=poporange,colframe=popink,boxrule=2.8pt,arc=13pt,
      drop shadow=popink,left=18pt,right=18pt,top=12pt,bottom=13pt,after skip=10pt]
    \poppill{#2 \textbullet\ #3}{popink}{popcream}\hspace{5pt}\poppill{#4}{white}{popink}\par\vspace{8pt}
    {\popdisplay\fontsize{14}{14}\selectfont\color{popink}LEÇON #5}\par\vspace{4pt}
    {\raggedright\hyphenpenalty=10000\exhyphenpenalty=10000\popdisplay\fontsize{25}{27}\selectfont\color{popcream}\MakeUppercase{#1}\par}
    \vspace{8pt}
    {\popxbold\fontsize{9.5}{9.5}\selectfont\color{popink}\MakeUppercase{Durée conseillée : #6}}
  \end{tcolorbox}
  \begin{tcolorbox}[enhanced,colback=white,colframe=popink,boxrule=2.2pt,arc=10pt,
      drop shadow=popink,left=16pt,right=16pt,top=10pt,bottom=10pt,after skip=8pt]
    \poppill{Dans cette leçon}{poppurple}{white}\par\vspace{5pt}
    \popsemi\fontsize{9.3}{12.6}\selectfont\color{popink}#7
  \end{tcolorbox}
  \noindent\begin{minipage}[t]{0.487\linewidth}
    \begin{tcolorbox}[enhanced,colback=popgreen,colframe=popink,boxrule=2.2pt,arc=10pt,
        drop shadow=popink,left=11pt,right=11pt,top=10pt,bottom=10pt,height=3.1cm,valign=center]
      \tcbox[colback=white,colframe=popink,boxrule=1.3pt,arc=5pt,boxsep=0pt,nobeforeafter,
        left=3pt,right=3pt,top=3pt,bottom=3pt,valign=center]{\qrcode[height=1.9cm]{https://play.google.com/store/apps/details?id=com.luvvix.onbuch&hl=fr}}%
      \hspace{8pt}\begin{minipage}[c]{0.5\linewidth}
        {\popdisplay\fontsize{11}{12.5}\selectfont\color{white}\MakeUppercase{Télécharge OnBuch}}\par\vspace{4pt}
        {\raggedright\popsemi\fontsize{8.4}{11}\selectfont\color{white}Gratuit \textbullet\ Google Play\\Tuteur IA, annales, concours, résultats\par}
      \end{minipage}
    \end{tcolorbox}
  \end{minipage}\hfill
  \begin{minipage}[t]{0.487\linewidth}
    \begin{tcolorbox}[enhanced,colback=poppurple,colframe=popink,boxrule=2.2pt,arc=10pt,
        drop shadow=popink,left=11pt,right=11pt,top=10pt,bottom=10pt,height=3.1cm,valign=center]
      \tikz[baseline=-0.7ex]\node[circle,fill=white,draw=popink,line width=1.3pt,minimum size=1.6cm,inner sep=0]{\popdisplay\fontsize{24}{24}\selectfont\color{poppurple}+};%
      \hspace{8pt}\begin{minipage}[c]{0.6\linewidth}
        {\popdisplay\fontsize{11}{12.5}\selectfont\color{white}\MakeUppercase{Passe à OnBuch+}}\par\vspace{4pt}
        {\raggedright\popsemi\fontsize{8.4}{11}\selectfont\color{white}Tous les cours signés, Tuteur IA illimité, fiches PDF — onglet OnBuch+ de l'app\par}
      \end{minipage}
    \end{tcolorbox}
  \end{minipage}
  \vspace{8pt}
  \popcontenu
  \vfill
  \signatureonbuch
  \restoregeometry
  \newpage
}

% Rangée « Ce cours contient » (page de garde)
\newcommand{\poptile}[3]{% #1 couleur #2 gros texte #3 légende
  \begin{tcolorbox}[enhanced,colback=#1,colframe=popink,boxrule=2pt,arc=9pt,shadow={2.5pt}{-2.5pt}{0pt}{popink},
      left=6pt,right=6pt,top=9pt,bottom=9pt,halign=center,valign=center,height=1.75cm,width=\linewidth,nobeforeafter]
    {\popdisplay\fontsize{12.5}{13}\selectfont\color{white}\MakeUppercase{#2}}\par\vspace{3pt}
    {\centering\popsemi\fontsize{7.8}{9.5}\selectfont\color{white}#3\par}
  \end{tcolorbox}}
\newcommand{\popcontenu}{%
  \par\noindent{\popxboldls\fontsize{7.5}{7.5}\selectfont\color{popmuted}CE COURS SIGNÉ CONTIENT}\par\vspace{7pt}
  \noindent\begin{minipage}[t]{0.235\linewidth}\poptile{popblue}{Cours}{complet, illustré, pas à pas}\end{minipage}\hfill
  \begin{minipage}[t]{0.235\linewidth}\poptile{poporange}{Méthodes}{et exercices résolus}\end{minipage}\hfill
  \begin{minipage}[t]{0.235\linewidth}\poptile{poppink}{Exercices}{type examen + corrigés}\end{minipage}\hfill
  \begin{minipage}[t]{0.235\linewidth}\poptile{popgreen}{Fiche bilan}{l'essentiel à retenir}\end{minipage}\par}

% Sommaire
\newtcolorbox{sommaire}{enhanced,colback=white,colframe=popink,boxrule=2.2pt,arc=10pt,
  drop shadow=popink,left=18pt,right=18pt,top=13pt,bottom=13pt,before skip=6pt,after skip=20pt,
  before upper={\poppill{Sommaire}{poppurple}{white}\par\vspace{9pt}\popsemi\fontsize{10.6}{14.5}\selectfont\color{popink}}}

% Signature de fin de cours — poussée tout en bas de la dernière page
\newcommand{\findecours}{%
  \par\vfill\nopagebreak
  \begin{center}\popsquiggle{poporange}\end{center}\vspace{4pt}
  \signatureonbuch
}
\AtBeginDocument{\def\llbracket{\mathopen{[\mkern-3.5mu[}}\def\rrbracket{\mathclose{]\mkern-3.5mu]}}} % intervalles d'entiers (glyphe ⟦ absent de la police)
% Glyphes absents de Fira Math : redéfinis à partir de glyphes disponibles
\AtBeginDocument{%
  \def\setminus{\mathbin{\backslash}}\let\smallsetminus\setminus
  \def\vdots{\mathord{\vbox{\baselineskip3.2pt\lineskiplimit0pt\kern4pt\hbox{.}\hbox{.}\hbox{.}}}}%
  \def\ddots{\mathinner{\mkern1mu\raise6.5pt\hbox{.}\mkern2mu\raise3.6pt\hbox{.}\mkern2mu\raise0.7pt\hbox{.}\mkern1mu}}%
  \def\triangle{\mathord{\scalebox{0.9}{$\symup{\Delta}$}}}%
  \def\nabla{\mathord{\raisebox{\depth}{\scalebox{1}[-1]{$\symup{\Delta}$}}}}%
  \def\checkmark{\popcheck{popdark}}%
  \def\star{\mathbin{\ast}}\def\bigstar{\mathord{\ast}}%
  \def\diamond{\mathbin{\scalebox{0.6}{$\Diamond$}}}%
  \def\bigcup{\mathop{\mathchoice{\raisebox{-0.3em}{\scalebox{1.7}{$\cup$}}}{\scalebox{1.25}{$\cup$}}{\cup}{\cup}}\displaylimits}%
  \def\bigcap{\mathop{\mathchoice{\raisebox{-0.3em}{\scalebox{1.7}{$\cap$}}}{\scalebox{1.25}{$\cap$}}{\cap}{\cap}}\displaylimits}%
  \def\lesseqqgtr{\mathrel{\lessgtr}}\def\bigoplus{\mathop{\oplus}\displaylimits}%
}
\providecommand{\ssi}{si et seulement si} % abréviation fréquente
\AtBeginDocument{\providecommand{\wideparen}{\overparen}} % arc de cercle

\begin{document}

\pagegarde{Lesson 47 — Vocabulary: Words and expresssions related to campaigning and voting for leaders}{Anglais}{Tle SH}{Chapter 4 : Citizenship and Human Rights}{EN47}{2 h}{Cette leçon de 2 heures propose une immersion lexicale complète dans le processus électoral, de la campagne au dépouillement. Elle structure le vocabulaire par champs sémantiques, traite les collocations indispensables, déjoue les faux-amis et entraîne à l'emploi en contexte pour réussir l'épreuve du Baccalauréat.
\vspace{4pt}
\begin{multicols}{2}\small
\begin{itemize}
\item Identifier et classer le vocabulaire clé des étapes du processus électoral (campagne, scrutin, résultats).
\item Employer correctement les collocations verbales et nominales usuelles (ex: 'cast a vote', 'run for office').
\item Différencier les faux-amis critiques (ex: 'candidate' vs 'candidat', 'scrutin' vs 'scrutiny').
\item Manipuler les familles de mots (vote, voter, voting ; elect, election, elective, electorate).
\item Prononcer avec l'accent tonique correct les termes longs (campaigning, constituency, referendum).
\item Réinvestir le lexique dans un dialogue simulant un débat électoral ou un reportage de bureau de vote.
\end{itemize}\end{multicols}}

\begin{sommaire}
\begin{multicols}{2}
\begin{enumerate}
\item 1. Le processus électoral : Chronologie et Acteurs (Word Families \& Mindmap)
\item 2. La Campagne : Discours, Promesses et Médias (Collocations \& Idioms)
\item 3. Le Jour du Scrutin : Procédures, Matériel et Verbes d'Action (Voting Process)
\item 4. Après le Vote : Résultats, Contentieux et Installation (Post-Election Lexicon)
\item 5. Atelier d'Intégration : Production Orale \& Écrite (Exam Practice)
\item Activité d'intégration
\item Méthodes et exercices résolus
\item Exercices
\item Corrigés des exercices
\item Fiche bilan
\end{enumerate}
\end{multicols}
\end{sommaire}

\begin{objectifs}
\begin{itemize}
\item Identifier et classer le vocabulaire clé des étapes du processus électoral (campagne, scrutin, résultats).
\item Employer correctement les collocations verbales et nominales usuelles (ex: 'cast a vote', 'run for office').
\item Différencier les faux-amis critiques (ex: 'candidate' vs 'candidat', 'scrutin' vs 'scrutiny').
\item Manipuler les familles de mots (vote, voter, voting ; elect, election, elective, electorate).
\item Prononcer avec l'accent tonique correct les termes longs (campaigning, constituency, referendum).
\item Réinvestir le lexique dans un dialogue simulant un débat électoral ou un reportage de bureau de vote.
\item Rédiger un court article d'opinion ou une lettre ouverte en utilisant le registre formel approprié.
\item Analyser un extrait de discours politique ou un reportage médiatique pour en extraire le vocabulaire fonctionnel.
\end{itemize}
\end{objectifs}

\begin{prerequis}
\begin{itemize}
\item Connaissance du système politique de base (Président, Parlement, Maire) vue en Première (Citizenship).
\item Maîtrise de la formation des mots (préfixes/suffixes : re-, -tion, -ment, -or/-er).
\item Notions de voix passive (The bill was passed) et reported speech pour le reportage médiatique.
\item Vocabulaire de base de l'argumentation (agree, disagree, argue, claim, evidence).
\item Culture générale : différence entre 'first-past-the-post' (UK/US/Cameroon) et représentation proportionnelle.
\end{itemize}
\end{prerequis}

\newpage

\coursec{Le processus électoral : Chronologie et Acteurs (Word Families \& Mindmap)}

Imaginez la scène : c'est la veille des élections municipales à Bamenda ou à Yaoundé. Les rues bourdonnent au rythme des \eng{campaign rallies}, les \eng{manifestos} sont distribués aux carrefours, et chacun scrute le \eng{ballot paper} en se demandant : \eng{Who will get the mandate?} Que vous soyez un \eng{voter} à Buea, un \eng{polling agent} à Lagos ou un \eng{constituent} à Londres, le langage de la démocratie a ses codes précis. Maîtriser ce lexique, c'est s'armer pour comprendre les enjeux, débattre sans ambiguïté et exercer sa citoyenneté pleinement, du \eng{grassroots} jusqu'au \eng{parliament}.

Cette première section pose les fondations : nous allons structurer le temps électoral en étapes claires, identifier les acteurs avec leurs termes techniques exacts, et construire des familles de mots pour multiplier votre pouvoir d'expression.

\courssub{La frise chronologique : De la pré-campagne à l'investiture}

Tout processus électoral suit une séquence logique immuable. La comprendre permet de situer n'importe quel mot de vocabulaire dans son contexte temporel. Observons la frise ci-dessous, qui va de la décision de se présenter jusqu'à la prise de fonction.

\begin{popfigure}
\begin{tikzpicture}[
    node distance=1.2cm and 1.5cm,
    every node/.style={font=\small, align=center, rounded corners=4pt, inner sep=6pt},
    phase/.style={draw=popblue, thick, fill=popblueL, text width=2.8cm, minimum height=1.8cm},
    arrow/.style={-Stealth, thick, popdark}
]
% Nodes
\node[phase, fill=popblue!10, align=center] (A){\textbf{Pre-Campaign}\\\eng{Nomination / Primaries}\\\eng{Registration}};
\node[phase, right=of A, align=center] (B){\textbf{Campaign Trail}\\\eng{Rallies / Debates}\\\eng{Canvassing / Manifestos}};
\node[phase, right=of B, align=center] (C){\textbf{Election Day}\\\eng{Polling Stations}\\\eng{Casting Ballots}};
\node[phase, right=of C, align=center] (D){\textbf{Counting / Results}\\\eng{Tallying}\\\eng{Concession / Victory}};
\node[phase, fill=popgreenL, right=of D, align=center] (E){\textbf{Inauguration}\\\eng{Swearing-in Ceremony}\\\eng{Mandate Begins}};

% Arrows
\draw[arrow] (A) -- (B) node[midway, above, font=\tiny] {\eng{Official Launch}};
\draw[arrow] (B) -- (C) node[midway, above, font=\tiny] {\eng{Silence Period}};
\draw[arrow] (C) -- (D) node[midway, above, font=\tiny] {\eng{Close of Polls}};
\draw[arrow] (D) -- (E) node[midway, above, font=\tiny] {\eng{Certification}};
\end{tikzpicture}
\legende{Frise chronologique du processus électoral type (modèle Westminster / Cameroun)}
\end{popfigure}

\begin{aretenir}[Les 5 étapes clés et leur vocabulaire associé]
\begin{enumerate}
    \item \textbf{Pre-Campaign (Pré-campagne)} : \eng{Nomination} (investiture), \eng{Primaries} (primaires), \eng{Voter registration} (inscription sur les listes électorales), \eng{Manifesto drafting} (rédaction du programme).
    \item \textbf{Campaign Trail (Campagne active)} : \eng{Rally} (meeting), \eng{Canvassing} (porte-à-porte / démarchage), \eng{Debate} (débat), \eng{Endorsement} (soutien officiel), \eng{Swing voter} (électeur indécis).
    \item \textbf{Election Day (Jour du scrutin)} : \eng{Polling station} (bureau de vote), \eng{Ballot paper} (bulletin de vote), \eng{Booth} (isoloir), \eng{Queue} (file d'attente), \eng{Indelible ink} (encre indélébile).
    \item \textbf{Counting / Results (Dépouillement / Résultats)} : \eng{Tallying} (comptage), \eng{Turnout} (taux de participation), \eng{Landslide victory} (victoire écrasante), \eng{Concede defeat} (reconnaître sa défaite).
    \item \textbf{Inauguration (Investiture)} : \eng{Swearing-in} (prestation de serment), \eng{Oath of office} (serment), \eng{Mandate} (mandat), \eng{Transition of power} (transition pacifique).
\end{enumerate}
\end{aretenir}

\courssub{Les Acteurs clés : Identification et Familles de mots}

Le vocabulaire des acteurs distingue soigneusement la fonction, la personne et le groupe. C'est ici que les \textbf{familles de mots} (word families) deviennent votre meilleur outil pour passer du singulier au collectif, du verbe au nom.

\begin{definition}[Electorate (n.)]
The body of all people entitled to vote in an election. \textit{Ne pas confondre avec 'électeur' (individuel) = \eng{Voter} / \eng{Elector}.}
\eng{The electorate in the North West region turned out in large numbers.} « L'électorat de la région du Nord-Ouest s'est déplacé en grand nombre. »
\end{definition}

\begin{propriete}[Word Family Table : Acteurs et Processus]
\begin{center}
\begin{tabularx}{\linewidth}{|X|X|X|X|}
\hline
\rowcolor{popblueL}
\textbf{Verb (Action)} & \textbf{Noun (Concept/Event)} & \textbf{Noun (Person/Agent)} & \textbf{Adjective} \\
\hline
\eng{Elect} & \eng{Election} & \eng{Elector} / \eng{Candidate} & \eng{Electoral} \\
\hline
\eng{Vote} & \eng{Vote} / \eng{Voting} & \eng{Voter} & \eng{Voting} (adj) \\
\hline
\eng{Campaign} & \eng{Campaign} & \eng{Campaigner} & \eng{Campaign} (adj) \\
\hline
\eng{Nominate} & \eng{Nomination} & \eng{Nominee} & \eng{Nominative} \\
\hline
\eng{Represent} & \eng{Representation} & \eng{Representative} & \eng{Representative} \\
\hline
\eng{Constituent} (rare as verb) & \eng{Constituency} & \eng{Constituent} & \eng{Constituent} (adj) \\
\hline
\eng{Observe} & \eng{Observation} & \eng{Observer} & \eng{Observational} \\
\hline
\end{tabularx}
\end{center}
\textbf{Note pédagogique} : Remarquez que \eng{Candidate} n'a pas de verbe direct *to candidate* (on dit \eng{to run for office} ou \eng{to stand for election}). \eng{Constituent} fonctionne comme nom (l'électeur) et adjectif (élément constitutif), le verbe étant \eng{constitute}.
\end{propriete}

\begin{exemplebox}[Acteurs en contexte : Exemples traduits]
\begin{itemize}
    \item \eng{The \textbf{incumbent} mayor decided to \textbf{run for a second term}.} « Le maire \textbf{sortant} a décidé de \textbf{briguer un second mandat}. »
    \item \eng{His main \textbf{challenger} is a former \textbf{MP} (Member of Parliament).} « Son principal \textbf{adversaire / challenger} est un ancien \textbf{député}. »
    \item \eng{Each candidate chooses a \textbf{running mate} for the ticket.} « Chaque candidat choisit un \textbf{colistier} (vice-président / adjoint) pour le ticket. »
    \item \eng{International \textbf{observers} monitored the \textbf{polling stations}.} « Des \textbf{observateurs} internationaux ont surveillé les \textbf{bureaux de vote}. »
    \item \eng{The \textbf{polling staff} arrived early to set up the \textbf{booths}.} « Le \textbf{personnel de bureau de vote} est arrivé tôt pour installer les \textbf{isoloirs}. »
\end{itemize}
\end{exemplebox}

\courssub{Les Actions du processus : Verbes, Collocations et Expressions}

Au-delà des noms d'acteurs, ce sont les \textbf{verbes d'action} et leurs \textbf{collocations} (associations naturelles de mots) qui font la fluidité. En français, on « fait campagne » ; en anglais, on \eng{runs a campaign} ou on \eng{campaigns}. On ne « fait » pas une élection (\eng{*make an election}), on \eng{holds an election}.

\begin{textebox}[The Campaign Trail]
The \textbf{incumbent} President launched his \textbf{re-election bid} at a massive \textbf{rally} in the capital. His main \textbf{challenger} promised to \textbf{tackle} unemployment and \textbf{root out} corruption. Both candidates \textbf{canvassed} rural \textbf{constituencies} hoping to \textbf{swing} undecided \textbf{voters}.
\end{textebox}

\begin{center}
\begin{tabularx}{\linewidth}{|X|X|X|}
\hline
\rowcolor{popblueL}
\textbf{Action Verb / Expression} & \textbf{Collocation Naturelle} & \textbf{Traduction / Contexte} \\
\hline
\eng{To campaign} & \eng{to campaign vigorously / door-to-door} & Faire campagne (activement / porte-à-porte) \\
\hline
\eng{To canvass} & \eng{to canvass voters / support} & Démarcher, faire du porte-à-porte pour obtenir des voix \\
\hline
\eng{To endorse} & \eng{to endorse a candidate / a policy} & Soutenir officiellement, donner son parrainage \\
\hline
\eng{To nominate} & \eng{to nominate sb for a post} & Investir, désigner officiellement un candidat \\
\hline
\eng{To register} & \eng{to register to vote / on the electoral roll} & S'inscrire (sur les listes électorales) \\
\hline
\eng{To cast a ballot} & \eng{to cast one's ballot / vote} & Déposer son bulletin dans l'urne (formel) \\
\hline
\eng{To concede (defeat)} & \eng{to concede defeat / the election} & Reconnaître sa défaite (geste démocratique fort) \\
\hline
\eng{To swear in} & \eng{to swear in the President} & Investir, faire prêter serment (passif: \eng{be sworn in}) \\
\hline
\end{tabularx}
\end{center}

\begin{exemplebox}[Collocations en phrases complètes]
\begin{itemize}
    \item \eng{The party \textbf{held a primary} to \textbf{choose its nominee}.} « Le parti a \textbf{organisé une primaire} pour \textbf{désigner son candidat investi}. »
    \item \eng{Voters \textbf{turned out} in record numbers despite the rain.} « Les électeurs \textbf{se sont déplacés / ont voté} en nombre record malgré la pluie. » (\eng{Turnout} = participation).
    \item \eng{The opposition leader \textbf{called for a recount} after the \textbf{narrow margin}.} « Le chef de l'opposition a \textbf{demandé un recomptage} après l'\textbf{écart serré}. »
    \item \eng{The new President was \textbf{sworn in} before the \textbf{Chief Justice}.} « Le nouveau Président a \textbf{prêté serment} devant le \textbf{Président de la Cour Suprême}. »
\end{itemize}
\end{exemplebox}

\courssub{Focus Prononciation et Orthographe : Pièges classiques}

La maîtrise orale et écrite passe par l'attention aux détails phonétiques et graphiques.

\begin{propriete}[Accent tonique (Word Stress) : Mots pièges]
En anglais, l'accent tonique change souvent par rapport au français. Marquez-le fort ( \cle{' } ).
\begin{center}
\begin{tabular}{|l|l|l|}
\hline
\rowcolor{popblueL}
\textbf{Mot} & \textbf{Phonétique simplifiée} & \textbf{Erreur fréquente (à la française)} \\
\hline
\eng{Candidate} & \eng{'kæn.dɪ.deɪt} → « \textbf{KAN}-di-deït » & *kan-di-\textbf{DAÏT} \\
\hline
\eng{Constituency} & \eng{kən'stɪt.ju.ən.si} → « kon-\textbf{STI}-tchou-en-si » & *cons-ti-\textbf{TUAN}-si \\
\hline
\eng{Referendum} & \eng{.ref.ə'ren.dəm} → « re-fé-\textbf{REN}-dème » & *ré-fé-ren-\textbf{DUM} \\
\hline
\eng{Campaigning} & \eng{kæm'peɪ.nɪŋ} → « kam-\textbf{PEI}-ning » & *\textbf{CAM}-pai-ning \\
\hline
\eng{Electoral} & \eng{ɪ'lek.tər.əl} → « i-\textbf{LEK}-te-ral » & *é-\textbf{LEC}-to-ral \\
\hline
\end{tabular}
\end{center}
\end{propriete}

\begin{attention}[Orthographe : Doubles consonnes et Terminaisons]
\begin{itemize}
    \item \textbf{Double L / Double T} : \eng{Campaign} (un seul \textbf{n}, pas *campaing), \eng{Ballot} (double \textbf{l}, un \textbf{t}), \eng{Polling} (double \textbf{l} car \eng{poll} + \eng{-ing}), \eng{Controlled} (double \textbf{l} en anglais britannique \eng{controlled} mais attention \eng{control} $\to$ \eng{controlling}).
    \item \textbf{-ment vs -ance/-ence} : \eng{Government} (gouvernement) $\neq$ \eng{Governance} (gouvernance / mode de gestion). \eng{Parliament} (parlement) $\to$ \eng{Parliamentary} (parlementaire).
    \item \textbf{Suffixes agents} : \eng{Voter} (électeur), \eng{Campaigner} (militant), \eng{Observer} (observateur), \eng{Candidate} (candidat — pas de suffixe \eng{-er}).
\end{itemize}
\end{attention}

\courssub{Atelier Méthodologique : La technique « Cluster \& Story »}

\begin{methode}[Méthode 'Cluster \& Story' pour mémoriser le lexique électoral]
\begin{enumerate}
    \item \textbf{Groupez (Cluster)} : Créez 3 colonnes sur une feuille : \eng{BEFORE} (Pre-campaign), \eng{DURING} (Campaign + Election Day), \eng{AFTER} (Results + Inauguration). Placez chaque mot appris dans la bonne colonne.
    \item \textbf{Reliez (Collocations)} : Tracez des flèches entre les verbes et leurs noms partenaires (ex: \eng{hold} $\to$ \eng{an election}, \eng{cast} $\to$ \eng{a ballot}, \eng{concede} $\to$ \eng{defeat}).
    \item \textbf{Racontez (Story)} : Écrivez un mini-récit de 5 phrases (une par étape de la frise) en utilisant \textbf{au moins 3 mots} de votre liste par phrase.
    \item \textbf{Vérifiez} : Relisez-vous à voix haute en marquant l'accent tonique. Corrigez les faux-amis (\eng{candidate FOR}, pas *candidate to).
\end{enumerate}
\eng{Example story start: "During the PRE-CAMPAIGN, the NOMINEE filed his papers. During the CAMPAIGN, he CANVASSED rural CONSTITUENCIES..."}
\end{methode}

\begin{savaistu}[Le vocabulaire électoral au Cameroun anglophone]
Au Cameroun anglophone (régions du Nord-Ouest et du Sud-Ouest), le système électoral calqué sur le modèle britannique (Westminster) impose des termes spécifiques :
\begin{itemize}
    \item \eng{Ward} : Circonscription municipale (pour les élections locales / \eng{Council elections}). Équivalent francophone : \textbf{Arrondissement / Commune}.
    \item \eng{Constituency} : Circonscription législative (pour l'Assemblée Nationale / \eng{Parliamentary elections}).
    \item \eng{Councillor} : Conseiller municipal (élu au \eng{Council}). Ne pas confondre avec \eng{Counselor} (conseiller juridique / psychologue).
    \item \eng{Returning Officer} : Officiel responsable de la supervision du scrutin dans une circonscription (annonce les résultats).
\end{itemize}
Cette terminologie diffère parfois du français officiel camerounais (\eng{Sub-division}, \eng{Division}) : soyez précis selon le contexte (local vs national).
\end{savaistu}

\begin{exoresolu}[Entraînement : Complétion et Choix lexical]
\textbf{Énoncé :} Complétez le texte suivant avec les mots de la liste ci-dessous. Attention, deux mots sont en trop. Mettez le verbe à la forme correcte si nécessaire.

\eng{List: incumbent, challenger, constituency, canvass, endorse, nominate, ballot, concede, swear in, turnout, ward, manifesto.}

\eng{1. The \_\_\_\_\_\_\_\_\_\_\_\_ President addressed the nation yesterday.}
\eng{2. His main \_\_\_\_\_\_\_\_\_\_\_\_ promised to reduce taxes.}
\eng{3. Party delegates will \_\_\_\_\_\_\_\_\_\_\_\_ her as the official candidate tomorrow.}
\eng{4. Candidates must \_\_\_\_\_\_\_\_\_\_\_\_ every \_\_\_\_\_\_\_\_\_\_\_\_ to meet voters.}
\eng{5. The \_\_\_\_\_\_\_\_\_\_\_\_ was high, reaching 75\% of registered voters.}
\eng{6. She refused to \_\_\_\_\_\_\_\_\_\_\_\_ defeat until the final count was verified.}
\eng{7. The Chief Justice will \_\_\_\_\_\_\_\_\_\_\_\_ the winner next Monday.}

\tcblower
\textbf{Solution détaillée :}
\begin{enumerate}
    \item \eng{incumbent} (Adjectif placé avant le nom, signifie « sortant »). \textit{Traduction : Le Président sortant...}
    \item \eng{challenger} (Nom, l'adversaire principal). \textit{Traduction : Son principal adversaire...}
    \item \eng{nominate} $\to$ \eng{nominate} (Verbe à l'infinitif après modale \eng{will}). \textit{Traduction : ...désigneront / investiront...}
    \item \eng{canvass} (Verbe à l'infinitif après \eng{must}) + \eng{ward} ou \eng{constituency} (Nom, objet de l'action). \textit{Choix contextuel :} \eng{ward} (local) ou \eng{constituency} (législatif). Les deux sont acceptés. \textit{Traduction : ...doivent faire du porte-à-porte dans chaque quartier/circonscription...}
    \item \eng{turnout} (Nom, sujet de \eng{was}). \textit{Traduction : La participation a été élevée...}
    \item \eng{concede} (Verbe à l'infinitif après \eng{refused to}). \textit{Traduction : Elle a refusé de reconnaître sa défaite...}
    \item \eng{swear in} (Verbe phrasal à l'infinitif après \eng{will}). \textit{Traduction : Le Président de la Cour Suprême investira / fera prêter serment au vainqueur...}
\end{enumerate}
\textbf{Mots non utilisés :} \eng{ballot, endorse, manifesto}.
\end{exoresolu}

\begin{exercice}[Entraînement personnel : Production écrite guidée]
Rédigez un paragraphe d'une cinquantaine de mots en anglais pour décrire la scène d'un bureau de vote le jour de l'élection (Election Day). Utilisez \textbf{au moins 6 mots ou expressions} vus dans cette section (acteurs, verbes, noms de lieu). Soulignez-les.
\textit{Indication : Pensez à \eng{polling station, queue, ballot paper, booth, polling staff, cast a vote, observer, indelible ink}.}
\end{exercice}

\begin{corrige}[Exercice 1 - Proposition de production]
\eng{At the \textbf{polling station}, a long \textbf{queue} formed early. The \textbf{polling staff} checked IDs and handed out \textbf{ballot papers}. Voters entered the \textbf{booth}, marked their choice, and \textbf{cast their ballots} into the box. An \textbf{observer} watched the process. Finally, voters dipped their fingers in \textbf{indelible ink}.}
\textit{Traduction : Au bureau de vote, une longue file s'est formée tôt. Le personnel a vérifié les identités et distribué les bulletins. Les électeurs sont entrés dans l'isoloir, ont coché leur choix et ont déposé leurs bulletins dans l'urne. Un observateur surveillait le processus. Enfin, les électeurs ont trempé leur doigt dans l'encre indélébile.}
\end{corrige}

\coursec{La Campagne : Discours, Promesses et Médias (Collocations \& Idioms)}

Dans la section précédente, nous avons posé le décor : la chronologie d'une élection (de la convocation du corps électoral à la proclamation des résultats) et ses acteurs (candidates, voters, polling officials, observers). Mais avant le jour du scrutin, il se passe quelque chose d'essentiel : la \cle{campagne électorale}. C'est le moment où les candidats parlent, promettent, débattent — et où les médias commentent. Or la campagne possède en anglais un vocabulaire très codifié : des \cle{collocations} (associations verbe + nom qu'on ne peut pas traduire mot à mot), des \cle{expressions idiomatiques} imagées, et un lexique médiatique propre. C'est l'objet de cette section.

\courssub{Observer d'abord : un extrait de presse}

Lisons d'abord ce court article, tel qu'on pourrait le lire dans l'édition anglaise d'un journal camerounais la veille d'un scrutin.

\begin{textebox}[Media Coverage Extract — “A Race Too Close to Call”]
With only three days to go before polling day, the campaign has entered its final stretch. The ruling party held a massive rally in Yaoundé on Saturday, where its candidate addressed an enthusiastic crowd of supporters and unveiled the final version of the party manifesto. “We pledge to create jobs for the youth and to fight corruption,” he vowed, promising to commit more resources to rural roads.

The latest opinion polls suggest a neck-and-neck race. The incumbent enjoys high name recognition, but the challenger is gaining traction among young, first-time voters. She was once considered a dark horse, but after a strong performance in last week's televised debate — where she rebutted her opponent's claims and refused to dodge difficult questions — many analysts now call her the front-runner.

The campaign has not been clean. Both sides accuse each other of mudslinging and of running a smear campaign on social media, where fake news spreads faster than fact-checking can correct it. Analysts also warn of potential vote-buying and ballot stuffing in remote polling stations. International observers have been deployed to monitor the process, and journalists are demanding fair airtime and unbiased coverage for all candidates. Whoever wins, it will certainly not be a landslide victory.
\end{textebox}

\textbf{Glossaire (mots difficiles) :}

\begin{center}
\begin{tabularx}{\linewidth}{|l|l|X|}
\hline
\rowcolor{popblueL} \textbf{Mot anglais} & \textbf{Nature} & \textbf{Traduction} \\ \hline
polling day & n. & jour du scrutin \\ \hline
to hold a rally & v. + n. & tenir un meeting \\ \hline
to address a crowd & v. + n. & s'adresser à une foule \\ \hline
to unveil a manifesto & v. + n. & dévoiler un programme électoral \\ \hline
to pledge / to vow & v. & s'engager / jurer solennellement \\ \hline
incumbent & n. / adj. & sortant (titulaire du poste) \\ \hline
challenger & n. & challenger, adversaire du sortant \\ \hline
to gain traction & expr. & gagner du terrain, séduire de plus en plus \\ \hline
to rebut & v. & réfuter \\ \hline
to dodge a question & expr. & esquiver une question \\ \hline
front-runner & n. & favori, candidat en tête \\ \hline
vote-buying & n. & achat de voix \\ \hline
ballot stuffing & n. & bourrage d'urnes \\ \hline
airtime & n. & temps d'antenne \\ \hline
\end{tabularx}
\end{center}

\textbf{Questions de compréhension (reading) :}
\begin{enumerate}
\item What did the ruling party candidate do at the rally in Yaoundé?
\item Why do analysts now call the challenger the front-runner?
\item Mention two forms of electoral fraud mentioned in the article.
\item What role do international observers play?
\end{enumerate}

\courssub{Les collocations fortes : verbe + nom}

\begin{definition}[Collocation]
Une \cle{collocation} est une association habituelle et « naturelle » de deux mots (souvent verbe + nom) qu'un anglophone attend ensemble. On dit \eng{hold a rally} et non \eng{*do a rally} ; \eng{launch a manifesto} et non \eng{*open a manifesto}. Les collocations s'apprennent comme des blocs, jamais mot à mot.
\end{definition}

Voici les collocations essentielles de la campagne électorale :

\begin{center}
\begin{tabularx}{\linewidth}{|l|X|X|}
\hline
\rowcolor{popblueL} \textbf{Collocation} & \textbf{Traduction} & \textbf{Exemple} \\ \hline
to launch a manifesto & lancer un programme & \eng{The party launched its manifesto in Douala.} \\ \hline
to unveil a platform & dévoiler un programme & \eng{She unveiled her platform on health care.} \\ \hline
to make pledges & faire des promesses & \eng{Candidates make pledges they rarely keep.} \\ \hline
to hold a rally & tenir un meeting & \eng{They held a rally at the Bamenda stadium.} \\ \hline
to address a crowd & s'adresser à une foule & \eng{He addressed a crowd of 5,000 supporters.} \\ \hline
to run a campaign & mener une campagne & \eng{She ran a clean, issue-based campaign.} \\ \hline
to whip up support & mobiliser les soutiens & \eng{He toured the region to whip up support.} \\ \hline
to cast a vote & voter (déposer son bulletin) & \eng{Millions cast their votes peacefully.} \\ \hline
to concede defeat & concéder la défaite & \eng{The loser conceded defeat at midnight.} \\ \hline
to win a landslide & remporter une victoire écrasante & \eng{She won a landslide in her home region.} \\ \hline
\end{tabularx}
\end{center}

\begin{attention}[Calque : « faire campagne »]
Les francophones traduisent souvent « faire campagne » par \eng{*do campaign} ou \eng{*make a campaign} : c'est \textbf{faux}. Les formes correctes sont :
\begin{itemize}
\item \eng{to campaign} (verbe intransitif) : \eng{She campaigned in the North-West region.} → Elle a fait campagne dans la région du Nord-Ouest.
\item \eng{to run a campaign} : \eng{He ran a very aggressive campaign.} → Il a mené une campagne très agressive.
\item \eng{to go on the campaign trail} : \eng{The candidates are on the campaign trail.} → Les candidats sont sur le terrain de la campagne.
\end{itemize}
De même, « promesses » se dit bien \eng{promises}, mais dans un registre soutenu on préfère \eng{manifesto pledges} (engagements du programme). Et attention : « programme » électoral = \eng{manifesto} (UK) ou \eng{platform} (US), jamais \eng{*program} seul dans ce sens politique précis.
\end{attention}

\begin{popfigure}
\begin{tikzpicture}[node distance=0.9cm and 1.6cm,
verb/.style={rectangle, rounded corners, draw=popblue, fill=popblueL, font=\small\bfseries, align=center},
noun/.style={rectangle, rounded corners, draw=poporange, fill=poporangeL, font=\small, align=center},
arr/.style={-{Stealth[length=2mm]}, thick, popdark}]
\node[verb] (v1) {launch};
\node[noun, right=of v1] (n1) {a manifesto};
\node[verb, below=of v1] (v2) {hold};
\node[noun, right=of v2] (n2) {a rally};
\node[verb, below=of v2] (v3) {cast};
\node[noun, right=of v3] (n3) {a vote};
\node[verb, below=of v3] (v4) {concede};
\node[noun, right=of v4] (n4) {defeat};
\node[verb, below=of v4] (v5) {win};
\node[noun, right=of v5] (n5) {a landslide};
\draw[arr] (v1) -- (n1);
\draw[arr] (v2) -- (n2);
\draw[arr] (v3) -- (n3);
\draw[arr] (v4) -- (n4);
\draw[arr] (v5) -- (n5);
\node[font=\small\itshape, popmuted, above=0.15cm of v1, xshift=1.2cm] {verb + noun = strong collocation};
\end{tikzpicture}
\legende{Les collocations fortes de la campagne : le verbe « appelle » son nom.}
\end{popfigure}

\courssub{Les expressions idiomatiques de la politique}

La presse politique anglophone adore les images. Beaucoup viennent des courses de chevaux (\eng{neck and neck}, \eng{dark horse}, \eng{front-runner}) ou de la nature (\eng{landslide}, \eng{grassroots}). Il faut connaître le sens littéral pour mémoriser le sens figuré.

\begin{propriete}[Idioms \& Metaphors Table]
\begin{center}
\begin{tabularx}{\linewidth}{|l|X|X|}
\hline
\rowcolor{popblueL} \textbf{Expression} & \textbf{Sens littéral} & \textbf{Sens figuré (contexte électoral)} \\ \hline
a landslide victory & glissement de terrain & victoire écrasante, massive \\ \hline
neck and neck & cou à cou (chevaux) & égalité parfaite dans les sondages \\ \hline
a dark horse & cheval inconnu & candidat outsider qui gagne soudain \\ \hline
mudslinging & lancer de boue & attaques personnelles, calomnies \\ \hline
grassroots support & racines de l'herbe & soutien populaire, à la base \\ \hline
to toe the party line & placer ses orteils sur une ligne (inspection militaire) & obéir strictement à la discipline du parti \\ \hline
a shoo-in & cheval qu'on fait entrer au box sans effort & grand favori, victoire assurée \\ \hline
to whip up support & fouetter (la crème) & mobiliser, galvaniser les soutiens \\ \hline
a smear campaign & tacher, salir & campagne de diffamation organisée \\ \hline
\end{tabularx}
\end{center}
\end{propriete}

\begin{exemplebox}[Exemples traduits]
\begin{enumerate}
\item \eng{The party unveiled its manifesto yesterday.} → Le parti a dévoilé son programme hier.
\item \eng{He pledged to reduce taxes if elected.} → Il s'est engagé à baisser les impôts s'il est élu.
\item \eng{Negative campaigning turns voters off.} → La campagne négative dégoûte les électeurs.
\item \eng{She is the clear front-runner.} → Elle est la favorite évidente.
\item \eng{The two candidates are neck and neck in the polls.} → Les deux candidats sont au coude-à-coude dans les sondages.
\item \eng{Nobody expected the dark horse to win the nomination.} → Personne ne s'attendait à ce que l'outsider remporte l'investiture.
\item \eng{He was once considered a shoo-in for re-election.} → Il était autrefois considéré comme un grand favori pour sa réélection.
\item \eng{MPs are expected to toe the party line.} → Les députés sont censés suivre la ligne du parti.
\end{enumerate}
\end{exemplebox}

\courssub{Registre de langue : manifesto ou platform ?}

\begin{aretenir}[Manifesto (UK) vs Platform (US)]
\begin{itemize}
\item \eng{manifesto} : terme britannique et Commonwealth — déclaration publique du programme d'un parti avant une élection. Prononciation : « ma-ni-fès-to ».
\item \eng{platform} : terme américain pour le même concept.
\item Au \textbf{Cameroun} (pays de tradition Commonwealth dans les régions anglophones), on emploie \eng{manifesto} ou \eng{party programme}. Dans un examen camerounais, \eng{manifesto} est le choix le plus sûr.
\end{itemize}
\end{aretenir}

\begin{definition}[Manifesto]
\eng{Manifesto} (n.) : a public declaration of policy and aims, especially one issued before an election by a political party. \emph{Synonyme américain : platform.} Exemple : \eng{The manifesto promises free primary education.} → Le programme promet l'école primaire gratuite.
\end{definition}

\courssub{Le vocabulaire des médias et du débat}

La campagne se joue aussi dans les médias. Voici le lexique indispensable pour commenter la couverture médiatique d'une élection :

\begin{center}
\begin{tabularx}{\linewidth}{|l|X|X|}
\hline
\rowcolor{popblueL} \textbf{Mot} & \textbf{Traduction} & \textbf{Exemple} \\ \hline
airtime & temps d'antenne & \eng{All candidates deserve equal airtime.} \\ \hline
coverage & couverture médiatique & \eng{The election received wide coverage.} \\ \hline
bias & parti pris, partialité & \eng{The newspaper was accused of bias.} \\ \hline
spin & présentation orientée & \eng{His team put a positive spin on the defeat.} \\ \hline
fake news & fausses informations & \eng{Fake news spreads fast on social media.} \\ \hline
fact-checking & vérification des faits & \eng{Fact-checking teams worked all night.} \\ \hline
opinion poll & sondage d'opinion & \eng{The latest opinion poll shows a tie.} \\ \hline
exit poll & sondage sortie des urnes & \eng{Exit polls predict a narrow victory.} \\ \hline
\end{tabularx}
\end{center}

Et dans les débats télévisés, les verbes d'action sont précis :

\begin{center}
\begin{tabularx}{\linewidth}{|l|X|X|}
\hline
\rowcolor{popblueL} \textbf{Verbe} & \textbf{Traduction} & \textbf{Nuance} \\ \hline
to debate & débattre & discuter publiquement \\ \hline
to rebut & réfuter & démontrer qu'un argument est faux \\ \hline
to clarify & clarifier & préciser sa position \\ \hline
to dodge a question & esquiver une question & éviter de répondre (péjoratif) \\ \hline
to pledge & s'engager & promesse ferme, formelle \\ \hline
to vow & jurer & promesse solennelle, émotionnelle \\ \hline
to commit to & s'engager à & suivi d'un nom ou d'un gérondif \\ \hline
\end{tabularx}
\end{center}

\begin{attention}[Pièges de construction]
\begin{itemize}
\item \eng{to commit to} est suivi d'un \textbf{nom ou d'un verbe en -ing} : \eng{She committed to reducing poverty.} (et non \eng{*to reduce}).
\item \eng{to pledge} et \eng{to vow} sont suivis de l'\textbf{infinitif} : \eng{He vowed to fight corruption.}
\item \eng{biased} se construit avec \eng{against} ou \eng{in favour of} : \eng{The channel is biased against the opposition.}
\item Ne confondez pas \eng{opinion poll} (avant le vote, sur les intentions) et \eng{exit poll} (le jour du vote, à la sortie des bureaux).
\end{itemize}
\end{attention}

\begin{savaistu}[Qu'est-ce qu'un « spin doctor » ?]
Un \eng{spin doctor} est un conseiller en communication politique qui présente les événements sous l'angle le plus favorable à son candidat — une forme de propagande douce. Le mot \eng{spin} désigne la « rotation » qu'on donne à une information, comme on donne de l'effet à une balle. Le terme est très courant dans la presse anglophone camerounaise (\emph{The Guardian Post}, l'édition anglaise de \emph{Cameroon Tribune}) pour décrire les stratèges de communication des partis pendant les campagnes.
\end{savaistu}

\courssub{Exercice résolu}

\begin{exoresolu}[Complétez avec le mot ou l'expression correct(e)]
Énoncé — Complete the sentences with: \emph{landslide, neck and neck, rebutted, airtime, manifesto, mudslinging, pledged, exit poll}.
\begin{enumerate}
\item The two candidates are \dots\ in the latest opinion poll: 48 \% each.
\item The president won a \dots\ victory with over 80 \% of the votes.
\item She \dots\ to build three new hospitals if elected.
\item During the debate, he calmly \dots\ all the accusations against him.
\item Voters are tired of \dots ; they want to hear real ideas.
\item The party launched its \dots\ at a press conference in Buea.
\item According to the \dots , the challenger is leading by two points.
\item Opposition parties complained about unfair \dots\ on state television.
\end{enumerate}
\tcblower
Solution détaillée :
\begin{enumerate}
\item \eng{neck and neck} — 48 \% chacun : égalité parfaite (image des courses de chevaux).
\item \eng{landslide} — plus de 80 \% : victoire écrasante.
\item \eng{pledged} — promesse formelle + infinitif (\eng{pledged to build}).
\item \eng{rebutted} — il a réfuté les accusations (prétérit régulier : rebut $\rightarrow$ rebutted, on double le \emph{t}).
\item \eng{mudslinging} — les électeurs en ont assez des attaques personnelles.
\item \eng{manifesto} — le programme électoral du parti (registre UK/Cameroun).
\item \eng{exit poll} — sondage réalisé à la sortie des bureaux de vote.
\item \eng{airtime} — temps d'antenne inégal à la télévision publique.
\end{enumerate}
\end{exoresolu}

\begin{exercice}[À vous de jouer]
\begin{enumerate}
\item Traduisez en anglais : « Les deux candidats sont au coude-à-coude, mais le sortant reste le favori. »
\item Trouvez trois collocations avec le nom \eng{campaign} (ex. \eng{run a campaign}).
\item Expliquez en anglais la différence entre un \eng{opinion poll} et un \eng{exit poll}.
\item Écrivez une phrase avec \eng{toe the party line} et une avec \eng{dark horse}.
\end{enumerate}
\end{exercice}

\begin{corrige}[Exercice « À vous de jouer »]
\begin{enumerate}
\item \eng{The two candidates are neck and neck, but the incumbent remains the front-runner.}
\item Par exemple : \eng{to run a campaign}, \eng{to launch a campaign}, \eng{a smear campaign} (ou \eng{an election campaign}, \eng{a negative campaign}).
\item \eng{An opinion poll measures voting intentions before polling day; an exit poll is conducted on polling day, as voters leave the polling station, to predict the result.}
\item Exemples : \eng{Backbench MPs are expected to toe the party line on the budget vote.} / \eng{The young mayor of a small town emerged as the dark horse of the presidential race.}
\end{enumerate}
\end{corrige}

\begin{aretenir}[L'essentiel de la section]
\begin{itemize}
\item Les collocations s'apprennent en blocs : \eng{launch a manifesto}, \eng{hold a rally}, \eng{address a crowd}, \eng{run a campaign} — jamais \eng{*do a campaign}.
\item Les idiomes politiques sont imagés : \eng{neck and neck} (égalité), \eng{landslide victory} (victoire écrasante), \eng{dark horse} (outsider), \eng{mudslinging} (calomnies), \eng{toe the party line} (discipline de parti).
\item \eng{Manifesto} (UK/Cameroun) vs \eng{platform} (US).
\item Médias : \eng{airtime}, \eng{coverage}, \eng{bias}, \eng{spin}, \eng{fake news}, \eng{fact-checking}, \eng{opinion poll} $\neq$ \eng{exit poll}.
\item Débat : \eng{rebut}, \eng{clarify}, \eng{dodge a question} ; promesses : \eng{pledge to + inf.}, \eng{vow to + inf.}, \eng{commit to + V-ing}.
\end{itemize}
\end{aretenir}

Maintenant que nous savons parler de la campagne, de ses promesses et de son traitement médiatique, passons au moment décisif : le jour du scrutin, ses procédures, son matériel et ses verbes d'action — l'objet de la section suivante.

\coursec{Le Jour du Scrutin : Procédures, Matériel et Verbes d'Action (Voting Process)}

Après des semaines de campagne — meetings, slogans, promesses et débats télévisés — vient le moment de vérité : \cle{Election Day}. Toute l'énergie déployée par les candidats se concentre désormais sur un geste simple mais solennel : celui de l'électeur qui dépose son bulletin dans l'urne. Cette section vous donne tout le lexique nécessaire pour décrire, à l'oral comme à l'écrit (notamment dans les sujets de type « process description » du Baccalauréat), le déroulement complet du jour du vote : le matériel utilisé, les acteurs présents, les verbes d'action et la terminologie des résultats.

\courssub{Observation : le script du président du bureau de vote}

Avant toute règle, observons la langue authentique. Voici le discours type qu'un \cle{Presiding Officer} (président de bureau de vote) adresse aux électeurs qui se présentent, par exemple à un bureau de vote de Buea ou de Bamenda.

\begin{textebox}[Presiding Officer's Script (Simulation)]
“Good morning. Please \textbf{join the queue}. Have your \textbf{Voter's Card} and \textbf{ID} ready. \textbf{Approach the table} when called. The \textbf{Polling Clerk} will \textbf{verify} your name on the \textbf{Electoral Roll}. You will receive a \textbf{stamped ballot paper}. Proceed to the \textbf{booth}, \textbf{mark} your choice clearly, \textbf{fold} it, and \textbf{deposit} it in the \textbf{box}. Dip your finger in the \textbf{indelible ink}. Thank you for voting.”
\end{textebox}

\textbf{Traduction.} « Bonjour. Veuillez vous mettre dans la file d'attente. Préparez votre carte d'électeur et votre pièce d'identité. Approchez-vous de la table quand on vous appelle. Le secrétaire de bureau vérifiera votre nom sur la liste électorale. Vous recevrez un bulletin de vote estampillé. Rendez-vous dans l'isoloir, marquez clairement votre choix, pliez le bulletin et déposez-le dans l'urne. Trempez votre doigt dans l'encre indélébile. Merci d'avoir voté. »

\textbf{Glossaire.}

\begin{center}
\begin{tabularx}{\linewidth}{|l|l|X|}
\hline
\rowcolor{popblueL} Mot anglais & Nature & Traduction française \\
\hline
to join the queue & loc. verb. & faire la queue, se mettre dans la file \\
\hline
Voter's Card & nom & carte d'électeur \\
\hline
ID (identity document) & nom & pièce d'identité \\
\hline
Polling Clerk & nom & secrétaire de bureau de vote \\
\hline
to verify & verbe & vérifier (prononcé « vé-ri-faï ») \\
\hline
Electoral Roll & nom & liste électorale \\
\hline
stamped ballot paper & nom & bulletin de vote estampillé (paraphé) \\
\hline
booth & nom & isoloir (prononcé « bouze ») \\
\hline
to mark & verbe & marquer, cocher \\
\hline
to fold & verbe & plier \\
\hline
to deposit & verbe & déposer \\
\hline
indelible ink & nom & encre indélébile (prononcé « in-dé-li-beul ») \\
\hline
\end{tabularx}
\end{center}

\textbf{Questions de compréhension.}
\begin{enumerate}
\item What two documents must a voter have ready? \emph{(The Voter's Card and an ID.)}
\item Who checks the voter's name, and on what document? \emph{(The Polling Clerk, on the Electoral Roll.)}
\item In what order does the voter handle the ballot paper? \emph{(Receive it, mark it, fold it, deposit it.)}
\item Why is indelible ink used? \emph{(To show that someone has already voted and prevent double voting.)}
\end{enumerate}

\courssub{Le matériel et les lieux du vote}

\begin{definition}[Voting equipment and places]
\begin{itemize}
\item \cle{ballot paper} (UK) / \cle{ballot} (US) : le bulletin de vote, le document sur lequel l'électeur exprime son choix.
\item \cle{ballot box} : l'urne (électorale).
\item \cle{polling booth} / \cle{voting booth} : l'isoloir, petit espace fermé où l'électeur vote en secret.
\item \cle{polling station} : le bureau de vote (le lieu, souvent une école ou une mairie).
\item \cle{indelible ink} : l'encre indélébile, qui ne s'efface pas pendant plusieurs jours.
\item \cle{voter's card} : la carte d'électeur ; \cle{ID card} : la carte nationale d'identité.
\end{itemize}
\end{definition}

\begin{exemplebox}[Le matériel en contexte]
\eng{Voters queued outside the polling station in Molyko, Buea, before dawn.} « Les électeurs faisaient la queue devant le bureau de vote de Molyko, à Buea, avant l'aube. »

\eng{Each voter received one ballot paper and entered the polling booth alone.} « Chaque électeur a reçu un bulletin et est entré seul dans l'isoloir. »

\eng{After casting her ballot, the old lady dipped her finger in the indelible ink.} « Après avoir déposé son bulletin, la vieille dame a trempé son doigt dans l'encre indélébile. »
\end{exemplebox}

\begin{attention}[Faux amis et calques du français]
\begin{itemize}
\item \textbf{« Scrutin » n'est pas \eng{scrutiny}.} Le \emph{scrutin} (le vote) se dit \eng{the vote}, \eng{the ballot} ou \eng{the election}. \eng{Scrutiny} signifie « examen minutieux » : \eng{The counting took place under close scrutiny.} « Le dépouillement a eu lieu sous étroite surveillance. »
\item \textbf{« Bulletin de vote » = \eng{ballot paper} (UK) / \eng{ballot} (US)}, jamais \eng{bulletin} (qui signifie « bulletin d'information » ou « bulletin scolaire »).
\item \textbf{« Isoloir » = \eng{polling booth} / \eng{voting booth}}, jamais \eng{isolator}.
\item \textbf{« Urne » = \eng{ballot box}} ; \eng{urn} seul désigne une urne funéraire ou une fontaine à thé.
\item \textbf{« Bureau de vote » = \eng{polling station}} (le lieu) ; \eng{polling office} est rare. Attention : \eng{polling station} désigne le bâtiment, pas le « bureau » au sens de meuble.
\end{itemize}
\end{attention}

\courssub{Les verbes de procédure : la chronologie du vote}

Le jour du scrutin suit une séquence rigide. Chaque étape correspond à un verbe précis que l'examinateur attend dans vos productions.

\begin{center}
\begin{tabularx}{\linewidth}{|l|X|X|}
\hline
\rowcolor{popblueL} Verbe anglais & Traduction & Exemple en contexte \\
\hline
to queue / to line up & faire la queue & \eng{Voters lined up at 6 a.m.} \\
\hline
to verify identity & vérifier l'identité & \eng{The clerk verified her identity.} \\
\hline
to issue a ballot paper & remettre un bulletin & \eng{A stamped ballot paper is issued.} \\
\hline
to mark the ballot & marquer le bulletin & \eng{He marked his ballot with an X.} \\
\hline
to fold the ballot paper & plier le bulletin & \eng{She folded the ballot paper twice.} \\
\hline
to cast / drop the ballot & déposer le bulletin & \eng{They cast their ballots peacefully.} \\
\hline
to spoil a ballot & rendre un bulletin nul & \eng{He spoilt his ballot by signing it.} \\
\hline
to count / tally the votes & compter les voix & \eng{The votes were tallied overnight.} \\
\hline
\end{tabularx}
\end{center}

\begin{popfigure}
\begin{tikzpicture}[node distance=0.45cm, every node/.style={font=\small}]
\node[draw, rounded corners, fill=popblue!30, align=center, text width=4.6cm] (a) {Arrival at Polling Station};
\node[draw, rounded corners, fill=popcream, align=center, text width=4.6cm, below=of a] (b) {Verification / Accreditation};
\node[draw, rounded corners, fill=popcream, align=center, text width=4.6cm, below=of b] (c) {Receive Ballot Paper};
\node[draw, rounded corners, fill=popcream, align=center, text width=4.6cm, below=of c] (d) {Enter Polling Booth};
\node[draw, rounded corners, fill=popcream, align=center, text width=4.6cm, below=of d] (e) {Mark Choice (Thumbprint / X)};
\node[draw, rounded corners, fill=popcream, align=center, text width=4.6cm, below=of e] (f) {Fold Ballot};
\node[draw, rounded corners, fill=popcream, align=center, text width=4.6cm, below=of f] (g) {Cast in Ballot Box};
\node[draw, rounded corners, fill=popcream, align=center, text width=4.6cm, below=of g] (h) {Ink Finger};
\node[draw, rounded corners, fill=popgreen!40, align=center, text width=4.6cm, below=of h] (i) {Exit};
\draw[-{Stealth}, thick, popdark] (a) -- (b);
\draw[-{Stealth}, thick, popdark] (b) -- (c);
\draw[-{Stealth}, thick, popdark] (c) -- (d);
\draw[-{Stealth}, thick, popdark] (d) -- (e);
\draw[-{Stealth}, thick, popdark] (e) -- (f);
\draw[-{Stealth}, thick, popdark] (f) -- (g);
\draw[-{Stealth}, thick, popdark] (g) -- (h);
\draw[-{Stealth}, thick, popdark] (h) -- (i);
\end{tikzpicture}
\legende{The voting process: nine steps from arrival to exit.}
\end{popfigure}

\begin{propriete}[Key verbs: transitive vs intransitive]
Certains verbes électoraux exigent un complément direct (transitifs), d'autres une préposition (intransitifs). C'est une source majeure d'erreurs pour les francophones.

\begin{center}
\begin{tabular}{|l|l|l|}
\hline
\rowcolor{popblueL} Verbe & Transitif (objet direct) & Intransitif / préposition \\
\hline
\textbf{vote} & \eng{vote for} a candidate & \eng{vote in} an election \\
\hline
\textbf{elect} & \eng{elect} a president (passif : \eng{be elected}) & --- \\
\hline
\textbf{cast} & \eng{cast} a vote / \eng{cast} a ballot & --- \\
\hline
\textbf{run} & \eng{run} a campaign & \eng{run for} office / \eng{run against} sb \\
\hline
\textbf{concede} & \eng{concede} defeat / \eng{concede} the election & --- \\
\hline
\end{tabular}
\end{center}

\eng{She was elected mayor of Douala.} « Elle a été élue maire de Douala. » (passif, sans \eng{by} si l'agent est évident.)

\eng{He ran for parliament but conceded defeat.} « Il s'est présenté aux législatives mais a reconnu sa défaite. »
\end{propriete}

\courssub{Les acteurs officiels du bureau de vote}

\begin{definition}[Officials at the polling station]
\begin{itemize}
\item \cle{Presiding Officer} : le président du bureau de vote ; il dirige les opérations et annonce la clôture (\eng{The Presiding Officer declared the poll closed.}).
\item \cle{Polling Clerk} : le secrétaire ; il vérifie les identités et remet les bulletins.
\item \cle{Polling Agent} : le représentant d'un parti ou d'un candidat dans le bureau ; il surveille la régularité du vote.
\item \cle{Observer} : l'observateur (national ou international) ; il ne participe pas, il constate et rédige un rapport.
\item \cle{Security personnel} : le personnel de sécurité (policiers, gendarmes) chargé du maintien de l'ordre.
\end{itemize}
\end{definition}

\begin{exemplebox}[Les rôles en contexte]
\eng{International observers praised the calm atmosphere at polling stations in Bamenda.} « Des observateurs internationaux ont salué le calme des bureaux de vote de Bamenda. »

\eng{A polling agent objected when a voter tried to photograph his ballot.} « Un représentant de parti a protesté lorsqu'un électeur a tenté de photographier son bulletin. »
\end{exemplebox}

\courssub{La terminologie des résultats}

\begin{definition}[Results vocabulary]
\begin{itemize}
\item \cle{turnout} : le taux de participation (nom invariable, sans pluriel).
\item \cle{valid votes} : les suffrages valablement exprimés.
\item \cle{invalid / spoilt votes} : les bulletins nuls.
\item \cle{majority} (ou \cle{absolute majority}) : la majorité absolue (plus de 50 \%) .
\item \cle{plurality} / \cle{relative majority} : la majorité relative (le plus grand nombre de voix, sans atteindre 50 \%) .
\item \cle{runoff} / \cle{second round} : le second tour, le ballottage.
\item \cle{declaration of results} : la proclamation des résultats.
\end{itemize}
\end{definition}

\begin{definition}[Spoilt ballot / Spoiled ballot]
(n.) A ballot paper that has been marked incorrectly, torn, or defaced, and therefore cannot be counted. \emph{Synonyme : invalid vote.} — « Bulletin nul : bulletin marqué incorrectement, déchiré ou gribouillé, qui ne peut donc pas être comptabilisé. »

\eng{Seventy spoilt ballots were found in the box.} « Soixante-dix bulletins nuls ont été trouvés dans l'urne. »
\end{definition}

\begin{attention}[Focus : le mot \eng{turnout}]
\eng{Turnout} est un \textbf{nom invariable} : on ne dit jamais \eng{turnouts}.

On dit \eng{high turnout} (forte participation) et \eng{low turnout} (faible participation).

Ne dites pas \eng{turnout rate} — c'est un pléonasme, car \eng{turnout} contient déjà l'idée de taux. Dites simplement :

\eng{Turnout was 70\%.} « Le taux de participation était de 70 \%. »

\eng{There was a low turnout in the rural areas.} « La participation a été faible dans les zones rurales. »
\end{attention}

\courssub{Méthode : décrire la procédure au passif impersonnel}

\begin{methode}[Process Description (Passif impersonnel)]
Pour décrire la procédure électorale dans une production écrite type Bac, le style attendu est le \textbf{passif au présent}, impersonnel et objectif :
\begin{enumerate}
\item Repérez les verbes d'action du processus (issue, verify, mark, fold, count, announce).
\item Mettez-les au passif présent : \eng{Ballot papers \textbf{are issued}. Votes \textbf{are counted}. Results \textbf{are announced}.}
\item Reliez les étapes par des connecteurs chronologiques : \eng{first, then, next, after that, finally}.
\item Évitez les sujets personnels (\eng{we, they}) : la procédure, pas les personnes, est au centre.
\end{enumerate}
Modèle : \eng{First, the voter's identity is verified. Then, a stamped ballot paper is issued. Next, the choice is marked in the booth. Finally, the ballot is cast and the finger is inked.}
\end{methode}

\begin{exoresolu}[Transformez au passif impersonnel]
Énoncé : réécrivez ces phrases actives au passif présent pour une description de procédure.
1. The clerk verifies the voter's name. 2. The officer stamps the ballot paper. 3. The voter folds the ballot. 4. Officials count the votes. 5. The council announces the results.
\tcblower
Solution détaillée : on déplace le complément d'objet en tête, on conjugue \eng{be} au présent, puis le participe passé.
1. \eng{The voter's name is verified (by the clerk).} 2. \eng{The ballot paper is stamped.} 3. \eng{The ballot is folded.} 4. \eng{The votes are counted.} 5. \eng{The results are announced.}
Remarque : le complément d'agent introduit par \eng{by} est généralement omis dans une description de procédure, car l'agent est évident ou sans importance.
\end{exoresolu}

\begin{exercice}[À vous de jouer]
1. Complétez avec le mot juste : \eng{ballot box, polling booth, turnout, spoilt, runoff}. (a) Voters mark their choice inside the \ldots. (b) A \ldots election is needed when no candidate wins an absolute majority. (c) All ballots are dropped into the \ldots. (d) A ballot with two crosses is \ldots. (e) \ldots was high in the cities but low in the villages.
2. Traduisez : « Les électeurs font la queue devant le bureau de vote. Leur identité est vérifiée, puis ils reçoivent un bulletin estampillé. »
\end{exercice}

\begin{corrige}[Exercice : À vous de jouer]
1. (a) \eng{polling booth} ; (b) \eng{runoff} ; (c) \eng{ballot box} ; (d) \eng{spoilt} ; (e) \eng{Turnout}.
2. \eng{Voters are queueing outside the polling station. Their identity is verified, then they are issued a stamped ballot paper.} (ou, à l'actif : \eng{Voters queue outside the polling station. Their identity is verified, then they receive a stamped ballot paper.})
\end{corrige}

\begin{savaistu}[L'encre indélébile en Afrique de l'Ouest et centrale]
Au Cameroun, l'encre indélébile (\eng{indelible ink}) est appliquée sur l'index gauche de l'électeur. Au Ghana et au Nigeria, c'est souvent le pouce qui est marqué. Le terme technique employé par les observateurs est \eng{finger marking} ou simplement \eng{inking}. Cette encre, qui résiste au lavage pendant plusieurs jours, empêche le vote multiple : \eng{The ink prevents double voting.} « L'encre empêche le double vote. »
\end{savaistu}

\begin{aretenir}[L'essentiel de la section]
\begin{itemize}
\item Le matériel : \eng{ballot paper} (bulletin), \eng{ballot box} (urne), \eng{polling booth} (isoloir), \eng{polling station} (bureau de vote), \eng{indelible ink} (encre indélébile).
\item La séquence : \eng{queue} $\rightarrow$ \eng{verify} $\rightarrow$ \eng{receive} $\rightarrow$ \eng{mark} $\rightarrow$ \eng{fold} $\rightarrow$ \eng{cast} $\rightarrow$ \eng{ink}.
\item Les acteurs : \eng{Presiding Officer, Polling Clerk, Polling Agent, Observer, Security personnel}.
\item Les résultats : \eng{turnout} (invariable, jamais \eng{turnout rate}), \eng{valid / spoilt votes}, \eng{absolute majority} vs \eng{relative majority}, \eng{runoff}.
\item À l'écrit, décrivez la procédure au \textbf{passif présent} : \eng{Votes are counted, results are announced.}
\end{itemize}
\end{aretenir}

\coursec{Après le Vote : Résultats, Contentieux et Installation (Post-Election Lexicon)}

Dans la section précédente, nous avons suivi les électeurs jusqu'au \eng{polling station} : ils ont vérifié leur nom sur la \eng{voters' register}, reçu leur \eng{ballot paper}, marqué leur choix dans le \eng{polling booth} et glissé le bulletin dans la \eng{ballot box}. Mais une élection ne s'arrête pas à la fermeture des bureaux de vote. C'est même là que tout commence : le dépouillement (\eng{counting}), l'annonce des résultats, d'éventuels recours devant les tribunaux, puis la prestation de serment du vainqueur. Cette quatrième section vous donne tout le lexique nécessaire pour raconter cette « après-élection » en anglais correct et naturel — un vocabulaire très fréquent dans les sujets du baccalauréat (compréhension de texte et essay writing).

\courssub{Observer d'abord : un discours de concession}

\begin{textebox}[Concession Speech Extract — a losing candidate addresses his supporters]
“My friends, the \textbf{returns} are in. The \textbf{electorate} has spoken. I have just called \textbf{President-elect} Mr Ayuk to \textbf{concede} this race and \textbf{congratulate} him. We \textbf{contested} this election vigorously but we \textbf{respect the verdict} of the people. Our observers reported some \textbf{irregularities} in three divisions, yet the \textbf{evidence} is not strong enough to \textbf{challenge} the overall outcome before the Constitutional Council. I therefore urge all my supporters to \textbf{accept the outcome} peacefully and to reject any call to violence. Our democracy is bigger than one man and one party. To those who feel \textbf{disenfranchised}, I say: your voice matters, and we will keep fighting for you through Parliament, not in the streets. The \textbf{transition of power} begins today, and I pledge my full cooperation so that it is smooth and dignified. Finally, I congratulate the Electoral Commission for organising a \textbf{free and fair} poll. God bless our nation, and God bless you all.”
\end{textebox}

\textbf{Glossaire (words in context) :}

\begin{center}
\begin{tabularx}{\linewidth}{|l|l|X|}
\hline
\rowcolor{popblueL} \textbf{English word} & \textbf{Nature} & \textbf{Français} \\
\hline
returns & n. pl. & résultats (au fur et à mesure du dépouillement) \\
\hline
electorate & n. & corps électoral, ensemble des électeurs \\
\hline
to concede & v. & concéder la défaite, admettre sa défaite \\
\hline
verdict & n. & verdict, décision (du peuple ou d'un tribunal) \\
\hline
irregularities & n. pl. & irrégularités (de scrutin) \\
\hline
to challenge & v. & contester (un résultat, une décision) \\
\hline
outcome & n. & issue, résultat final \\
\hline
disenfranchised & adj. & privé du droit de vote, exclu du processus électoral \\
\hline
transition of power & expr. & transition / passation du pouvoir \\
\hline
to pledge & v. & promettre solennellement, s'engager \\
\hline
\end{tabularx}
\end{center}

\textbf{Comprehension questions (oral practice) :}
\begin{enumerate}
\item What did the candidate do as soon as the returns came in?
\item Why did he decide not to file a petition?
\item What does he ask his supporters to do?
\end{enumerate}

\begin{definition}[President-elect]
\eng{President-elect} (n.) : the person who has won the presidential election but has not yet taken the oath of office. \emph{Ne traduisez pas simplement par « le président élu » sans nuance : en anglais, le mot insiste sur la période d'attente entre la victoire et l'investiture.} On dit aussi \eng{the governor-elect}, \eng{the mayor-elect}. Le trait d'union et la position du mot \eng{elect} APRÈS le nom sont obligatoires : \eng{the President-elect}, jamais \eng{the elect President}. Prononciation : « pré-zi-dent i-lèkt ».
\end{definition}

\courssub{Les verbes du résultat : win, lose, concede, contest}

Une fois les bulletins comptés (\eng{counted / tallied}), les résultats sont annoncés. Voici les verbes essentiels, avec leurs constructions exactes.

\begin{propriete}[Verbes de résultat : forme et emploi]
\begin{itemize}
\item \eng{to win (an election / a seat)} : gagner. \eng{to win by a landslide} = gagner par un raz-de-marée (écrasante majorité). \eng{to win by a narrow margin} = gagner d'une courte tête. Passé et participe passé : \eng{won, won}.
\item \eng{to lose (an election)} : perdre. Passé et participe passé : \eng{lost, lost}. On dit \eng{the losing candidate} (le candidat perdant).
\item \eng{to concede (defeat / the race)} : concéder la défaite publiquement. Construction : \eng{concede defeat}, \eng{concede the election to somebody}. Le discours correspondant est \eng{a concession speech}.
\item \eng{to contest} : attention, DEUX sens ! (1) \eng{to contest an election} = se porter candidat, briguer un mandat ; (2) \eng{to contest the results} = contester les résultats. Le contexte décide.
\item \eng{to challenge (the results / a decision)} : contester officiellement, souvent devant un tribunal. \eng{to lodge / file a challenge}.
\item \eng{to swear in (somebody)} : faire prêter serment (phrasal verb séparable : \eng{to swear the President in}). À la voix passive : \eng{to be sworn in} = prêter serment. Passé et participe passé : \eng{swore, sworn} (« souorn »).
\item \eng{to inaugurate} : investir officiellement ; le nom est \eng{the inauguration} (l'investiture).
\end{itemize}
\end{propriete}

\begin{exemplebox}[Exemples traduits]
\begin{enumerate}
\item \eng{The incumbent won by a landslide, with over seventy per cent of the votes.} → Le président sortant a gagné par un raz-de-marée, avec plus de soixante-dix pour cent des voix.
\item \eng{She lost the election by a narrow margin.} → Elle a perdu l'élection d'une courte tête.
\item \eng{The opposition filed a petition challenging the results.} → L'opposition a déposé un recours contestant les résultats.
\item \eng{The Constitutional Council upheld the election.} → Le Conseil Constitutionnel a validé (confirmé) l'élection.
\item \eng{The transition of power was peaceful.} → La transition du pouvoir s'est déroulée pacifiquement.
\item \eng{He was sworn in yesterday at the National Assembly.} → Il a prêté serment hier à l'Assemblée nationale.
\item \eng{Three candidates contested the municipal election in Bamenda.} → Trois candidats ont brigué (se sont présentés à) l'élection municipale de Bamenda.
\end{enumerate}
\end{exemplebox}

\begin{attention}[Piège : « contester » et le verbe \eng{to contest}]
Les francophones hésitent entre \eng{contest} et \eng{challenge}. Règle pratique : pour un recours juridique formel, les deux sont possibles (\eng{to contest / to challenge the results in court}), mais \eng{challenge} est plus fréquent dans la presse. En revanche, ne dites jamais \eng{to dispute an election} pour « se présenter à une élection » : \eng{dispute} signifie uniquement « contester, remettre en cause ». Et rappelez-vous : \eng{to win} ne se construit JAMAIS avec une personne comme objet direct — on \eng{wins an election}, on \eng{defeats / beats an opponent} (on ne « gagne » pas quelqu'un !).
\end{attention}

\courssub{Le contentieux électoral : recours, tribunaux et décisions}

Quand un candidat refuse les résultats, il saisit le juge. Voici le vocabulaire juridique indispensable.

\begin{definition}[Noms du contentieux électoral]
\begin{itemize}
\item \eng{a petition} (n.) : un recours, une requête (devant le juge électoral).
\item \eng{an appeal} (n.) : un appel (contre une décision) ; verbe : \eng{to appeal (against) a ruling}.
\item \eng{a tribunal / a court} : un tribunal ; \eng{the Constitutional Council} (Conseil Constitutionnel), \eng{the Supreme Court} (Cour suprême).
\item \eng{a ruling / a judgment} : une décision de justice, un arrêt.
\item \eng{evidence of fraud} : des preuves de fraude. Attention : \eng{evidence} est INCOMPTABLE en anglais — jamais \eng{evidences} ! On dit \eng{a piece of evidence} (une preuve).
\item \eng{irregularities} : irrégularités ; \eng{malpractice} : manœuvres frauduleuses, abus ; \eng{electoral fraud} : fraude électorale.
\item \eng{disenfranchisement} : privation du droit de vote (exclure des électeurs des listes, par exemple).
\end{itemize}
\end{definition}

\begin{propriete}[Legal Collocations — la justice électorale]
En anglais, chaque action juridique a SON verbe. Apprenez les couples verbe + nom par cœur :
\begin{center}
\begin{tabularx}{\linewidth}{|X|X|X|}
\hline
\rowcolor{popblueL} \textbf{Action (français)} & \textbf{Collocation correcte} & \textbf{Erreur fréquente (francophone)} \\
\hline
Déposer un recours & \eng{file / lodge a petition} & \eng{make / do a petition} \\
\hline
Juger une affaire & \eng{hear / adjudicate a case} & calque de « juger un recours » \\
\hline
Rendre une décision & \eng{deliver / hand down a ruling} & \eng{give a judgment} (lourd, peu naturel ici) \\
\hline
Valider les résultats & \eng{certify / validate the results} & \eng{confirm the results} (moins précis) \\
\hline
Confirmer (en appel) & \eng{uphold the election / the ruling} & \eng{maintain the election} \\
\hline
Annuler une élection & \eng{annul / nullify / overturn the election} & \eng{cancel the election} (= annuler un événement prévu) \\
\hline
Entrer en fonction & \eng{assume office / take office} & \eng{enter in function} (calque) \\
\hline
Passation du pouvoir & \eng{transition / transfer of power} & \eng{passation of power} (faux ami) \\
\hline
\end{tabularx}
\end{center}
\end{propriete}

\begin{attention}[\eng{Majority} vs \eng{Plurality} : un piège de civilisation]
\eng{Absolute majority} = plus de 50 \% des voix (majorité absolue). \eng{Relative majority} (GB) ou \eng{plurality} (US) = plus de voix que chacun des autres candidats, mais moins de 50 \% (majorité relative). Ne traduisez donc jamais « majorité relative » par \eng{plurality} dans un contexte britannique sans vérifier : les Britanniques disent \eng{relative majority}, les Américains \eng{plurality}. Autre piège : \eng{the majority} prend un verbe au singulier ou au pluriel selon le sens — \eng{The majority has decided} (le bloc) mais \eng{The majority of voters were young} (les individus).
\end{attention}

\courssub{Les institutions et la chronologie de l'après-vote}

\begin{definition}[Institutions et cérémonies]
\begin{itemize}
\item \eng{the Electoral Commission} : la commission électorale (au Cameroun : \eng{ELECAM}, Elections Cameroon). Elle \eng{organises the poll}, \eng{announces provisional results} et \eng{certifies the results}.
\item \eng{the Constitutional Council} : le Conseil Constitutionnel, qui \eng{hears petitions} et \eng{proclaims the final results}.
\item \eng{the Supreme Court} : la Cour suprême ; \eng{Parliament} : le Parlement (\eng{the National Assembly} + \eng{the Senate} au Cameroun).
\item \eng{the inauguration ceremony / the swearing-in ceremony} : la cérémonie d'investiture, où le vainqueur \eng{takes the oath of office} (prête serment). Prononciation de \eng{oath} : « euse » (un seul son, th sourd).
\end{itemize}
\end{definition}

\begin{popfigure}
\begin{tikzpicture}[
  node distance=0.9cm and 1.6cm,
  box/.style={draw=popblue, fill=popblueL, rounded corners, align=center, text width=3.1cm, minimum height=0.9cm, font=\small},
  dec/.style={draw=popgold, fill=popgoldL, diamond, aspect=2.2, align=center, font=\small, inner sep=1pt},
  good/.style={draw=popgreen, fill=popgreenL, rounded corners, align=center, text width=3.1cm, minimum height=0.9cm, font=\small},
  bad/.style={draw=poppink, fill=poppinkL, rounded corners, align=center, text width=3.1cm, minimum height=0.9cm, font=\small},
  arr/.style={-{Stealth[length=2.5mm]}, thick, popdark}
]
\node[box] (count) {Counting / Tallying of votes};
\node[dec, right=of count] (res) {Results};
\node[good, above right=0.2cm and 1.6cm of res] (conc) {Concession speech};
\node[box, below right=0.2cm and 1.6cm of res] (pet) {Petition to the Court};
\node[good, right=of conc] (swear) {Swearing-in / Inauguration};
\node[box, right=of pet] (rul) {Court ruling};
\node[bad, below=of rul] (rerun) {Re-run / By-election};
\draw[arr] (count) -- (res);
\draw[arr] (res) -- node[above, sloped, font=\scriptsize]{clear winner} (conc);
\draw[arr] (res) -- node[below, sloped, font=\scriptsize]{dispute} (pet);
\draw[arr] (conc) -- (swear);
\draw[arr] (pet) -- (rul);
\draw[arr] (rul) -- node[above, sloped, font=\scriptsize]{upheld} (swear);
\draw[arr] (rul) -- node[right, font=\scriptsize]{overturned} (rerun);
\end{tikzpicture}
\legende{From the ballot box to the inauguration: the post-election process.}
\end{popfigure}

Ce schéma résume la logique de l'après-vote : après le dépouillement (\eng{counting / tallying}), deux chemins sont possibles. S'il y a un vainqueur clair (\eng{a clear winner}), le perdant prononce son \eng{concession speech} et on passe à l'investiture. En cas de contestation (\eng{dispute}), un recours est déposé (\eng{a petition is filed}) ; le tribunal peut soit confirmer (\eng{uphold}) le résultat — et le vainqueur est investi — soit l'annuler (\eng{overturn / annul}), ce qui entraîne un nouveau scrutin (\eng{a re-run} pour une présidentielle, \eng{a by-election} pour un siège vacant au Parlement).

\courssub{Les adjectifs de légitimité : qualifier une élection}

Pour commenter un scrutin à l'oral ou dans un essay, il faut savoir le qualifier avec précision.

\begin{center}
\begin{tabularx}{\linewidth}{|X|X|X|}
\hline
\rowcolor{popblueL} \textbf{Adjectif / expression} & \textbf{Français} & \textbf{Exemple en contexte} \\
\hline
\eng{free and fair} & libre et équitable & \eng{Observers declared the poll free and fair.} \\
\hline
\eng{credible} & crédible & \eng{The results are credible.} \\
\hline
\eng{transparent} & transparent & \eng{a transparent counting process} \\
\hline
\eng{inclusive} & inclusif & \eng{an inclusive election} \\
\hline
\eng{peaceful} & pacifique, calme & \eng{The vote was peaceful nationwide.} \\
\hline
\eng{disputed / contested} & contesté & \eng{a disputed outcome} \\
\hline
\eng{marred by violence} & entaché de violences & \eng{The poll was marred by violence.} \\
\hline
\eng{rigged} & truqué & \eng{He claimed the election was rigged.} \\
\hline
\eng{boycotted} & boycotté & \eng{The opposition boycotted the poll.} \\
\hline
\end{tabularx}
\end{center}

\begin{attention}[Faux amis et calques à bannir]
\begin{itemize}
\item « Les élections se sont bien passées » ne se traduit pas par \eng{the elections passed well} mais par \eng{the elections went smoothly / were peaceful}.
\item \eng{sensible} est un faux ami : il signifie « raisonnable », pas « sensible ». Inutile ici, mais il revient souvent dans les essays sur les élections !
\item On dit \eng{voter turnout} (taux de participation) : \eng{Voter turnout was high / low.} — jamais \eng{the participation rate was high} dans un style journalistique naturel.
\item \eng{to rig an election} (truquer) est le verbe ; le nom est \eng{rigging} : \eng{allegations of vote-rigging}.
\end{itemize}
\end{attention}

\courssub{Exercice résolu et entraînement}

\begin{exoresolu}[Fill in the gaps with the correct post-election vocabulary]
Complétez avec : \eng{concede, upheld, sworn in, petition, turnout, landslide, transition, irregularities}.
\begin{enumerate}
\item The President-elect won by a \ldots\ldots\ldots, with 78 \% of the votes.
\item The defeated candidate refused to \ldots\ldots\ldots defeat and went to court.
\item His lawyers filed a \ldots\ldots\ldots alleging widespread \ldots\ldots\ldots in the North-West region.
\item The Constitutional Council \ldots\ldots\ldots the results after hearing the case.
\item The new governor was \ldots\ldots\ldots in a colourful ceremony in Buea.
\item Voter \ldots\ldots\ldots was low in urban areas.
\item The \ldots\ldots\ldots of power went smoothly.
\end{enumerate}
\tcblower
\textbf{Solution détaillée :}
\begin{enumerate}
\item \eng{landslide} — « gagner par un raz-de-marée » : \eng{to win by a landslide}.
\item \eng{concede} — collocation : \eng{to concede defeat} (concéder la défaite).
\item \eng{petition} ; \eng{irregularities} — collocations : \eng{to file a petition} ; \eng{alleging irregularities} (irrégularités alléguées).
\item \eng{upheld} — \eng{to uphold the results} = confirmer/valider les résultats (passé irrégulier : \eng{upheld}).
\item \eng{sworn in} — voix passive : \eng{to be sworn in} (prêter serment) ; participe passé irrégulier de \eng{swear} : \eng{sworn}.
\item \eng{turnout} — \eng{voter turnout} = taux de participation.
\item \eng{transition} — \eng{the transition of power} = la passation du pouvoir.
\end{enumerate}
\end{exoresolu}

\begin{exercice}[Your turn — sentence building]
Traduisez en anglais, en utilisant le lexique de cette section :
\begin{enumerate}
\item L'opposition a déposé un recours contestant les résultats.
\item Le tribunal a annulé l'élection et ordonné un nouveau scrutin.
\item Le taux de participation était élevé et le scrutin s'est déroulé pacifiquement.
\item Elle a prêté serment la semaine dernière et a immédiatement pris ses fonctions.
\end{enumerate}
\end{exercice}

\begin{corrige}[Exercice « Your turn »]
\begin{enumerate}
\item \eng{The opposition filed a petition challenging the results.}
\item \eng{The court annulled the election and ordered a re-run.}
\item \eng{Voter turnout was high and the poll was peaceful.}
\item \eng{She was sworn in last week and immediately assumed office.}
\end{enumerate}
\end{corrige}

\begin{savaistu}[The “lame duck”]
En anglais américain et britannique, on appelle \eng{a lame duck} (littéralement « un canard boiteux ») un dirigeant élu qui termine son mandat en attendant l'arrivée de son successeur : son pouvoir réel est affaibli, car tout le monde sait qu'il part. On parle de \eng{a lame-duck president} ou \eng{a lame-duck mayor}. Dans la presse anglophone camerounaise, l'expression apparaît parfois pour décrire les maires ou députés sortants dans les derniers mois avant de nouvelles élections. Attention à ne pas l'utiliser dans un essay formel : c'est une expression journalistique et légèrement moqueuse !
\end{savaistu}

\begin{aretenir}[L'essentiel de la section]
\begin{itemize}
\item Résultats : \eng{to win by a landslide / by a narrow margin}, \eng{to lose}, \eng{to concede defeat}, \eng{a concession speech}.
\item Contentieux : \eng{to file a petition}, \eng{to challenge the results}, \eng{to hear a case}, \eng{to deliver a ruling}, \eng{to uphold} (confirmer) ou \eng{to overturn / annul} (annuler).
\item Installation : \eng{to certify the results}, \eng{to be sworn in}, \eng{to take the oath of office}, \eng{to assume office}, \eng{transition of power}.
\item Légitimité : \eng{free and fair}, \eng{credible}, \eng{transparent}, \eng{peaceful}, \eng{disputed}, \eng{marred by violence}, \eng{rigged}.
\item Pièges : \eng{evidence} est incomptable ; on \eng{wins an election} mais on \eng{defeats an opponent} ; « annuler une élection » = \eng{annul}, pas \eng{cancel}.
\end{itemize}
\end{aretenir}

\coursec{Atelier d'Intégration : Production Orale \& Écrite (Exam Practice)}

Après avoir exploré, dans la section précédente, le lexique de l'après-vote — \eng{concede defeat}, \eng{file a petition}, \eng{swearing-in ceremony} — il est temps de \cle{réemployer} tout ce vocabulaire dans des situations de communication réelles, exactement comme au baccalauréat. L'objectif de cet atelier est double : d'une part, faire passer le lexique électoral de la \emph{reconnaissance} (je comprends le mot quand je le lis) à la \emph{production} (je l'utilise spontanément à l'oral et à l'écrit) ; d'autre part, s'entraîner sur des tâches calquées sur les épreuves officielles.

\begin{popfigure}
\begin{center}
\begin{tikzpicture}[node distance=0.55cm, every node/.style={font=\small},
box/.style={draw=popblue, fill=popblueL, rounded corners=2pt, align=center, inner sep=4pt, text width=3.4cm}]
\node[box] (a) {Input: Text / Audio};
\node[box, below=of a] (b) {Vocab Extraction};
\node[box, below=of b, align=center] (c){Structured Practice:\\ Gap-fill / Matching};
\node[box, below=of c, align=center] (d){Guided Production:\\ Dialogue / Paragraph};
\node[box, below=of d, align=center] (e){Free Production:\\ Debate / Article};
\node[box, below=of e, fill=poppurpleL, draw=poppurple, align=center] (f){Feedback:\\ Peer / Teacher Checklist};
\draw[-{Stealth}, thick, popdark] (a) -- (b);
\draw[-{Stealth}, thick, popdark] (b) -- (c);
\draw[-{Stealth}, thick, popdark] (c) -- (d);
\draw[-{Stealth}, thick, popdark] (d) -- (e);
\draw[-{Stealth}, thick, popdark] (e) -- (f);
\end{tikzpicture}
\end{center}
\legende{La spirale d'apprentissage du lexique : de la réception à la production libre, puis à la rétroaction.}
\end{popfigure}

\courssub{Tâche 1 — Interaction orale : Simulated Radio Debate}

\begin{experience}[Radio Debate : “Tackling Youth Unemployment”]
La classe simule un débat radiophonique en anglais, sur le modèle des émissions de \eng{CRTV Radio} ou de la \eng{BBC Focus on Africa}.

\textbf{Rôles (groupes de 4) :}
\begin{itemize}
\item \textbf{The Host} (l'animateur) : pose les questions, relance, impose le temps de parole.
\item \textbf{The Incumbent} (le candidat sortant) : défend son bilan (\eng{track record}).
\item \textbf{The Challenger} (le candidat de l'opposition) : attaque le bilan et présente son \eng{manifesto}.
\item \textbf{The Civil Society Representative} : pose des questions critiques au nom des électeurs (\eng{voters, constituents}).
\end{itemize}

\textbf{Thème :} \eng{Youth Unemployment in Our Constituency.}

\textbf{Contraintes linguistiques obligatoires :}
\begin{enumerate}
\item Employer au moins \textbf{5 idiomes ou collocations} de la leçon : \eng{to run for office}, \eng{to launch a campaign}, \eng{to deliver on one's promises}, \eng{to win over undecided voters}, \eng{a landslide victory}, \eng{to go to the polls}, \eng{empty promises}, \eng{to drum up support}.
\item Employer une \textbf{famille de mots complète} : \eng{elect} (v.), \eng{election} (n.), \eng{electoral} (adj.), \eng{electorate} (n.), \eng{elector} (n.).
\item Respecter un \textbf{registre formel} : pas de \eng{wanna, gonna} ; utiliser \eng{I would argue that…}, \eng{With all due respect…}.
\end{enumerate}

\textbf{Amorces utiles (sentence starters) :}
\begin{itemize}
\item \eng{Mr Chairman, my opponent has failed to deliver on his promises.} « Monsieur le Président, mon adversaire n'a pas tenu ses promesses. »
\item \eng{If elected, I pledge to create jobs for the youth of this constituency.} « Si je suis élu, je m'engage à créer des emplois pour les jeunes de cette circonscription. »
\item \eng{Let me remind the listeners that turnout among young voters was low in the last election.} « Permettez-moi de rappeler aux auditeurs que la participation des jeunes était faible au dernier scrutin. »
\end{itemize}
\end{experience}

\begin{methode}[Réussir un débat oral formel en anglais]
\begin{enumerate}
\item \textbf{Préparer} ses arguments en notes, jamais en texte lu : trois idées clés, chacune appuyée par une collocation électorale.
\item \textbf{Structurer} chaque prise de parole : opinion (\eng{I firmly believe that…}) → justification (\eng{because / since…}) → exemple (\eng{for instance…}).
\item \textbf{Réagir} à l'adversaire avec des formules polies : \eng{I take your point, however…}, \eng{With all due respect, the facts show otherwise.}
\item \textbf{Conclure} par une phrase mémorable : \eng{Ultimately, the choice is in the hands of the electorate.}
\end{enumerate}
\end{methode}

\courssub{Tâche 2 — Production écrite : Letter to the Editor}

\textbf{Sujet :} \eng{“Ensuring Free and Fair Elections in My Constituency”} — rédiger une lettre au rédacteur en chef d'un journal anglophone (environ 250 mots).

\textbf{Plan imposé :} Hook (accroche) → Context (\eng{turnout}, \eng{civic education}) → Problems (\eng{vote buying}, \eng{intimidation}) → Solutions (\eng{observers}, \eng{biometric voting}) → Call to action (appel à l'action).

\begin{textebox}[Writing Model : Letter to the Editor (Extrait)]
\textbf{Subject: Safeguarding the Integrity of the Upcoming Polls}

Sir, As a concerned \textbf{constituent}, I write to urge the \textbf{Electoral Commission} to ensure \textbf{transparent} voter registration. Reports of \textbf{ghost voters} on the \textbf{electoral roll} threaten the \textbf{credibility} of the process. We demand \textbf{biometric verification} at every \textbf{polling station} to prevent \textbf{multiple voting}. A \textbf{free and fair} election is the bedrock of our democracy. Yours faithfully, \textbf{A Voter}.
\end{textebox}

\begin{definition}[Ghost voter]
\eng{Ghost voter} (n.) : a fictitious name on the electoral roll, used to inflate figures or commit fraud. « Électeur fantôme » : nom fictif inscrit sur la liste électorale pour gonfler les chiffres ou frauder. \emph{Synonymes : dead voter, phantom voter.}
\end{definition}

\begin{propriete}[Connecteurs logiques pour l'argumentation (Writing)]
\begin{center}
\begin{tabularx}{\linewidth}{|l|X|}\hline
\rowcolor{popblueL} \textbf{Fonction} & \textbf{Connecteurs (registre formel)} \\ \hline
Ajouter & Furthermore, Moreover, In addition, Additionally \\ \hline
Contraster & However, Nevertheless, Conversely, On the other hand \\ \hline
Causalité & Consequently, Therefore, As a result, Hence \\ \hline
Exemplifier & For instance, Such as, Notably \\ \hline
Conclure & To conclude, Ultimately, In light of the above \\ \hline
\end{tabularx}
\end{center}
\emph{Exemple :} \eng{Vote buying undermines democracy; moreover, it discourages honest candidates from running for office.} « L'achat de voix sape la démocratie ; de plus, il décourage les candidats honnêtes de se présenter. »
\end{propriete}

\begin{attention}[Erreurs fréquentes au Bac]
\begin{itemize}
\item \eng{The population voted massively.} → \eng{Massively} suggère la force physique. Préférer : \eng{Turnout was high.} / \eng{Voters turned out in large numbers.} / \eng{There was a massive turnout.}
\item \eng{The candidates campaigned hardly.} → \eng{Hardly} signifie « à peine » ! Dire : \eng{campaigned hard} / \eng{campaigned vigorously}.
\item \eng{make campaign} → calque du français « faire campagne ». Dire : \eng{to campaign} ou \eng{to run / conduct a campaign}.
\end{itemize}
\end{attention}

\courssub{Tâche 3 — Compréhension et transfert : Reformulation journalistique}

Un communiqué administratif est sec, passif et nominalisé ; un article de presse est vivant, actif et concret. Savoir passer de l'un à l'autre prouve la maîtrise du lexique.

\begin{center}
\begin{tabularx}{\linewidth}{|X|X|}\hline
\rowcolor{popblueL} \textbf{Style administratif (communiqué)} & \textbf{Style journalistique (article)} \\ \hline
The election was conducted on Sunday. & Voters went to the polls on Sunday. \\ \hline
The counting of ballots was effected. & Officials counted the votes late into the night. \\ \hline
The incumbent candidate was defeated. & The incumbent lost his seat in a shock result. \\ \hline
A petition was filed by the opposition. & The opposition challenged the results in court. \\ \hline
Voter participation was satisfactory. & Turnout was high, with long queues at polling stations. \\ \hline
\end{tabularx}
\end{center}

\begin{exoresolu}[Du communiqué à l'article]
\textbf{Énoncé.} Réécrire en style journalistique : \eng{“The Electoral Commission announces that the election was conducted peacefully and that the incumbent was re-elected with a majority of the votes cast.”}

\tcblower

\textbf{Solution détaillée.} \eng{“Voters went to the polls peacefully on Sunday and handed the incumbent a fresh mandate, as he cruised to re-election with a clear majority of the votes cast.”}

\textbf{Commentaire :} on remplace le passif administratif (\eng{was conducted}) par un verbe d'action concret (\eng{went to the polls}) ; \eng{announces that} disparaît au profit du récit direct ; \eng{was re-elected} devient la collocation vivante \eng{cruised to re-election} (« a été réélu sans difficulté »). La phrase gagne en rythme et en images, sans perdre l'information.
\end{exoresolu}

\courssub{Révision Bac : sujets types}

\begin{itemize}
\item \textbf{Commentaire de texte (discours) :} un extrait de discours de campagne à analyser — repérer les promesses (\eng{pledge, vow, commit to}), les procédés de persuasion et le ton.
\item \textbf{Rédaction guidée (lettre au maire) :} \eng{Write a letter to your mayor complaining about the poor organisation of the last local elections.} — registre formel, plan problème/solutions.
\item \textbf{Traduction thématique (extrait de code électoral) :} « Tout citoyen âgé de vingt ans révolus et inscrit sur la liste électorale a le droit de vote. » → \eng{Every citizen aged twenty and above, who is registered on the electoral roll, has the right to vote.}
\end{itemize}

\begin{methode}[Checklist Rédaction Bac (Vocabulary Focus)]
\begin{enumerate}
\item Ai-je utilisé au moins 3 termes techniques précis (ex. \eng{ballot box, incumbent, manifesto}) ?
\item Ai-je évité les faux amis (\eng{scrutiny, bulletin, make campaign}) ?
\item Mes collocations sont-elles correctes (\eng{cast a vote, file a petition}) ?
\item L'orthographe des mots longs est-elle vérifiée (\eng{constituency, referendum, campaigning}) ?
\item Le registre est-il soutenu (pas de contractions \eng{don't, can't} dans l'écrit formel) ?
\end{enumerate}
\end{methode}

\begin{exercice}[Entraînement final]
\begin{enumerate}
\item Complétez : \eng{The opposition decided to \ldots\ a petition against the results.} (file / fill)
\item Traduisez : « Les électeurs se sont rendus aux urnes en grand nombre. »
\item Réécrivez en style journalistique : \eng{The candidate conceded defeat after the proclamation of the results.}
\end{enumerate}
\end{exercice}

\begin{corrige}[Exercice « Entraînement final »]
\begin{enumerate}
\item \eng{file} : \eng{The opposition decided to file a petition against the results.} (\eng{fill} = remplir, faux ami de « déposer ».)
\item \eng{Voters went to the polls in large numbers.} / \eng{Voters turned out in large numbers.}
\item \eng{The candidate graciously admitted he had lost, shortly after the official results were announced.} — on remplace la nominalisation administrative par un verbe actif et un adverbe journalistique.
\end{enumerate}
\end{corrige}

\courssub{Auto-évaluation : grille “Can-do”}

Cochez honnêtement chaque ligne ; toute case non cochée devient un objectif de révision avant l'épreuve.

\begin{center}
\begin{tabularx}{\linewidth}{|X|c|c|}\hline
\rowcolor{popblueL} \textbf{Je sais…} & \textbf{Oui} & \textbf{Pas encore} \\ \hline
définir \eng{gerrymandering} & & \\ \hline
utiliser \eng{concede defeat} dans une phrase & & \\ \hline
épeler \eng{ballot box} et \eng{constituency} sans erreur & & \\ \hline
prononcer \eng{constituency} (« con-sti-tu-èn-si », accent sur « sti ») & & \\ \hline
rédiger une lettre au rédacteur en chef avec 5 connecteurs formels & & \\ \hline
transformer un communiqué administratif en article de presse & & \\ \hline
\end{tabularx}
\end{center}

\begin{savaistu}[Gerrymandering]
\eng{Gerrymandering} (US/UK) désigne le découpage électoral truqué visant à favoriser un parti — le mot vient du gouverneur américain Elbridge \emph{Gerry}, dont une circonscription redécoupée avait la forme d'une salamandre (\eng{salamander}) ! Au Cameroun, le redécoupage des circonscriptions (\eng{redistricting / delimitation}) par ELECAM fait régulièrement débat. Un terme précieux pour briller dans une composition sur la citoyenneté.
\end{savaistu}

\newpage

\coursec{Activité d'intégration}

\courssub{Objectifs de l'activité}

À l'issue de cette simulation intégrée, l'élève sera capable de :
\begin{itemize}
\item \cle{réemployer} spontanément, à l'oral, le lexique complet du processus électoral étudié dans la leçon (\eng{polling station, presiding officer, ballot paper, ballot box, indelible ink, voter register, polling agent, observer}) ;
\item \cle{décrire} une procédure en temps réel à l'aide des verbes d'action du vote (\eng{cast a ballot, tick, fold, drop, verify, count}) et des connecteurs chronologiques (\eng{first, then, afterwards, finally}) ;
\item \cle{rédiger} un court rapport d'incident (\eng{incident report}) au passif, en utilisant le lexique du contentieux électoral (\eng{allege, irregularity, spoilt ballot, file a report, investigate}) ;
\item \cle{gérer} une situation de communication imprévue et stressante (réclamation, accusation, incident matériel) avec des expressions fonctionnelles appropriées ;
\item \cle{éviter} les faux amis et les calques du français (\eng{to vote for somebody}, et non « vote somebody » ; \eng{a polling station}, et non « a polling office »).
\end{itemize}

\courssub{Situation}

La classe se transforme en bureau de vote le jour d'une élection municipale fictive à Buea, dans la région du Sud-Ouest du Cameroun. Chaque groupe de six ou sept élèves constitue un \eng{polling station} (bureau de vote) et reçoit une fiche de rôle tirée au sort. La simulation se déroule en temps réel : ouverture à 7h00, incidents imprévus, fermeture et rédaction du rapport final.

\begin{textebox}[Election Day Live — Role Cards, Buea Municipal Election]
\textbf{Polling Station No. 12, Government High School Buea — Election Day, 7:00 a.m.}

\textbf{Presiding Officer} — You open the station. Read the opening procedure aloud: check the empty ballot box, seal it, display the voter register, and announce: “This polling station is now open. Voting will close at 6 p.m.”

\textbf{Polling Clerk} — You verify each voter's identity card against the voter register, hand out one ballot paper per voter, and mark each voter's finger with indelible ink.

\textbf{Polling Agents (two, from rival parties)} — You monitor the procedure. You may challenge any irregularity, but you must remain polite: “I would like to raise an objection…”

\textbf{Observer} — You watch silently, take notes on every incident, and write the final incident report.

\textbf{Voter 1 (first-time voter)} — You are nervous. Ask questions: “Where do I tick my choice? Is my vote secret?”

\textbf{Voter 2 (elderly voter)} — You cannot walk easily. Ask for assistance at the entrance.

\textbf{Unexpected events (drawn by the teacher at random):}
\begin{itemize}
\item A voter presents a damaged voter card — the photo is torn.
\item A polling agent accuses the clerk of bias: “You are helping the other party!”
\item The indelible ink has dried up and cannot be used.
\end{itemize}
\end{textebox}

\begin{definition}[Glossaire des rôles]
\begin{itemize}
\item \eng{presiding officer} (n.) : président de bureau de vote
\item \eng{polling clerk} (n.) : assesseur, secrétaire de bureau de vote
\item \eng{polling agent} (n.) : représentant d'un parti dans le bureau de vote
\item \eng{observer} (n.) : observateur (national ou international)
\item \eng{voter register} (n.) : liste électorale
\item \eng{indelible ink} (n.) : encre indélébile
\item \eng{bias} (n.) : partialité, parti pris
\end{itemize}
\end{definition}

\courssub{Consignes}

\begin{enumerate}
\item \textbf{Préparation des rôles (5 min).} Tirez les rôles au sort. Chaque élève relit sa fiche et note trois expressions du lexique de la leçon qu'il devra obligatoirement employer pendant la simulation. Le \eng{Presiding Officer} prépare son annonce d'ouverture avec les connecteurs chronologiques (\eng{first of all, then, next, finally}).

\item \textbf{Ouverture du bureau de vote (5 min).} Le \eng{Presiding Officer} lit la procédure d'ouverture à voix haute : vérification de l'urne vide (\eng{the ballot box is empty}), scellage (\eng{the box is sealed}), affichage de la liste électorale. Le reste du groupe écoute et relève le vocabulaire utilisé.

\item \textbf{Vote des électeurs (10 min).} Les deux \eng{voters} accomplissent toutes les étapes : vérification de l'identité, remise du bulletin, passage dans l'isoloir (\eng{polling booth}), vote, dépôt du bulletin plié dans l'urne, marquage du doigt à l'encre indélébile. Les \eng{polling agents} et l'\eng{observer} surveillent et formulent leurs remarques en anglais.

\item \textbf{Gestion de l'imprévu (5 min).} Le professeur distribue une carte « incident » à chaque groupe. Le groupe doit réagir en temps réel, en anglais : le \eng{Presiding Officer} tranche, les agents peuvent déposer une objection formelle (\eng{I wish to file a complaint}).

\item \textbf{Rédaction du rapport (10 min).} Chaque groupe rédige un \eng{Incident Report / Closing Statement} d'environ 150 mots, rédigé \textbf{au passif} chaque fois que possible, et contenant au moins cinq mots du lexique du contentieux : \eng{allege, irregularity, spoilt ballot, investigate, file a report, uphold, reject}.

\item \textbf{Présentation orale (5 min par groupe).} L'\eng{observer} de chaque groupe lit le rapport devant la classe. Les autres élèves complètent une grille d'écoute : nombre de mots du lexique correctement employés, emploi du passif, fluidité.
\end{enumerate}

\begin{methode}[Rédiger un rapport d'incident au passif]
\begin{enumerate}
\item Commencez par l'en-tête : lieu, date, numéro du bureau de vote (\eng{Polling Station No. 12, Buea — Election Day}).
\item Décrivez les faits au passif : \eng{The ballot box was sealed at 7 a.m.} — « L'urne a été scellée à 7 heures. »
\item Rapportez les accusations sans les affirmer : \eng{It was alleged that…} / \eng{The clerk was accused of bias.} — « Il a été allégué que… » / « L'assesseur a été accusé de partialité. »
\item Indiquez les mesures prises : \eng{The incident was recorded and a report was filed.} — « L'incident a été consigné et un rapport a été déposé. »
\item Concluez par une recommandation : \eng{We recommend that the matter be investigated.} — « Nous recommandons que l'affaire fasse l'objet d'une enquête. »
\end{enumerate}
\end{methode}

\courssub{Ressources à mobiliser}

\begin{aretenir}[Boîte à outils lexicale et grammaticale]
\begin{itemize}
\item \textbf{Acteurs} : \eng{presiding officer, polling clerk, polling agent, observer, voter, candidate, electorate}.
\item \textbf{Matériel et lieux} : \eng{ballot paper, ballot box, polling booth, polling station, voter register, indelible ink}.
\item \textbf{Verbes d'action} : \eng{to cast a ballot, to tick a box, to fold the ballot paper, to drop it into the box, to verify identity, to count the votes, to announce the results}.
\item \textbf{Contentieux} : \eng{irregularity, spoilt ballot, damaged card, bias, intimidation, to allege, to file a complaint, to investigate, to uphold / to reject a complaint}.
\item \textbf{Grammaire} : le passif (\eng{be + participe passé}) pour décrire une procédure sans désigner de coupable ; le discours rapporté pour rapporter les objections (\eng{The agent claimed that the clerk was biased.}).
\item \textbf{Connecteurs} : \eng{first of all, then, next, afterwards, finally, however, consequently}.
\end{itemize}
\end{aretenir}

\begin{attention}[Pièges à éviter pendant la simulation]
\begin{itemize}
\item On dit \eng{to vote \textbf{for} a candidate} — « voter pour un candidat » ; jamais « to vote a candidate ».
\item \eng{A ballot} est le bulletin de vote ; ne confondez pas avec \eng{a ballet} (un ballet de danse) — attention à la prononciation : « ba-lot », le \emph{t} final ne se prononce pas.
\item \eng{To assist} signifie « aider » (\eng{The elderly voter was assisted.}) et non « assister à » ; pour « assister à », utilisez \eng{to attend} ou \eng{to observe}.
\item Un \eng{election result} se prononce « ri-zolte » ; ne dites pas « résultat » à la française.
\item Évitez le calque « the ink is finished » : dites \eng{the ink has dried up} (l'encre a séché) ou \eng{the ink has run out} (l'encre est épuisée).
\end{itemize}
\end{attention}

\courssub{Proposition de production et corrigé commenté}

\begin{exoresolu}[Modèle de Incident Report / Closing Statement]
Rédigez le rapport final de l'\eng{observer} du bureau de vote n° 12 de Buea (environ 150 mots), en utilisant le passif et au moins cinq termes du lexique du contentieux.
\tcblower
\textbf{Incident Report — Polling Station No. 12, Government High School Buea}

Polling Station No. 12 was opened at 7:00 a.m. by the Presiding Officer. The ballot box was checked, sealed and displayed in front of the polling agents and the observer. The voter register was verified before the first ballot was cast.

Three incidents were recorded during the day. First, a voter presented a damaged voter card; his identity was confirmed against the register and he was allowed to vote. Secondly, the clerk was accused of bias by a polling agent. However, no evidence was provided, and the complaint was rejected by the Presiding Officer. Thirdly, the indelible ink dried up at midday; it was immediately replaced, and no voter was registered twice.

It is alleged that one ballot paper was spoilt, but this claim could not be confirmed. Overall, the election was conducted peacefully and no intimidation was observed. We recommend that the rejected complaint be investigated by the electoral commission.

\textit{(152 mots)}

\textbf{Commentaire des choix de langue :}
\begin{itemize}
\item \textbf{Passif systématique} : \eng{was opened, was checked, sealed and displayed, was verified, was cast, were recorded, was confirmed, was accused, was provided, was rejected, was replaced, was registered, is alleged, was spoilt, was conducted, was observed, be investigated}. Le passif convient au rapport officiel : il décrit les faits sans accuser personne directement.
\item \textbf{Lexique du contentieux} : \eng{accused of bias, complaint was rejected, it is alleged, spoilt ballot, investigated} — cinq termes exigés, tous correctement employés.
\item \textbf{Prudence du rapporteur} : \eng{It is alleged that…} et \eng{this claim could not be confirmed} montrent que l'observateur rapporte sans affirmer — compétence clé du discours rapporté.
\item \textbf{Connecteurs} : \eng{first, secondly, thirdly, however, overall} structurent le récit chronologique.
\item \textbf{Faux amis évités} : \eng{was allowed to vote} (et non « was permitted to vote » dans ce registre), \eng{was conducted} pour « s'est déroulée » (et non « was realised », faux ami de « réaliser »).
\end{itemize}
\end{exoresolu}

\courssub{Bilan de l'activité}

Cette simulation a permis de mobiliser \textbf{l'ensemble du lexique de la leçon} dans une situation de communication authentique et contraignante : le temps réel interdit la traduction mot à mot et force l'emploi direct de l'anglais. Trois compétences ont été évaluées simultanément : la \cle{justesse lexicale} (le bon mot pour le bon rôle et la bonne procédure), la \cle{maîtrise grammaticale} (passif et discours rapporté dans le rapport écrit) et l'\cle{aisance orale} (réactions aux imprévus, objections, annonces). Pour vérifier votre progression, demandez-vous : suis-je capable, sans notes, d'expliquer en anglais le parcours complet d'un électeur, de l'ouverture du bureau de vote à la proclamation des résultats ? Suis-je capable de rapporter une irrégularité au passif sans calquer le français ? Si oui, l'objectif de la leçon est atteint ; sinon, reprenez la boîte « Ressources à mobiliser » et rejouez la scène avec un autre rôle.

\newpage

\coursec{Méthodes et exercices résolus}

\coursec{Lesson 47 — Vocabulary: Words and expressions related to campaigning and voting for leaders}

\courssub{Savoir-faire 1 : Identifier et employer le vocabulaire du thème (Processus électoral : chronologie et acteurs)}

\begin{definition}[Les trois temps du processus électoral]
Le vocabulaire électoral s'organise naturellement autour de trois phases chronologiques. Les maîtriser permet de situer immédiatement le contexte d'un texte ou d'un dialogue.
\begin{center}
\begin{tabularx}{\linewidth}{|X|X|X|}
\hline
\rowcolor{popblueL} \textbf{Phase} & \textbf{Mots-clés (Noms / Verbes / Adjectifs)} & \textbf{Acteurs principaux} \\
\hline
\cle{Avant le vote} (Pre-election) & \eng{campaign, manifesto, rally, candidate, to run for office, to canvass, polling day eve} & \eng{candidate, campaigner, party leader, voter, elector} \\
\hline
\cle{Le jour du vote} (Polling day) & \eng{polling station, ballot paper, ballot box, booth, indelible ink, voter register, to cast a vote, to queue} & \eng{presiding officer, polling clerk, observer, voter, agent} \\
\hline
\cle{Après le vote} (Post-election) & \eng{counting, results, turnout, landslide, concession speech, petition, inauguration, swearing-in} & \eng{returning officer, winner, MP / Councillor, president-elect, coalition} \\
\hline
\end{tabularx}
\end{center}
\end{definition}

\begin{textebox}[Script d'écoute simulé — Journal de 13 h, Cameroon Radio Television]
“Good afternoon. The Electoral Board has announced that the presidential campaign officially opens tomorrow. Three candidates have been cleared to run for office. They will criss-cross the country holding rallies and presenting their manifestos. On polling day, over seven million registered voters are expected to cast their ballots at polling stations nationwide. Observers from the African Union and the Commonwealth will monitor the process. Results are expected within forty-eight hours after the counting.”
\end{textebox}

\begin{methode}[Identifier et employer le vocabulaire électoral en contexte]
Pour reconnaître et utiliser correctement le lexique de la campagne et du vote dans un texte ou un dialogue, suis cette démarche :
\begin{enumerate}
\item \textbf{Repère le moment du processus électoral} décrit dans le texte : avant le vote (\cle{campaign}, \cle{manifesto}, \cle{rally}), le jour du vote (\cle{polling station}, \cle{ballot paper}, \cle{ballot box}), après le vote (\cle{counting}, \cle{results}, \cle{swearing-in}). Le vocabulaire se range naturellement dans ces trois familles.
\item \textbf{Identifie les acteurs} : \cle{candidate} (candidat), \cle{voter} / \cle{elector} (électeur), \cle{presiding officer} (président du bureau de vote), \cle{observer} (observateur), \cle{returning officer} (responsable de la proclamation des résultats).
\item \textbf{Classe chaque mot par nature grammaticale} : nom (\emph{a vote}, \emph{a ballot}), verbe (\emph{to vote}, \emph{to cast a ballot}), adjectif (\emph{electoral}, \emph{eligible}). Cela t'évite de confondre \eng{to elect} (verbe) et \eng{an election} (nom).
\item \textbf{Vérifie la collocation exacte} avant d'écrire : on dit \eng{to cast a vote} (et non « make a vote »), \eng{to run for office} (se porter candidat), \eng{to hold an election} (organiser une élection).
\item \textbf{Réemploie le mot dans une phrase personnelle} liée au contexte camerounais ou africain : \eng{Voters in Yaoundé queued patiently outside the polling station.}
\end{enumerate}
\end{methode}

\begin{exoresolu}[Lexique en contexte — complétion de texte]
\textbf{Énoncé.} Complete the following passage with the correct words from the box. Use each word once.

\begin{center}
ballot box — campaign — candidates — cast — polling station — results
\end{center}

“On election day in Bamenda, voters went to the nearest (1) \ldots\ early in the morning. Each voter (2) \ldots\ his or her vote by putting the ballot paper into the (3) \ldots\ . Three (4) \ldots\ had taken part in a long (5) \ldots\ , travelling from village to village. In the evening, everyone waited for the official (6) \ldots\ on the radio.”
\tcblower
\textbf{Solution détaillée.}
\begin{enumerate}
\item \textbf{polling station} : le lieu où l'on vote est le \cle{polling station} (bureau de vote). “Went to the nearest \ldots” exige un lieu ; \emph{ballot box} est impossible car on ne va pas “à” une urne.
\item \textbf{cast} : la collocation correcte est \eng{to cast a vote} (exprimer son suffrage). Le texte est au prétérit (\emph{went}, \emph{had taken}) ; le prétérit de \emph{cast} est \emph{cast} (verbe irrégulier invariable : cast — cast — cast).
\item \textbf{ballot box} : on glisse le bulletin (\emph{ballot paper}) dans l'urne, \cle{ballot box}. Collocation : \eng{to put the ballot paper into the ballot box}.
\item \textbf{candidates} : “Three \ldots” exige un nom pluriel de personnes ; les \cle{candidates} sont ceux qui se présentent à l'élection.
\item \textbf{campaign} : la collocation est \eng{to take part in a campaign} (participer à une campagne électorale). “Travelling from village to village” confirme l'idée de campagne.
\item \textbf{results} : après le vote, on attend les \cle{results} (résultats). Attention au pluriel : on dit \emph{the results are}, jamais “the result are”.
\end{enumerate}
\textbf{Texte complété :} “On election day in Bamenda, voters went to the nearest \textbf{polling station} early in the morning. Each voter \textbf{cast} his or her vote by putting the ballot paper into the \textbf{ballot box}. Three \textbf{candidates} had taken part in a long \textbf{campaign}, travelling from village to village. In the evening, everyone waited for the official \textbf{results} on the radio.”
\textbf{Conclusion :} le choix d'un mot se justifie toujours par trois critères : le sens, la nature grammaticale exigée par la phrase, et la collocation correcte.
\end{exoresolu}

\courssub{Savoir-faire 2 : Maîtriser les collocations et expressions idiomatiques (Campagne : discours, promesses, médias)}

\begin{propriete}[Collocations verbales indispensables — Tableau de révision]
En anglais, le verbe qui accompagne un nom est souvent figé (collocation). Apprends ces blocs comme des unités indivisibles.
\begin{center}
\begin{tabularx}{\linewidth}{|l|X|X|}
\hline
\rowcolor{popblueL} \textbf{Nom clé} & \textbf{Verbe(s) correct(s)} & \textbf{Exemple traduit} \\
\hline
\cle{election} & \eng{to hold}, \eng{to organise} & \eng{The country holds elections every five years.} (Le pays organise des élections tous les cinq ans.) \\
\cle{campaign} & \eng{to launch}, \eng{to conduct}, \eng{to run}, \eng{to wage} & \eng{She launched her campaign in Douala.} (Elle a lancé sa campagne à Douala.) \\
\cle{candidate / office} & \eng{to run for}, \eng{to stand for} (UK), \eng{to contest} & \eng{He decided to run for president.} (Il a décidé de se présenter à la présidence.) \\
\cle{speech} & \eng{to make}, \eng{to deliver}, \eng{to give} & \eng{The leader made a powerful speech.} (Le leader a prononcé un discours puissant.) \\
\cle{promise} & \eng{to make}, \eng{to keep}, \eng{to break}, \eng{to fulfil} & \eng{Politicians often break their promises.} (Les politiciens trahissent souvent leurs promesses.) \\
\cle{vote / ballot} & \eng{to cast} & \eng{Every citizen has the right to cast a vote.} (Tout citoyen a le droit d'exprimer son vote.) \\
\cle{referendum} & \eng{to hold}, \eng{to call} & \eng{The government called a referendum.} (Le gouvernement a convoqué un référendum.) \\
\hline
\end{tabularx}
\end{center}
\end{propriete}

\begin{definition}[Idiomes et expressions figées du langage politique]
Ces expressions sont fréquentes dans la presse et les commentaires électoraux. Leur sens n'est pas littéral.
\begin{itemize}
\item \cle{a landslide victory} : victoire écrasante (litt. “victoire par glissement de terrain”).
\item \cle{a close race} / \cle{a tight race} : course serrée, scrutin indécis.
\item \cle{to win by a wide margin} : gagner avec une large avance.
\item \cle{a hung parliament} : parlement sans majorité absolue (Royaume-Uni, Cameroun : Assemblée sans majorité claire).
\item \cle{to concede defeat} : reconnaître sa défaite (acte démocratique fort).
\item \cle{a swing vote} / \cle{swing voter} : vote / électeur indécis qui fait basculer l'élection.
\item \cle{to throw one's hat in the ring} : annoncer sa candidature (litt. “jeter son chapeau dans l'arène”).
\item \cle{grassroots campaign} : campagne de terrain, porte-à-porte, mobilisation de la base.
\end{itemize}
\end{definition}

\begin{methode}[Apprendre les collocations et les idiomes électoraux]
\begin{enumerate}
\item \textbf{Apprends les verbes qui accompagnent les noms clés} (voir tableau ci-dessus).
\item \textbf{Mémorise les idiomes fréquents} en les associant à une image mentale (\emph{landslide} = glissement de terrain = tout emporte).
\item \textbf{Ne traduis jamais mot à mot} : “faire campagne” se dit \eng{to campaign} ou \eng{to run a campaign}, jamais “to make campaign”.
\item \textbf{Entraîne-toi à reformuler} : remplace un idiome par son équivalent simple, et inversement. \eng{She won by a wide margin} = \eng{She won with many more votes than her opponent.}
\item \textbf{Crée des phrases personnelles} avec chaque idiome pour l'ancrer en mémoire.
\end{enumerate}
\end{methode}

\begin{exoresolu}[Collocations et idiomes — choix et reformulation]
\textbf{Énoncé.}
\textbf{A.} Choose the correct verb in each sentence.
\begin{enumerate}
\item The president (made / held / did) a speech in front of thousands of supporters.
\item Our country (holds / makes / takes) municipal elections every five years.
\item The candidate (did / made / gave) many promises during the campaign, but she broke them all.
\end{enumerate}
\textbf{B.} Rewrite each sentence using the idiom in brackets.
\begin{enumerate}
\item Mrs Etame won the election with 78 \% of the votes. (a landslide victory)
\item The two candidates obtained almost the same number of votes. (a close race)
\end{enumerate}
\tcblower
\textbf{Solution détaillée.}
\textbf{A.}
\begin{enumerate}
\item \textbf{made} : la collocation correcte est \eng{to make a speech} (prononcer un discours). “Did a speech” est un calque du français “faire un discours” ; “held a speech” n'existe pas en anglais standard. \eng{The president made a speech in front of thousands of supporters.}
\item \textbf{holds} : on dit \eng{to hold elections} (organiser des élections). Le sujet est “our country” (3\textsuperscript{e} personne du singulier), d'où \emph{holds} avec un \emph{-s}. \eng{Our country holds municipal elections every five years.}
\item \textbf{made} : la collocation est \eng{to make a promise} (faire une promesse), et son contraire \eng{to break a promise}. “Did promises” est un calque du français. \eng{The candidate made many promises during the campaign, but she broke them all.}
\end{enumerate}
\textbf{B.}
\begin{enumerate}
\item \eng{Mrs Etame won the election in a landslide victory, with 78 \% of the votes.} L'idiome \cle{a landslide victory} exprime une victoire écrasante ; 78 \% des voix justifie pleinement son emploi.
\item \eng{The election was a close race: the two candidates obtained almost the same number of votes.} \cle{A close race} (une course serrée) décrit un scrutin dont le résultat est très serré.
\end{enumerate}
\textbf{Conclusion :} en anglais, le verbe qui accompagne un nom est souvent différent du français : on \emph{makes} a speech et a promise, on \emph{holds} an election. Ces collocations doivent être apprises comme des blocs.
\end{exoresolu}

\courssub{Savoir-faire 3 : Exploiter les familles de mots, la prononciation et l'orthographe}

\begin{propriete}[Familles de mots électoraux — Dérivation et suffixes]
La dérivation (word formation) est cruciale pour le \emph{Use of English} et la rédaction. Repère les suffixes productifs.
\begin{center}
\begin{tabular}{|l|l|l|l|l|}
\hline
\rowcolor{popblueL} \textbf{Verbe / Base} & \textbf{Nom Action / Concept (-tion, -ment)} & \textbf{Nom Personne (-or, -er, -ist)} & \textbf{Adjectif (-al, -ive, -ory)} & \textbf{Nom Qualité / Fonction (-ship, -cy)} \\
\hline
\eng{to elect} & \eng{election} & \eng{elector} & \eng{electoral} & \eng{electorate} (ensemble des électeurs) \\
\eng{to vote} & \eng{vote} & \eng{voter} & \eng{voting} (adj.) & — \\
\eng{to campaign} & \eng{campaign} & \eng{campaigner} & \eng{campaign} (adj.) & — \\
\eng{to lead} & \eng{leadership} & \eng{leader} & \eng{leading} & \eng{leadership} \\
\eng{to govern} & \eng{government} & \eng{governor} & \eng{governmental} & \eng{governance} \\
\eng{to nominate} & \eng{nomination} & \eng{nominee} & \eng{nominative} & — \\
\eng{to preside} & \eng{presidency} & \eng{president} & \eng{presidential} & — \\
\eng{to legislate} & \eng{legislation} & \eng{legislator} & \eng{legislative} & — \\
\hline
\end{tabular}
\end{center}
\end{propriete}

\begin{aretenir}[Orthographe et Pièges fréquents]
\begin{itemize}
\item \cle{Government} : s'écrit avec un \textbf{n} avant le \textbf{m} (g-o-v-e-r-\textbf{n}-m-e-n-t). Ne pas écrire “goverment”.
\item \cle{Candidate} : se termine par \textbf{-date} (comme \emph{date}), pas “-dat” ni “-dant”.
\item \cle{Ballot} : double \textbf{l} (b-a-\textbf{ll}-o-t).
\item \cle{Referendum} : pluriel \textbf{referendums} (anglais) ou \textbf{referenda} (latin), pas “referendums” avec un seul \emph{m} ? Si, \emph{referendums} prend un \emph{s}. Attention : \emph{referendum} (singulier) un seul \emph{m} ? Non, \textbf{referendum} (un seul \emph{m} au milieu, un seul \emph{m} final). \textbf{Referendum}.
\item \cle{Manifesto} : pluriel \textbf{manifestos}.
\end{itemize}
\end{aretenir}

\begin{exemplebox}[Prononciation — Accent tonique (Word Stress)]
En anglais, l'accent tonique change souvent le sens ou la nature du mot. Marque l'accent sur la syllabe en \textbf{gras} et entraîne-toi à le prononcer fort et long.
\begin{center}
\begin{tabular}{|l|l|l|}
\hline
\rowcolor{popblueL} \textbf{Mot} & \textbf{Phonétique approximative (français)} & \textbf{Règle / Note} \\
\hline
\eng{candidate} & \textbf{kan}-di-dèt & Accent sur 1\textsuperscript{re} syllabe. \\
\eng{campaign} & kam-\textbf{pègne} & \emph{g} muet, accent final. \\
\eng{election} & i-\textbf{lèk}-cheun & Accent sur 2\textsuperscript{e} syllabe. \\
\eng{electoral} & i-\textbf{lèk}-to-ral & Accent sur 2\textsuperscript{e} (même radical). \\
\eng{government} & \textbf{go}-veurn-ment & 1\textsuperscript{re} syllabe, \emph{n} souvent muet à l'oral rapide. \\
\eng{president} & \textbf{pré}-zi-dent & Accent sur 1\textsuperscript{re}. \\
\eng{presidential} & pré-zi-\textbf{den}-chl & Accent sur 3\textsuperscript{e} (règle \emph{-ial}). \\
\eng{referendum} & ré-fé-\textbf{ren}-dèm & Accent sur 3\textsuperscript{e}. \\
\hline
\end{tabular}
\end{center}
\end{exemplebox}

\begin{exoresolu}[Familles de mots — formation et emploi]
\textbf{Énoncé.} Complete each sentence with the correct form of the word in capitals.
\begin{enumerate}
\item Every \ldots has the right to choose his or her leaders. (ELECT)
\item The \ldots of the new mayor will take place next Saturday. (ELECT)
\item Good \ldots requires honesty and courage. (LEAD)
\item The \ldots has promised to build new classrooms. (GOVERN)
\item Thousands of \ldots marched through Douala to support their candidate. (CAMPAIGN)
\end{enumerate}
\tcblower
\textbf{Solution détaillée.}
\begin{enumerate}
\item \textbf{elector} : “Every \ldots has” exige un nom singulier de personne ; le suffixe \emph{-or} forme le nom d'agent : \emph{elect} $\rightarrow$ \emph{elector} (électeur). \eng{Every elector has the right to choose his or her leaders.}
\item \textbf{election} : “The \ldots of the new mayor” exige un nom d'action ; suffixe \emph{-tion} : \emph{elect} $\rightarrow$ \emph{election}. \eng{The election of the new mayor will take place next Saturday.}
\item \textbf{leadership} : l'adjectif “Good” appelle un nom abstrait ; le suffixe \emph{-ship} forme le nom de la qualité ou de la fonction : \emph{lead} $\rightarrow$ \emph{leader} $\rightarrow$ \emph{leadership} (le leadership, l'art de diriger). \eng{Good leadership requires honesty and courage.}
\item \textbf{government} : “The \ldots has promised” exige un nom singulier désignant l'institution : \emph{govern} $\rightarrow$ \emph{government}. Attention à l'orthographe : g-o-v-e-r-\textbf{n}-m-e-n-t. \eng{The government has promised to build new classrooms.}
\item \textbf{campaigners} : “Thousands of \ldots marched” exige un nom pluriel de personnes ; suffixe \emph{-er} : \emph{campaign} $\rightarrow$ \emph{campaigner}, pluriel \emph{campaigners}. \eng{Thousands of campaigners marched through Douala to support their candidate.}
\end{enumerate}
\textbf{Conclusion :} pour trouver la bonne forme, demande-toi d'abord quelle nature grammaticale la phrase exige (personne ? action ? adjectif ?), puis applique le suffixe correspondant.
\end{exoresolu}

\courssub{Savoir-faire 4 : Éviter les faux amis et les calques du français (Jour du scrutin : procédures, matériel, verbes d'action)}

\begin{propriete}[Faux amis et pièges de traduction — Tableau comparatif]
\begin{center}
\begin{tabularx}{\linewidth}{|l|X|X|X|}
\hline
\rowcolor{popblueL} \textbf{Mot anglais (Piège)} & \textbf{Sens RÉEL en anglais} & \textbf{Mot français visé} & \textbf{Traduction correcte en anglais} \\
\hline
\eng{to assist} & Aider, secourir & Assister à (être présent) & \eng{to attend} \\
\eng{to demand} & Exiger (avec autorité) & Demander (poliment) & \eng{to ask (for)}, \eng{to request} \\
\eng{a meeting} & Réunion (général) & Meeting politique, rassemblement & \eng{a rally}, \eng{a public meeting} \\
\eng{a deputy} & Adjoint, suppléant, lieutenant & Député (parlementaire) & \eng{a Member of Parliament (MP)}, \eng{a lawmaker} \\
\eng{a poll} & Scrutin, vote ; Sondage & Sondage d'opinion & \eng{an opinion poll}, \eng{a survey} \\
\eng{to vote} (+ objet) & Voter \textbf{pour} qqn / qqch & Voter une loi, voter qqn & \eng{to pass a law}, \eng{to vote for a candidate} \\
\eng{sensible} & Raisonnable, sensé & Sensible (émotif) & \eng{sensitive} \\
\eng{actually} & En fait, en réalité & Actuellement & \eng{currently}, \eng{at present} \\
\hline
\end{tabularx}
\end{center}
\end{propriete}

\begin{methode}[Repérer et corriger les faux amis et calques électoraux]
\begin{enumerate}
\item \textbf{Méfie-toi des mots qui ressemblent au français} : \eng{to demand} = exiger ; \eng{sensible} = raisonnable.
\item \textbf{Attention aux faux amis du domaine politique} (voir tableau ci-dessus) : \eng{deputy} $\neq$ député ; \eng{meeting} $\neq$ meeting politique (\eng{rally}) ; \eng{assist} $\neq$ assister à (\eng{attend}).
\item \textbf{Évite les calques de structure} :
\begin{itemize}
\item “Voter pour quelqu'un” $\rightarrow$ \eng{to vote \textbf{for} somebody} (préposition \textbf{for} obligatoire).
\item “Élire quelqu'un président” $\rightarrow$ \eng{to elect somebody president} (sans “as” ni “like”).
\item “Se porter candidat” $\rightarrow$ \eng{to run for office} / \eng{to stand as a candidate} (pas “to present oneself”).
\end{itemize}
\item \textbf{Vérifie toujours la préposition} : \eng{to vote for a candidate}, \eng{to vote on an issue} (voter sur une question), \eng{to campaign for reforms}, \eng{to be interested in politics}.
\item \textbf{Relis ta production} en te demandant : “Cette structure existe-t-elle vraiment en anglais, ou est-ce du français déguisé ?”
\end{enumerate}
\end{methode}

\begin{exoresolu}[Faux amis et calques — correction de phrases]
\textbf{Énoncé.} Each sentence contains a mistake caused by French interference. Correct it.
\begin{enumerate}
\item Many citizens assisted the candidate's meeting in Buea.
\item The people demanded an explanation, so they asked to see the mayor. (Correct the verb “demanded” if the meaning is simply “asked”).
\item I will vote him because he is honest.
\item The deputies of the opposition refused to vote the law.
\end{enumerate}
\tcblower
\textbf{Solution détaillée.}
\begin{enumerate}
\item \eng{Many citizens attended the candidate's rally in Buea.} Double correction : \eng{to assist} signifie “aider” ; “assister à” se dit \eng{to attend}. De plus, un “meeting politique” français correspond à \eng{a rally} (rassemblement électoral).
\item \eng{The people asked for an explanation and requested to see the mayor.} \eng{To demand} = exiger avec force ; pour “demander” simplement, on utilise \eng{to ask (for)} ou \eng{to request}.
\item \eng{I will vote \textbf{for} him because he is honest.} Le verbe \emph{vote} est ici intransitif : il exige la préposition \emph{for} devant la personne soutenue. “Vote him” est un calque direct.
\item \eng{The opposition MPs refused to pass the law.} (ou : \emph{refused to vote for the law}). Deux corrections : le “député” français se dit \eng{Member of Parliament (MP)}, car \eng{deputy} signifie “adjoint” ; et “voter une loi” se dit \eng{to pass a law} ou \eng{to vote for a law}, jamais “to vote the law”.
\end{enumerate}
\textbf{Conclusion :} chaque fois qu'un mot anglais ressemble à un mot français, vérifie son sens réel dans le contexte ; la ressemblance est souvent un piège.
\end{exoresolu}

\courssub{Savoir-faire 5 : Le lexique post-électoral — Résultats, contentieux et installation}

\begin{definition}[Vocabulaire clé : Après le vote (Post-Election)]
\begin{itemize}
\item \cle{The count} / \cle{counting} : le dépouillement.
\item \cle{Turnout} : taux de participation (\eng{a high/low turnout}). \eng{Voter turnout was 65\%.}
\item \cle{Results} (pluriel) : résultats. \eng{The results were announced yesterday.}
\item \cle{Landslide (victory)} : victoire écrasante.
\item \cle{Concession speech} : discours de concession (le perdant félicite le gagnant).
\item \cle{To concede defeat} : reconnaître sa défaite.
\item \cle{A petition} / \cle{to file a petition} : un recours / déposer un recours (contentieux électoral).
\item \cle{The Supreme Court} / \cle{The Constitutional Council} : instance de validation (Cameroun : Conseil Constitutionnel).
\item \cle{Inauguration} / \cle{Swearing-in ceremony} : investiture / prestation de serment.
\item \cle{President-elect} / \cle{Mayor-elect} : président élu / maire élu (avant l'investiture).
\item \cle{Coalition} / \cle{Coalition government} : coalition / gouvernement de coalition.
\end{itemize}
\end{definition}

\begin{textebox}[Article de presse simulé — The Cameroon Tribune, Post-Election]
“Following a high turnout of 72\% in Sunday's presidential election, the Electoral Commission began the counting process immediately. Provisional results give the incumbent a landslide victory with 68\% of the vote. The main opposition candidate has refused to concede defeat, alleging irregularities at several polling stations. His party announced it would file a petition at the Constitutional Council within the legal timeframe. The inauguration of the president-elect is scheduled for next month, provided the Court validates the results.”
\end{textebox}

\begin{exoresolu}[Lexique post-électoral — Complétion et reformulation]
\textbf{Énoncé.} Complete the sentences with the correct word/phrase from the box.
\begin{center}
concede defeat — turnout — petition — president-elect — swearing-in — landslide
\end{center}
\begin{enumerate}
\item Despite losing by a \ldots, the opposition leader refused to \ldots immediately.
\item The \ldots was surprisingly low in urban areas, barely reaching 40\%.
\item After the Court validated the results, the \ldots prepared for his \ldots ceremony.
\item The coalition filed a \ldots challenging the results in three regions.
\end{enumerate}
\tcblower
\textbf{Solution.}
\begin{enumerate}
\item \textbf{landslide victory} (ou \emph{landslide}) / \textbf{concede defeat}. (Collocation : \emph{landslide victory}, \emph{concede defeat}).
\item \textbf{turnout}. (Sujet singulier : \emph{was}).
\item \textbf{president-elect} / \textbf{swearing-in}. (Titre avant investiture ; nom de la cérémonie).
\item \textbf{petition}. (Nom comptable : \emph{a petition} / \emph{filed a petition}).
\end{enumerate}
\end{exoresolu}

\courssub{Savoir-faire 6 : Réemployer le lexique à l'oral et à l'écrit (Atelier d'intégration)}

\begin{methode}[Produire un court texte ou un discours électoral — Guide pas à pas]
\begin{enumerate}
\item \textbf{Prépare un plan en trois parties} : introduction (la situation électorale), développement (les actions : campagne, vote, résultats), conclusion (ton opinion ou tes espoirs).
\item \textbf{Mobilise au moins huit mots du champ lexical} : \emph{candidate, campaign, manifesto, rally, polling station, ballot paper, turnout, results, concede, inauguration}.
\item \textbf{Utilise des connecteurs} pour structurer : \emph{first of all}, \emph{then}, \emph{however}, \emph{as a result}, \emph{in conclusion}, \emph{furthermore}.
\item \textbf{Varie les temps} à bon escient : prétérit pour raconter le scrutin passé (\emph{Voters queued, candidates campaigned}), futur pour les promesses (\emph{The new mayor will build roads}), présent pour les vérités générales (\emph{Voting is a civic duty}).
\item \textbf{À l'oral}, soigne la prononciation des mots clés : \emph{candidate} “kan-di-dète”, \emph{government} “go-veurn-ment”, \emph{campaign} “kam-pègne”. Parle lentement et articule.
\item \textbf{Relis-toi} : accords sujet-verbe, prépositions après \emph{vote} et \emph{campaign}, orthographe de \emph{government} et \emph{ballot}.
\end{enumerate}
\end{methode}

\begin{experience}[Activité orale — Simulation : Débat télévisé / Interview]
\textbf{Contexte :} Tu es journaliste pour “Cameroon Today”. Interviewe un candidat (ton camarade) ou un électeur après le vote.
\textbf{Rôles :}
\begin{itemize}
\item \textbf{Journaliste} : Utilise \emph{manifesto, rally, promises, turnout, expectations, polling station}.
\item \textbf{Candidat} : Utilise \emph{campaign, grassroots, landslide, serve, constituency, leadership}.
\item \textbf{Électeur} : Utilise \emph{queue, ballot paper, cast, civic duty, change, transparency}.
\end{itemize}
\textbf{Contrainte :} Minimum 5 tours de parole chacun. Note les erreurs de collocation ou de faux amis pour correction collective.
\end{experience}

\begin{exoresolu}[Production écrite type examen — Article d'opinion (150 mots)]
\textbf{Énoncé.} Write an article for your school magazine entitled: “Why voting matters in a democracy”. Use at least ten words/expressions from the lesson. Structure: Introduction $\rightarrow$ Arguments (2-3) $\rightarrow$ Conclusion.
\tcblower
\textbf{Solution détaillée — Démarche :}
\begin{enumerate}
\item \textbf{Plan} : (1) Introduction : définition du vote comme pilier. (2) Argument 1 : Choix des leaders / reddition de comptes (\emph{accountability}). (3) Argument 2 : Paix sociale / légitimité (\emph{legitimacy, peaceful transition}). (4) Conclusion : Appel à l'inscription sur les listes (\emph{voter register}).
\item \textbf{Lexique visé} : \emph{civic duty, electorate, ballot, polling station, candidate, manifesto, accountability, legitimacy, turnout, inauguration}.
\item \textbf{Grammaire} : Présent (vérités), Modaux (\emph{must, should, can}), Passif (\emph{leaders are chosen}).
\end{enumerate}
\textbf{Réponse rédigée :}

“Voting is not merely a right; it is a fundamental \textbf{civic duty} for every citizen. In a democracy, the \textbf{electorate} holds the ultimate power: the power to choose its \textbf{leaders} through the \textbf{ballot}. When voters go to the \textbf{polling station}, study the \textbf{manifestos} and \textbf{cast their vote} for a \textbf{candidate}, they demand \textbf{accountability}. A high \textbf{turnout} grants \textbf{legitimacy} to the winner and ensures a peaceful \textbf{transition} of power. Conversely, abstention weakens the mandate of the \textbf{president-elect} and endangers the \textbf{inauguration} of a truly representative government. Young Cameroonians must register on the \textbf{voter register} and participate actively. As the saying goes, “If you don't vote, you lose the right to complain.” Let us make our voices heard in the next election.”

(Traduction : “Voter n'est pas seulement un droit ; c'est un devoir civique fondamental pour tout citoyen. Dans une démocratie, l'électorat détient le pouvoir ultime : celui de choisir ses dirigeants par le bulletin de vote. Lorsque les électeurs se rendent au bureau de vote, étudient les programmes et expriment leur vote pour un candidat, ils exigent la reddition de comptes. Une forte participation confère de la légitimité au vainqueur et assure une transition pacifique du pouvoir. Inversement, l'abstention affaiblit le mandat du président élu et met en péril l'investiture d'un gouvernement véritablement représentatif. Les jeunes Camerounais doivent s'inscrire sur la liste électorale et participer activement. Comme le dit le proverbe : “Si tu ne votes pas, tu perds le droit de te plaindre.” Faisons entendre notre voix lors de la prochaine élection.”)

\textbf{Vérification} : 12 mots du champ lexical employés ; structures variées (passif \emph{is granted}, modaux \emph{must}, gérondifs) ; connecteurs (\emph{Conversely, Consequently}) ; orthographe \emph{government, ballot, manifesto} respectée.
\end{exoresolu}

\begin{attention}[Les 7 erreurs qui coûtent des points au Bac]
\begin{enumerate}
\item \textbf{“To make a vote” au lieu de \eng{to cast a vote}.} Calque du français “faire un vote”. De même : \eng{to make a speech}, \eng{to make a promise}, mais \eng{to hold an election} — chaque nom a son verbe attitré.
\item \textbf{Oublier la préposition après \emph{vote}.} On écrit \eng{to vote \textbf{for} a candidate} (et jamais “to vote him”). Mais attention : \eng{to elect somebody president} se dit \emph{sans} préposition.
\item \textbf{Confondre \eng{deputy} et “député”.} Un “député” français est \eng{a Member of Parliament (MP)} ; \eng{a deputy} est un adjoint. Autre piège : \eng{to assist} = aider ; “assister à un meeting” = \eng{to attend a rally}.
\item \textbf{Mal orthographier les mots clés} : \emph{government} (avec \emph{n}), \emph{candidate} (avec \emph{-date}), \emph{ballot} (b-a-l-l-o-t), \emph{referendum}, \emph{manifesto}. Une faute sur le mot central du sujet pénalise l'ensemble de la copie.
\item \textbf{Employer le mauvais temps dans le récit électoral.} Le déroulement d'un scrutin passé se raconte au prétérit : \eng{Voters queued and cast their ballots.} Réserve le futur aux promesses (\eng{The elected mayor will repair the roads}) et le présent aux vérités générales (\eng{Voting is a right and a duty}).
\item \textbf{Confondre \eng{election} (l'événement) et \eng{electorate} (l'ensemble des électeurs).} \eng{The election was peaceful.} / \eng{The electorate is angry.}
\item \textbf{Traduire “majorité” par \eng{majority} sans préciser.} \eng{Absolute majority} (majorité absolue), \eng{simple majority} (majorité relative), \eng{qualified majority} (majorité qualifiée). \eng{Majority} seul peut prêter à confusion.
\end{enumerate}
\end{attention}

\begin{savaistu}[Contexte camerounais \& Culture anglophone]
\begin{itemize}
\item \textbf{ELECAM} (\emph{Elections Cameroon}) : organe chargé de l'organisation des élections au Cameroun. Vocabulaire associé : \eng{biometric voter register} (fichier biométrique), \eng{indelible ink} (encre indélébile), \eng{polling station} (bureau de vote).
\item \textbf{Scrutins} : Présidentielle (\emph{Presidential election}), Législatives (\emph{Legislative / Parliamentary elections}), Municipales (\emph{Municipal / Council elections}), Régionales (\emph{Regional elections}).
\item \textbf{Système} : Majoritaire à un tour pour la présidentielle ; mixte (majoritaire/proportionnel) pour les législatives et municipales.
\item \textbf{Anglophonie} : Au Royaume-Uni, le Premier ministre n'est pas élu directement ; c'est le chef du parti majoritaire à la Chambre des Communes (\emph{House of Commons}). Aux USA, le Président est élu au suffrage universel indirect via le \emph{Electoral College} (Grand Électeurs). Le vocabulaire \eng{swing state}, \eng{primaries}, \emph{caucus} est spécifique au système US.
\end{itemize}
\end{savaistu}

\newpage

\coursec{Exercices}

\courssub{Je vérifie mes connaissances}

\begin{exercice}[QCM : le bon mot]
Pour chaque phrase, choisissez la bonne réponse (A, B ou C).
\begin{enumerate}
\item Before voting, every citizen must \cle{register} on the electoral \ldots\ 
\begin{itemize}
\item[A.] list \item[B.] paper \item[C.] page
\end{itemize}
\item The \ldots\ is the candidate who already holds the office.
\begin{itemize}
\item[A.] challenger \item[B.] incumbent \item[C.] opponent
\end{itemize}
\item Voters mark their choice on a \ldots\ paper.
\begin{itemize}
\item[A.] bulletin \item[B.] ballot \item[C.] ticket
\end{itemize}
\item She won the election by a \ldots\ : more than 70\,\% of the votes.
\begin{itemize}
\item[A.] landslide \item[B.] earthquake \item[C.] flood
\end{itemize}
\item After losing, the candidate \ldots\ defeat in a televised speech.
\begin{itemize}
\item[A.] admitted \item[B.] conceded \item[C.] confessed
\end{itemize}
\end{enumerate}
\end{exercice}

\begin{exercice}[Vrai ou faux — justifiez]
Lisez l'extrait ci-dessous, puis dites si les affirmations sont vraies (V) ou fausses (F) en citant le mot-clé du texte qui justifie votre réponse.
\begin{textebox}[Excerpt from a concession speech]
“My fellow citizens, the people have spoken. I have just phoned my opponent to congratulate her on her victory, and tonight I concede this race. Our campaign was tough but fair, and I thank the thousands of volunteers who knocked on doors and handed out leaflets. Though we lost, our ideas will keep fighting.”
\end{textebox}
\begin{enumerate}
\item The speaker won the election.
\item The speaker called his opponent.
\item The campaign was violent.
\item Volunteers distributed leaflets.
\item The speaker abandons his ideas completely.
\end{enumerate}
\end{exercice}

\begin{exercice}[Familles de mots]
Complétez le tableau avec la forme manquante de chaque famille.
\begin{center}
\begin{tabular}{|l|l|l|}
\hline
\rowcolor{popblueL} Verbe & Nom & Adjectif \\
\hline
to elect & \ldots\ldots\ldots & electoral \\
\hline
to vote & \ldots\ldots\ldots & \ldots\ldots\ldots \\
\hline
\ldots\ldots\ldots & campaign & \ldots\ldots\ldots \\
\hline
to nominate & \ldots\ldots\ldots & \ldots\ldots\ldots \\
\hline
\ldots\ldots\ldots & participation & \ldots\ldots\ldots \\
\hline
\end{tabular}
\end{center}
\end{exercice}

\begin{exercice}[Prépositions à compléter]
Complétez chaque phrase avec la préposition correcte (\emph{for, in, on, to, by, with}).
\begin{enumerate}
\item She campaigned \ldots\ the theme of free education.
\item He was elected \ldots\ a large majority.
\item The turnout \ldots\ the last election was very low.
\item Citizens voted \ldots\ the candidate of their choice.
\item The results were certified \ldots\ international observers.
\end{enumerate}
\end{exercice}

\courssub{Je m'entraîne}

\begin{exercice}[Traduction thématique]
Traduisez en anglais les phrases suivantes. Attention aux faux amis et aux collocations.
\begin{enumerate}
\item Le scrutin s'est déroulé sans incident majeur.
\item Il a déposé un recours auprès du Conseil Constitutionnel.
\item Le taux de participation a été faible.
\item Elle a fait campagne sur le thème de l'éducation.
\item Les observateurs ont certifié les résultats.
\end{enumerate}
\end{exercice}

\begin{exercice}[Appariement : verbes et collocations]
Associez chaque verbe de la colonne A au nom de la colonne B avec lequel il forme une collocation forte.
\begin{center}
\begin{tabular}{|l|l|}
\hline
\rowcolor{popblueL} A — Verbes & B — Noms \\
\hline
1. to cast & a. a rally \\
2. to hold & b. defeat \\
3. to file & c. a landslide victory \\
4. to concede & d. a vote \\
5. to win & e. a petition \\
6. to rig & f. a speech \\
7. to deliver & g. an election \\
8. to count & h. the ballots \\
\hline
\end{tabular}
\end{center}
\end{exercice}

\begin{exercice}[Texte à trous : le jour du vote]
Complétez le texte avec les mots suivants : \emph{polling station — ballot box — indelible ink — identity card — queue — sealed — observers — turnout — presiding officer — booth}.
\begin{textebox}[Election day in Yaoundé]
On election day, voters arrived early at the \ldots\ldots\ldots and stood in a long \ldots\ldots\ldots. Each voter showed an \ldots\ldots\ldots before entering the voting \ldots\ldots\ldots to mark the ballot paper in secret. The paper was then folded and dropped into the \ldots\ldots\ldots. To prevent double voting, the voter's finger was marked with \ldots\ldots\ldots. The \ldots\ldots\ldots supervised the whole operation, while national and international \ldots\ldots\ldots watched carefully. At the end of the day, the ballot box was \ldots\ldots\ldots before counting. Officials noted that the \ldots\ldots\ldots was higher than expected.
\end{textebox}
\end{exercice}

\begin{exercice}[Transformation de phrases]
Réécrivez chaque phrase en utilisant le mot entre parenthèses, sans changer le sens.
\begin{enumerate}
\item The candidate admitted that he had lost the election. (concede)
\item More than six million people voted in the election. (turnout)
\item The party organised a big public meeting in Douala. (rally)
\item She presented her political programme to the press. (manifesto)
\item They elected him president for a second time. (re-elect)
\end{enumerate}
\end{exercice}

\courssub{Je me prépare au Bac}

\begin{exercice}[Compréhension et langue — type Bac]
Lisez le texte suivant, puis répondez aux questions.
\begin{textebox}[A historic vote in Buea]
When the polling stations opened at 8 a.m. in Buea, hundreds of voters were already waiting patiently in line. For many young people, it was the first time they would cast a ballot. “I registered last year, and I have been waiting for this day ever since,” said Mirabel, a nineteen-year-old student. “Voting is not just a right; it is a duty.”

The campaign had lasted three weeks. Candidates held rallies in markets and schools, distributed leaflets and debated on local radio stations. The incumbent promised to improve roads and hospitals, while her main challenger campaigned on job creation for the youth.

Throughout the day, national and international observers visited polling stations to make sure the vote was free and fair. By 6 p.m., the ballot boxes were sealed and taken to the counting centre. The electoral commission announced that turnout had reached 65\,\%, a record for the region. The defeated candidate phoned his opponent and conceded defeat, calling for calm and unity.
\end{textebox}
\textbf{A. Comprehension questions} (répondez en anglais, avec vos propres mots)
\begin{enumerate}
\item Why was this election special for Mirabel? (2 pts)
\item Mention two campaign activities described in the text. (2 pts)
\item What did the incumbent promise? (2 pts)
\item What role did the observers play? (2 pts)
\item How did the defeated candidate react? (2 pts)
\end{enumerate}
\textbf{B. Language work}
\begin{enumerate}
\item Find in the text the noun formed from the verb \emph{to elect}. (1 pt)
\item Find in the text a synonym of “to admit defeat”. (1 pt)
\item Give the opposite of “the incumbent”. (1 pt)
\item Put into the passive voice: “The electoral commission announced the results.” (2 pts)
\end{enumerate}
\end{exercice}

\begin{exercice}[Traduction — type Bac]
Traduisez en anglais le passage suivant. (10 pts)
\begin{textebox}[Texte à traduire]
Les élections sont un moment important dans la vie d'un pays. Avant le jour du scrutin, les candidats font campagne dans tout le pays et présentent leur programme aux électeurs. Le jour du vote, chaque citoyen doit se munir de sa carte d'électeur. Après le dépouillement, la commission électorale proclame les résultats. Le candidat perdant peut contester les résultats devant les tribunaux, mais il doit respecter la décision finale.
\end{textebox}
\end{exercice}

\begin{exercice}[Composition — type Bac]
Choisissez \textbf{un} des deux sujets suivants et rédigez une composition en anglais de \textbf{250 à 300 mots}. (20 pts)
\begin{enumerate}
\item \textbf{Sujet 1 :} “Young people should take an active part in elections.” Do you agree? Write an essay giving your opinion, with at least three arguments illustrated by examples from Cameroon or other countries. Use at least five words from the vocabulary of campaigning and voting (e.g. \emph{turnout, ballot, campaign, candidate, polling station, manifesto}).
\item \textbf{Sujet 2 :} You were an international observer during the last municipal elections in your town. Write a formal report to your organisation describing how the vote was conducted (registration, polling stations, counting, results) and stating whether the election was free and fair. Use at least two collocations, one word from the family of \emph{elect}, and appropriate linking words.
\end{enumerate}
\end{exercice}

\newpage

\coursec{Corrigés des exercices}

\coursec{Corrigés des exercices de la leçon 47}

\begin{corrige}[Exercice 1 — QCM : le bon mot]
\begin{enumerate}
\item \textbf{A. list} — \eng{Before voting, every citizen must register on the electoral list.} « Avant de voter, chaque citoyen doit s'inscrire sur la liste électorale. » La collocation correcte est \cle{electoral list} ; \eng{paper} et \eng{page} ne s'emploient pas dans ce contexte.
\item \textbf{B. incumbent} — \eng{The incumbent is the candidate who already holds the office.} « Le sortant est le candidat qui occupe déjà le poste. » Le \eng{challenger} est le candidat qui défie le sortant ; \eng{opponent} désigne tout adversaire.
\item \textbf{B. ballot} — \eng{Voters mark their choice on a ballot paper.} « Les électeurs indiquent leur choix sur un bulletin de vote. » Attention au faux ami : \eng{bulletin} signifie « bulletin d'information », jamais « bulletin de vote ».
\item \textbf{A. landslide} — \eng{She won the election by a landslide.} « Elle a remporté l'élection par une victoire écrasante (raz-de-marée). » \eng{Earthquake} et \eng{flood} ne forment pas de collocation avec \eng{win an election}.
\item \textbf{B. conceded} — \eng{After losing, the candidate conceded defeat in a televised speech.} « Après sa défaite, le candidat a reconnu sa défaite dans un discours télévisé. » La collocation idiomatique est \cle{to concede defeat}. \eng{Admit defeat} est possible, mais \eng{confess} (« avouer une faute ») est un faux ami du français « confesser ».
\end{enumerate}
\end{corrige}

\begin{corrige}[Exercice 2 — Vrai ou faux — justifiez]
\begin{enumerate}
\item \textbf{F} — The speaker lost: \eng{“tonight I concede this race”} et \eng{“Though we lost”}. Il reconnaît sa défaite.
\item \textbf{V} — \eng{“I have just phoned my opponent to congratulate her on her victory.”} Il a téléphoné à son adversaire.
\item \textbf{F} — \eng{“Our campaign was tough but fair.”} La campagne fut rude mais loyale, pas violente.
\item \textbf{V} — \eng{“volunteers who knocked on doors and handed out leaflets.”} \eng{Hand out} = distribuer.
\item \textbf{F} — \eng{“our ideas will keep fighting.”} Ses idées continueront de se battre ; il n'y renonce pas.
\end{enumerate}
\textbf{Point de méthode :} au Bac, la justification doit toujours être une citation exacte du texte, entre guillemets anglais (“ … ”), accompagnée d'une courte explication dans vos propres mots.
\end{corrige}

\begin{corrige}[Exercice 3 — Familles de mots]
\begin{center}
\begin{tabular}{|l|l|l|}
\hline
\rowcolor{popblueL} Verbe & Nom & Adjectif \\
\hline
to elect & \textbf{election} (ou \textbf{elector}) & electoral \\
\hline
to vote & \textbf{vote} (ou \textbf{voter}) & \textbf{voting} \\
\hline
\textbf{to campaign} & campaign & \textbf{campaigning} \\
\hline
to nominate & \textbf{nomination} (ou \textbf{nominee}) & \textbf{nominated} / \textbf{nominative} \\
\hline
\textbf{to participate} & participation & \textbf{participating} / \textbf{participatory} \\
\hline
\end{tabular}
\end{center}
\textbf{Remarques :}
\begin{itemize}
\item \eng{Election} se prononce « i-lèk-cheune » : attention, le \emph{-tion} se prononce « cheune » (/ʃən/), pas « sion ».
\item \eng{Nominee} (« le candidat désigné ») se prononce « no-mi-ni » (/ˌnɒmɪˈniː/) ; ne pas confondre avec \eng{nomination} (« la désignation »).
\item \eng{Voter} désigne à la fois l'électeur et, selon le contexte, le fait de voter n'est exprimé que par le nom \eng{vote} : \eng{a low vote} est rare ; on préfère \eng{a low turnout}.
\end{itemize}
\end{corrige}

\begin{corrige}[Exercice 4 — Prépositions à compléter]
\begin{enumerate}
\item \textbf{on} — \eng{She campaigned on the theme of free education.} « Elle a fait campagne sur le thème de l'éducation gratuite. » Collocation : \cle{to campaign on an issue}.
\item \textbf{by} — \eng{He was elected by a large majority.} « Il a été élu par une large majorité. » \cle{by a majority of...} ; on dit aussi \eng{with a large majority} pour insister sur le score obtenu.
\item \textbf{in} — \eng{The turnout in the last election was very low.} « Le taux de participation aux dernières élections était très faible. » \cle{turnout in an election} ; \eng{at an election} est aussi accepté en anglais britannique.
\item \textbf{for} — \eng{Citizens voted for the candidate of their choice.} « Les citoyens ont voté pour le candidat de leur choix. » \cle{to vote for} (contre : \eng{to vote against}).
\item \textbf{by} — \eng{The results were certified by international observers.} « Les résultats ont été certifiés par des observateurs internationaux. » Complément d'agent du passif introduit par \cle{by}.
\end{enumerate}
\end{corrige}

\begin{corrige}[Exercice 5 — Traduction thématique]
\begin{enumerate}
\item \eng{The poll went off without any major incident.} ou \eng{The vote took place peacefully, without any major incident.} — Faux ami : \eng{scrutin} ne se traduit pas par \eng{scrutiny} (« examen minutieux ») ; on dit \eng{poll}, \eng{vote} ou \eng{ballot}.
\item \eng{He filed a petition with the Constitutional Council.} ou \eng{He lodged an appeal with the Constitutional Council.} — Collocation : \cle{to file a petition / an appeal} ; \cle{with} introduit l'institution saisie.
\item \eng{The turnout was low.} ou \eng{Voter turnout was low.} — Ne dites pas \eng{the participation rate} (calque) ; le mot consacré est \cle{turnout}.
\item \eng{She campaigned on the theme of education.} — \cle{to campaign on} + thème ; \eng{to make campaign} est un calque du français « faire campagne ».
\item \eng{The observers certified the results.} — Attention à l'orthographe : \eng{certified} (\emph{-y} devient \emph{-i-} devant \emph{-ed}).
\end{enumerate}
\end{corrige}

\begin{corrige}[Exercice 6 — Appariement : verbes et collocations]
\begin{center}
\begin{tabular}{|l|l|l|}
\hline
\rowcolor{popblueL} Collocation & Traduction & Réponse \\
\hline
to cast a vote & exprimer son vote & 1 — d \\
\hline
to hold a rally & organiser un meeting & 2 — a \\
\hline
to file a petition & déposer un recours & 3 — e \\
\hline
to concede defeat & reconnaître sa défaite & 4 — b \\
\hline
to win a landslide victory & remporter une victoire écrasante & 5 — c \\
\hline
to rig an election & truquer une élection & 6 — g \\
\hline
to deliver a speech & prononcer un discours & 7 — f \\
\hline
to count the ballots & dépouiller les bulletins & 8 — h \\
\hline
\end{tabular}
\end{center}
\textbf{Points de vigilance :} on ne dit pas \eng{to do a speech} mais \cle{to deliver / to give a speech} ; \cle{to cast a vote / a ballot} est le verbe noble du vote ; \eng{to rig} est péjoratif et signifie « truquer, frauder ».
\end{corrige}

\begin{corrige}[Exercice 7 — Texte à trous : le jour du vote]
\begin{textebox}[Election day in Yaoundé — corrected version]
On election day, voters arrived early at the \textbf{polling station} and stood in a long \textbf{queue}. Each voter showed an \textbf{identity card} before entering the voting \textbf{booth} to mark the ballot paper in secret. The paper was then folded and dropped into the \textbf{ballot box}. To prevent double voting, the voter's finger was marked with \textbf{indelible ink}. The \textbf{presiding officer} supervised the whole operation, while national and international \textbf{observers} watched carefully. At the end of the day, the ballot box was \textbf{sealed} before counting. Officials noted that the \textbf{turnout} was higher than expected.
\end{textebox}
\textbf{Justifications :} \cle{polling station} = bureau de vote ; \cle{queue} = file d'attente (anglais britannique ; américain : \eng{line}) ; \cle{voting booth} = isoloir ; \cle{ballot box} = urne ; \cle{indelible ink} = encre indélébile ; \cle{presiding officer} = président du bureau de vote ; \cle{sealed} = scellée (participe passé, accord avec la voix passive \eng{was sealed}) ; \cle{turnout} = taux de participation.
\end{corrige}

\begin{corrige}[Exercice 8 — Transformation de phrases]
\begin{enumerate}
\item \eng{The candidate conceded (that he had lost) the election.} ou mieux : \eng{The candidate conceded defeat.} — \cle{to concede} + nom ou \emph{that}-proposition.
\item \eng{The turnout in the election was more than six million (voters).} ou \eng{Turnout reached over six million voters.} — \cle{turnout} remplace \eng{the number of people who voted}.
\item \eng{The party held a big rally in Douala.} — Collocation : \cle{to hold a rally}.
\item \eng{She presented her manifesto to the press.} — \cle{manifesto} = programme électoral (faux ami : pas « manifeste » au sens philosophique uniquement).
\item \eng{They re-elected him president.} — Le préfixe \emph{re-} exprime la répétition ; noter la construction \cle{to elect / re-elect someone president} (attribut du complément sans \emph{as}).
\end{enumerate}
\end{corrige}

\begin{corrige}[Exercice 9 — Compréhension et langue — type Bac]
\textbf{A. Comprehension questions} (réponses modèles)
\begin{enumerate}
\item \eng{It was special for her because it was the first time she would cast a ballot; she had registered the year before and had been waiting for this day ever since.} (2 pts : 1 pt pour l'idée « first vote », 1 pt pour la formulation correcte.)
\item \eng{Candidates held rallies in markets and schools, distributed leaflets and debated on local radio stations.} (2 pts : 1 pt par activité citée.)
\item \eng{The incumbent promised to improve roads and hospitals.} (2 pts.)
\item \eng{They visited polling stations to make sure the vote was free and fair.} (2 pts.)
\item \eng{He phoned his opponent, conceded defeat and called for calm and unity.} (2 pts.)
\end{enumerate}
\textbf{B. Language work}
\begin{enumerate}
\item \eng{election} (dans le titre implicite et le contexte ; on accepte aussi \eng{electoral} comme adjectif de la même famille — 1 pt).
\item \eng{to concede defeat} (1 pt).
\item \eng{the challenger} (ou \eng{the opponent} — 1 pt).
\item \eng{The results were announced by the electoral commission.} (2 pts : 1 pt pour l'auxiliaire \eng{were} au prétérit, 1 pt pour le participe passé \eng{announced} ; le complément d'agent peut être omis.)
\end{enumerate}
\textbf{Points de vigilance :} ne recopiez pas de longues phrases du texte pour les questions de compréhension : reformulez. Au passif, n'oubliez pas l'accord de l'auxiliaire \emph{be} avec le nouveau sujet (\eng{the results were}, pas \eng{was}).
\end{corrige}

\begin{corrige}[Exercice 10 — Traduction — type Bac]
\textbf{Proposition de traduction :}
\begin{textebox}[Model translation]
Elections are an important moment in the life of a country. Before polling day, candidates campaign all over the country and present their manifesto to the voters. On election day, every citizen must bring his or her voter's card. After the counting of the votes, the electoral commission proclaims the results. The defeated candidate may challenge the results in court, but he or she must respect the final decision.
\end{textebox}
\textbf{Barème indicatif (10 pts) :} 2 pts par phrase (1 pt pour le lexique correct, 1 pt pour la grammaire et la syntaxe).
\textbf{Points de vigilance :}
\begin{itemize}
\item « le jour du scrutin » : \eng{polling day / election day}, jamais \eng{the day of the scrutiny}.
\item « faire campagne » : \cle{to campaign} (verbe), pas \eng{to make campaign}.
\item « carte d'électeur » : \eng{voter's card} ; « se munir de » : \eng{to bring / to carry}.
\item « dépouillement » : \eng{the counting of the votes / the count}.
\item « contester les résultats » : \cle{to challenge the results} ; « devant les tribunaux » : \eng{in court} (sans article).
\item « proclamer » : \eng{to proclaim / to announce officially}.
\end{itemize}
\end{corrige}

\begin{corrige}[Exercice 11 — Composition — type Bac]
\textbf{Barème indicatif (20 pts) :} contenu et pertinence des arguments (6 pts), organisation et connecteurs (4 pts), richesse et exactitude du lexique électoral (5 pts), correction grammaticale et orthographique (5 pts).

\textbf{Modèle — Sujet 1 (opinion essay) :}
\begin{textebox}[Model essay — Subject 1]
In many countries, young people under twenty-five represent the largest part of the population, yet they often have the lowest turnout on election day. In my opinion, young people should take an active part in elections, for three main reasons.

First of all, voting is the most direct way for the youth to influence decisions that concern their future. When a candidate campaigns on job creation or free education, young voters can support the manifesto that best answers their needs. If they stay at home, politicians will simply ignore them.

Secondly, participating in elections builds responsible citizens. Registering on the electoral list, going to the polling station and casting a ballot teach young people that democracy is not a gift but a duty. In Cameroon, for instance, students who debate politics at school often become the most committed voters. I remember accompanying my mother to the polling station in 2018; that experience taught me the value of my vote.

Finally, young people can do more than vote: they can volunteer in campaigns, hand out leaflets, or even stand as candidates in municipal elections. A young councillor in Bamenda or Douala understands the problems of the youth better than anyone.

To conclude, abstaining is never a solution. If young people want leaders who listen to them, they must register, vote and take part in public life. The future belongs to those who show up at the ballot box.
\end{textebox}
(Environ 260 mots — nombre de mots respecté.)

\textbf{Modèle — Sujet 2 (rapport formel) :}
\begin{textebox}[Model report — Subject 2]
To: The Director, National Observation Network\\
From: A. Ngo Bell, International Observer\\
Date: 12 October 2025\\
Subject: Observation report on the municipal elections in Bafoussam

I was deployed as an international observer in Bafoussam for the municipal elections. This report describes how the vote was conducted and assesses whether the election was free and fair.

First, concerning registration, the electoral list had been updated before the poll, and most voters I interviewed had received their voter's cards. However, a few young electors complained that their names were missing.

Secondly, on polling day, the polling stations I visited opened on time. Voters queued calmly, showed their identity cards and marked their ballot papers in the voting booths. Indelible ink was used systematically to prevent double voting, and the presiding officers worked professionally.

Moreover, the counting of the ballots was conducted transparently, in the presence of party representatives and observers. The ballot boxes had been properly sealed before being opened. The electoral commission later announced the results, and turnout reached about 60\,\%.

In conclusion, despite minor registration problems, I consider that the election was generally free and fair. I recommend that the electoral register be updated more regularly before the next election. I also recommend civic education campaigns targeting first-time voters and the extension of biometric verification to all polling stations.
\end{textebox}
(Environ 260 mots — nombre de mots respecté.)

\textbf{Points de vigilance :}
\begin{itemize}
\item Respectez le nombre de mots (250--300) : indiquez-le à la fin de la copie.
\item Sujet 1 : annoncez clairement votre opinion dans l'introduction et utilisez des connecteurs (\eng{first of all, secondly, finally, to conclude}).
\item Sujet 2 : adoptez le format du rapport (destinataire, date, objet), un ton formel et le passif quand il convient (\eng{the ballot boxes had been sealed}).
\item Vérifiez les collocations exigées : \eng{to cast a ballot}, \eng{to hold a rally}, \eng{to concede defeat}, et un mot de la famille de \emph{elect} (\eng{election, electoral, elector}).
\item Relisez : accords sujet-verbe, prétérit des verbes irréguliers, pas de calque du français.
\end{itemize}
\end{corrige}

\newpage

\coursec{Fiche bilan}

\begin{popfigure}
\begin{tikzpicture}[
  mindmap,
  grow cyclic,
  every node/.style={concept, align=center, font=\small},
  root concept/.append style={concept color=popblue, font=\bfseries\small, text=white},
  level 1 concept/.append style={level distance=3.6cm, sibling angle=60, font=\scriptsize},
  level 2 concept/.append style={level distance=2.1cm, font=\tiny}
]
\node[root concept]{Campaigning \& Voting for Leaders}
  child[concept color=poporange]{ node {The Electoral Process}
    child[concept color=poporangeL]{ node {election, poll, ballot} }
    child[concept color=poporangeL]{ node {candidate, voter, term} }
  }
  child[concept color=popgreen]{ node {Campaigning}
    child[concept color=popgreenL]{ node {run for office, canvass} }
    child[concept color=popgreenL]{ node {manifesto, pledge, rally} }
  }
  child[concept color=popgold]{ node {Media \& Speeches}
    child[concept color=popgoldL]{ node {deliver a speech} }
    child[concept color=popgoldL]{ node {coverage, debate, slogan} }
  }
  child[concept color=poppurple]{ node {Voting Day}
    child[concept color=poppurpleL]{ node {polling station, booth} }
    child[concept color=poppurpleL]{ node {cast a vote, ballot box} }
  }
  child[concept color=popgreen!70!black]{ node {After the Vote}
    child[concept color=popgreenL]{ node {count, results, turnout} }
    child[concept color=popgreenL]{ node {win, lose, swear in} }
  }
  child[concept color=popblue!70]{ node {Faux amis \& Pièges}
    child[concept color=popblueL]{ node {élection $\neq$ selection} }
    child[concept color=popblueL]{ node {vote \emph{for}, stand \emph{for}} }
  };
\end{tikzpicture}
\legende{Carte mentale : le lexique de la campagne et du vote}
\end{popfigure}

\begin{bilan}
\textbf{1. Lexique clé (anglais -- français)}
\begin{itemize}
  \item \eng{an election / general elections} : une élection / des élections générales ; \eng{a poll} : un scrutin, un sondage ; \eng{a by-election} : une élection partielle.
  \item \eng{a candidate} : un candidat ; \eng{to run for office / to stand for election} : se porter candidat ; \eng{a running mate} : un colistier.
  \item \eng{a voter} : un électeur ; \eng{the electorate} : le corps électoral ; \eng{a constituency} : une circonscription électorale.
  \item \eng{a campaign} : une campagne ; \eng{to campaign / to canvass} : faire campagne / faire du porte-à-porte ; \eng{a rally} : un meeting ; \eng{a manifesto} : un programme électoral ; \eng{a pledge / a promise} : un engagement / une promesse ; \eng{a slogan} : un slogan.
  \item \eng{a polling station} : un bureau de vote ; \eng{a polling booth} : un isoloir ; \eng{a ballot paper} : un bulletin de vote ; \eng{a ballot box} : une urne ; \eng{to cast one's vote} : déposer son vote ; \eng{to spoil a ballot} : rendre un bulletin nul.
  \item \eng{the turnout} : le taux de participation ; \eng{to count the votes} : dépouiller les votes ; \eng{a landslide victory} : une victoire écrasante ; \eng{a narrow defeat} : une défaite serrée ; \eng{to concede defeat} : reconnaître sa défaite ; \eng{to swear in} : investir, faire prêter serment ; \eng{a term of office} : un mandat.
\end{itemize}

\textbf{2. Collocations et expressions à connaître}
\begin{itemize}
  \item \eng{to hold / to call an election} : organiser / convoquer une élection ; \eng{to go to the polls} : aller aux urnes.
  \item \eng{to vote for / against somebody} : voter pour / contre quelqu'un ; \eng{to win / lose an election} : gagner / perdre une élection.
  \item \eng{free and fair elections} : des élections libres et transparentes ; \eng{a rigged election} : une élection truquée.
  \item \eng{to deliver / to make a speech} : prononcer un discours ; \eng{to keep / to break a promise} : tenir / ne pas tenir une promesse.
\end{itemize}

\textbf{3. Règles et points de langue à connaître par c\oe ur}
\begin{center}
\begin{tabularx}{\linewidth}{|X|X|}
\hline
\rowcolor{popblueL} \textbf{Règle} & \textbf{Exemple} \\
\hline
\eng{vote} + \textbf{for/against} (jamais \eng{vote someone} sans préposition) & \eng{She voted \textbf{for} the opposition candidate.} \\
\hline
\eng{stand for} = se présenter à ; \eng{run for} = être candidat à & \eng{He is \textbf{running for} president.} \\
\hline
\eng{elect somebody} (+ fonction sans article) & \eng{They elected her \textbf{president}.} \\
\hline
Passif fréquent dans les résultats & \eng{The votes \textbf{were counted} overnight.} \\
\hline
\eng{the electorate} est singulier collectif & \eng{The electorate \textbf{has} spoken.} \\
\hline
\end{tabularx}
\end{center}

\textbf{4. Faux amis et pièges}
\begin{itemize}
  \item \eng{an election} = une élection ; \emph{sélection} (au sens de choix) = \eng{a selection}.
  \item \eng{actually} = en fait (pas « actuellement » = \eng{currently}) : \eng{The current president is actually very popular.}
  \item \eng{a deputy} = un député (pas « adjoint » = \eng{an assistant / a deputy director} selon le cas).
  \item \eng{sensible} = raisonnable ; « sensible » = \eng{sensitive}.
\end{itemize}

\textbf{5. Méthode express}
\begin{enumerate}
  \item Repérer le moment du processus (avant / pendant / après le vote) pour choisir le bon champ lexical.
  \item Vérifier la préposition après le verbe (\eng{vote for}, \eng{stand for}, \eng{appeal to}).
  \item Employer les collocations apprises plutôt que des calques du français.
  \item À l'écrit, varier : \eng{candidate / contender / hopeful} ; \eng{win / secure / clinch a victory}.
\end{enumerate}
\end{bilan}

\begin{aretenir}[Checklist avant le Bac]
\begin{itemize}
  \item $\square$ Je connais le vocabulaire des trois phases : campagne, jour du vote, résultats.
  \item $\square$ Je sais employer \eng{to run for office} et \eng{to stand for election} sans les confondre.
  \item $\square$ Je mets toujours \textbf{for/against} après \eng{vote}.
  \item $\square$ Je distingue \eng{ballot paper} (bulletin), \eng{ballot box} (urne) et \eng{polling booth} (isoloir).
  \item $\square$ Je sais décrire un résultat : \eng{landslide victory}, \eng{narrow defeat}, \eng{turnout}.
  \item $\square$ Je connais les verbes d'action : \eng{cast, count, concede, swear in}.
  \item $\square$ J'évite le faux ami \eng{actually} pour « actuellement ».
  \item $\square$ Je sais construire le passif pour rapporter des résultats officiels.
  \item $\square$ Je peux prononcer correctement \eng{candidate} (« kan-di-date »), \eng{constituency} (« kon-sti-tiou-en-si ») et \eng{manoeuvre}.
  \item $\square$ Je sais rédiger un court paragraphe sur une élection au Cameroun en utilisant au moins six mots du lexique.
\end{itemize}
\end{aretenir}

\findecours

\end{document}
